% Lutte contre le VIH/Sida — SVT 1ère D — Cours signé OnBuch+
% Programme officiel MINESEC. Compiler avec : tectonic ND08-lutte-vih-sida-d.tex
\newcommand{\DOCMATIERE}{SVT}
\newcommand{\DOCNIVEAU}{1ère D}
\newcommand{\DOCMODULE}{Module 2 : L'Éducation à la Santé}
\newcommand{\DOCLECON}{ND08}
\newcommand{\DOCTITRE}{Lutte contre le VIH/Sida}
% OnBuch+ — préambule « cours signés » ENRICHI (compilé avec tectonic / XeLaTeX)
% Système visuel « Pop » d'OnBuch+ : fond crème (jamais blanc), bordures encre
% épaisses, ombres « dures » décalées, titres Archivo Black en capitales.
% Version MATHÉMATIQUES : graphes (pgfplots/tikz), boîtes pédagogiques
% (définition, propriété, méthode, attention, le savais-tu, exercice résolu,
% exercice, TP, bilan), page de garde et signature OnBuch+ ancrée en bas.
%
% Macros attendues dans main.tex AVANT \input{preamble} :
%   \DOCMATIERE  \DOCNIVEAU  \DOCMODULE  \DOCLECON (numéro)  \DOCTITRE
\documentclass[11pt]{article}

\usepackage{fontspec}
\usepackage[french]{babel}
\usepackage[a4paper,margin=2cm,top=3.4cm,bottom=2.8cm,headheight=1.4cm,footskip=1.3cm]{geometry}
\usepackage{amsmath,amssymb,amsfonts}
\usepackage{mathtools} % \xmapsto, \coloneqq... souvent utilisés par les modèles
\usepackage{cancel} % simplifications barrées (bilans de forces)
% \abs{z} (module, valeur absolue) et \norm{u} (norme), utilisés par les modèles
\providecommand{\abs}{}\let\abs\relax\DeclarePairedDelimiter{\abs}{\lvert}{\rvert}
\providecommand{\norm}{}\let\norm\relax\DeclarePairedDelimiter{\norm}{\lVert}{\rVert}
\usepackage[table]{xcolor} % [table] : fournit aussi \rowcolors (alternance de lignes)
\usepackage{pagecolor}
\usepackage{tikz}
\usetikzlibrary{arrows.meta,positioning,calc,shapes.geometric,shapes.misc,shapes.arrows,decorations.pathmorphing,decorations.pathreplacing,decorations.markings,patterns,shadows,mindmap,trees,fit,backgrounds,matrix,intersections,angles,quotes}
\usepackage{pgfplots}
\usepackage[version=4]{mhchem} % équations nucléaires \ce{^{235}_{92}U}
\usepackage[european,siunitx]{circuitikz} % schémas électriques
\pgfplotsset{compat=1.18}
\usepgfplotslibrary{fillbetween,groupplots} % aires entre courbes (calcul intégral)
\usepackage{enumitem}
\usepackage{fancyhdr}
% [most] charge toutes les bibliothèques tcolorbox (listings, minted, raster,
% documentation...) et gonfle inutilement la mémoire TeX sur un document long
% (~300 boîtes) : on ne charge que ce qui est réellement utilisé.
\usepackage[skins,breakable]{tcolorbox}
\usepackage{graphicx}
\usepackage{colortbl}
\usepackage{booktabs}
\usepackage{tabularx}
\usepackage{diagbox} % case d'en-tête diagonale des tableaux à double entrée
% Colonnes extensibles centrée / gauche / droite (C, L, R), souvent utilisées par les modèles
\newcolumntype{C}{>{\centering\arraybackslash}X}
\newcolumntype{L}{>{\raggedright\arraybackslash}X}
\newcolumntype{R}{>{\raggedleft\arraybackslash}X}
\usepackage{array}
\usepackage{multicol}
\usepackage{siunitx}
% parse-numbers=false : accepte aussi \SI{2,5 \times 10^{-3}}{...}
\sisetup{locale=FR,output-decimal-marker={,},group-minimum-digits=4,parse-numbers=false}
\DeclareSIUnit\ohm{\ensuremath{\symup{\Omega}}} % U+2126 absent de la police
\DeclareSIUnit\division{div} % réglages d'oscilloscope (ms/div, V/div)
\DeclareSIUnit\tr{tr} % tours (vitesse de rotation, ex. \SI{1500}{\tr\per\minute})
\DeclareSIUnit\ch{ch} % cheval-vapeur (puissance moteur, unité usuelle hors SI)
\DeclareSIUnit\franccfa{FCFA} % franc CFA (coûts, contexte camerounais)
\DeclareSIUnit\dioptre{\ensuremath{\delta}} % vergence d'une lentille (dioptrie)
\usepackage{qrcode}
% \degree : commande fréquemment utilisée par les modèles pour un angle ou une
% température en dehors de \SI{}{} (non fournie par les paquets chargés).
\providecommand{\degree}{^{\circ}}
% \qed (fin de démonstration), utilisé par les modèles sans amsthm
\providecommand{\qed}{\hfill$\square$}
% \barre : double barre des valeurs interdites dans un tableau de variations (tkz-tab)
\providecommand{\barre}{\ensuremath{\|}}
% environnement proof (amsthm non chargé) utilisé par les modèles
\ifdefined\proof\else\newenvironment{proof}[1][Démonstration]{\par\noindent\textit{#1.}\ }{\qed\par}\fi
\usepackage{lastpage}
\usepackage{hyperref}
\hypersetup{hidelinks}

% ============================================================
% Palette « Pop » OnBuch+ — mêmes hex que la classe P (plus_ui.dart)
% ============================================================
\definecolor{poporange}{HTML}{F59321}
\definecolor{popdark}{HTML}{E07A0C}
\definecolor{popcream}{HTML}{FFF6E8}
\definecolor{popink}{HTML}{1C1714}
\definecolor{poppurple}{HTML}{7A5AE0}
\definecolor{popgreen}{HTML}{2E9E6A}
\definecolor{poppink}{HTML}{E0457B}
\definecolor{popblue}{HTML}{2F7DE1}
\definecolor{popgold}{HTML}{C98A12}
\definecolor{popmuted}{HTML}{8A7F76}
\definecolor{popline}{HTML}{EADFD2}
\definecolor{popgray}{HTML}{8A7F76}
\colorlet{popgrey}{popgray}
% Alias tolérants (le modèle nomme parfois les couleurs différemment)
\colorlet{popwhite}{white}
\colorlet{popblack}{popink}
\colorlet{popred}{poppink}
\colorlet{popyellow}{popgold}
% Teintes claires pour fonds de tableaux / zones
\colorlet{popblueL}{popblue!10!white}
\colorlet{poporangeL}{poporange!14!white}
\colorlet{popgreenL}{popgreen!12!white}
\colorlet{poppurpleL}{poppurple!12!white}
\colorlet{poppinkL}{poppink!12!white}
\colorlet{popgoldL}{popgold!14!white}

\pagecolor{popcream}

% ============================================================
% Typographie
% ============================================================
\defaultfontfeatures{Ligatures=TeX}
\setmainfont{PlusJakartaSans-Medium.ttf}[Path=../../../fonts/,BoldFont=PlusJakartaSans-Bold.ttf]
% Maths en sans-serif (Fira Math) avec chiffres et lettres droites de Plus
% Jakarta, pour que les formules se fondent dans le texte.
\usepackage[math-style=ISO,bold-style=ISO]{unicode-math}
\setmathfont{FiraMath-Regular.otf}[Path=../../../fonts/]
\setmathfont{PlusJakartaSans-Medium.ttf}[Path=../../../fonts/,range={up/num,up/{Latin,latin},\mathup}]
\usepackage{icomma} % virgule décimale sans espace en mode math
% Caractères mathématiques Unicode absents des polices de texte (Plus Jakarta,
% Archivo Black) : rendus par leur équivalent mathématique, en texte comme en
% formule — sinon ils s'affichent en carré vide (ex. « Arithmétique dans ℤ »).
\usepackage{newunicodechar}
\newunicodechar{ℝ}{\ensuremath{\mathbb{R}}}
\newunicodechar{ℤ}{\ensuremath{\mathbb{Z}}}
\newunicodechar{ℂ}{\ensuremath{\mathbb{C}}}
\newunicodechar{ℕ}{\ensuremath{\mathbb{N}}}
\newunicodechar{ℚ}{\ensuremath{\mathbb{Q}}}
\newunicodechar{𝔻}{\ensuremath{\mathbb{D}}}
\newunicodechar{α}{\ensuremath{\alpha}}
\newunicodechar{β}{\ensuremath{\beta}}
\newunicodechar{γ}{\ensuremath{\gamma}}
\newunicodechar{δ}{\ensuremath{\delta}}
\newunicodechar{ε}{\ensuremath{\varepsilon}}
\newunicodechar{θ}{\ensuremath{\theta}}
\newunicodechar{λ}{\ensuremath{\lambda}}
\newunicodechar{μ}{\ensuremath{\mu}}
\newunicodechar{σ}{\ensuremath{\sigma}}
\newunicodechar{τ}{\ensuremath{\tau}}
\newunicodechar{φ}{\ensuremath{\varphi}}
\newunicodechar{ψ}{\ensuremath{\psi}}
\newunicodechar{ω}{\ensuremath{\omega}}
\newunicodechar{Σ}{\ensuremath{\Sigma}}
% majuscules produites par \MakeUppercase dans les titres (α-aminés -> Α)
\newunicodechar{Α}{\ensuremath{\alpha}}
\newunicodechar{Β}{\ensuremath{\beta}}
\newunicodechar{Γ}{\ensuremath{\Gamma}}
\newunicodechar{Δ}{\ensuremath{\Delta}}
\newunicodechar{Ε}{\ensuremath{\varepsilon}}
\newunicodechar{Θ}{\ensuremath{\Theta}}
\newunicodechar{Λ}{\ensuremath{\Lambda}}
\newunicodechar{Μ}{\ensuremath{\mu}}
\newunicodechar{Τ}{\ensuremath{\tau}}
\newunicodechar{Φ}{\ensuremath{\Phi}}
\newunicodechar{Ψ}{\ensuremath{\Psi}}
\newunicodechar{Ω}{\ensuremath{\Omega}}
\newunicodechar{↦}{\ensuremath{\mapsto}}
\newunicodechar{∈}{\ensuremath{\in}}
\newunicodechar{∉}{\ensuremath{\notin}}
\newunicodechar{∧}{\ensuremath{\wedge}}
\newunicodechar{⇔}{\ensuremath{\Leftrightarrow}}
\newunicodechar{⟺}{\ensuremath{\Longleftrightarrow}}
\newunicodechar{⇒}{\ensuremath{\Rightarrow}}
\newunicodechar{⟹}{\ensuremath{\Longrightarrow}}
\newunicodechar{∘}{\ensuremath{\circ}}
\newunicodechar{∪}{\ensuremath{\cup}}
\newunicodechar{∩}{\ensuremath{\cap}}
\newunicodechar{≡}{\ensuremath{\equiv}}
\newunicodechar{⊂}{\ensuremath{\subset}}
\newunicodechar{⊥}{\ensuremath{\perp}}
\newunicodechar{∃}{\ensuremath{\exists}}
\newunicodechar{∀}{\ensuremath{\forall}}
\newunicodechar{∛}{\ensuremath{\sqrt[3]{\,}}}
\newunicodechar{∜}{\ensuremath{\sqrt[4]{\,}}}
\newunicodechar{⃗}{\ensuremath{{}^{\to}}}
\newunicodechar{ⁿ}{\ifmmode^{n}\else\textsuperscript{\ensuremath{n}}\fi}
\newunicodechar{ˣ}{\ifmmode^{x}\else\textsuperscript{\ensuremath{x}}\fi}
\newunicodechar{ᵇ}{\ifmmode^{b}\else\textsuperscript{\ensuremath{b}}\fi}
\newunicodechar{ʳ}{\ifmmode^{r}\else\textsuperscript{\ensuremath{r}}\fi}
\newunicodechar{ᵘ}{\ifmmode^{u}\else\textsuperscript{\ensuremath{u}}\fi}
\newunicodechar{ʸ}{\ifmmode^{y}\else\textsuperscript{\ensuremath{y}}\fi}
\newunicodechar{ᵃ}{\ifmmode^{a}\else\textsuperscript{\ensuremath{a}}\fi}
\newunicodechar{ᵉ}{\ifmmode^{e}\else\textsuperscript{\ensuremath{e}}\fi}
\newunicodechar{ᵏ}{\ifmmode^{k}\else\textsuperscript{\ensuremath{k}}\fi}
\newunicodechar{ᵖ}{\ifmmode^{p}\else\textsuperscript{\ensuremath{p}}\fi}
\newunicodechar{ᴺ}{\ifmmode^{N}\else\textsuperscript{\ensuremath{N}}\fi}
\newunicodechar{ᶜ}{\ifmmode^{c}\else\textsuperscript{\ensuremath{c}}\fi}
\newunicodechar{ʰ}{\ifmmode^{h}\else\textsuperscript{\ensuremath{h}}\fi}
\newunicodechar{ᵠ}{\ifmmode^{q}\else\textsuperscript{\ensuremath{q}}\fi}
\newunicodechar{ₙ}{\ifmmode_{n}\else\textsubscript{\ensuremath{n}}\fi}
\newunicodechar{ₐ}{\ifmmode_{a}\else\textsubscript{\ensuremath{a}}\fi}
\newunicodechar{ᵢ}{\ifmmode_{i}\else\textsubscript{\ensuremath{i}}\fi}
\newunicodechar{ᵣ}{\ifmmode_{r}\else\textsubscript{\ensuremath{r}}\fi}
\newunicodechar{ₖ}{\ifmmode_{k}\else\textsubscript{\ensuremath{k}}\fi}
\newunicodechar{ᵦ}{\ifmmode_{\beta}\else\textsubscript{\ensuremath{\beta}}\fi}
\newunicodechar{ₓ}{\ifmmode_{x}\else\textsubscript{\ensuremath{x}}\fi}
\newfontfamily\popdisplay{ArchivoBlack-Regular.ttf}[Path=../../../fonts/]
\newfontfamily\poplogo{SpaceGrotesk-Bold.ttf}[Path=../../../fonts/]
\newfontfamily\popsemi{PlusJakartaSans-SemiBold.ttf}[Path=../../../fonts/]
\newfontfamily\popxbold{PlusJakartaSans-ExtraBold.ttf}[Path=../../../fonts/]
\newfontfamily\popxboldls{PlusJakartaSans-ExtraBold.ttf}[Path=../../../fonts/,LetterSpace=3.0]

\setlist[enumerate]{leftmargin=1.7em,itemsep=5pt,topsep=4pt}
\setlist[itemize]{leftmargin=1.7em,itemsep=4pt,topsep=4pt,label={\color{poporange}\textbullet}}
\setlength{\parindent}{0pt}
\setlength{\parskip}{5pt}
\renewcommand{\arraystretch}{1.25}

% pgfplots : style Pop par défaut
\pgfplotsset{
  every axis/.append style={axis line style={popink,line width=1pt},tick style={popink},
    grid style={popline,line width=0.5pt},label style={font=\small},tick label style={font=\footnotesize}},
  popcurve/.style={poporange,line width=1.8pt,no markers,smooth},
}
% Même style disponible dans un \draw TikZ simple (hors pgfplots)
\tikzset{popcurve/.style={poporange,line width=1.8pt}}
\tikzset{popgrid/.style={popline,very thin}}
\colorlet{popcurve}{poporange} % le nom du style est parfois utilisé comme couleur

% ============================================================
% Pill / squiggle
% ============================================================
\newcommand{\poppill}[3]{%
  \tikz[baseline=-0.55ex]\node[draw=popink,line width=1.2pt,rounded corners=2.6pt,%
    fill=#2,inner xsep=6.5pt,inner ysep=3pt]{%
    \popxboldls\fontsize{7}{7}\selectfont\color{#3}\MakeUppercase{#1}};%
}
\newcommand{\popsquiggle}[1]{%
  \tikz[baseline=-0.4ex]{
    \draw[#1,line width=2.6pt,line cap=round]
      plot [smooth] coordinates {(0,0.09) (0.34,0.01) (0.68,0.21) (1.02,0.01) (1.36,0.10)};
  }%
}
% Mot-clé surligné (marqueur orange sous le texte)
\newcommand{\cle}[1]{{\popxbold\color{popdark}#1}}

% ============================================================
% En-tête / pied
% ============================================================
\pagestyle{fancy}
\fancyhf{}
\renewcommand{\headrulewidth}{0pt}
\renewcommand{\footrulewidth}{0pt}
\lhead{\raisebox{-0.15ex}{\poplogo\fontsize{13}{13}\selectfont\color{popink}OnBuch\color{poporange}+}}
\rhead{\poppill{\DOCMATIERE\ \textbullet\ \DOCNIVEAU\ \textbullet\ Leçon \DOCLECON}{white}{popink}}
\lfoot{\popxbold\fontsize{7.6}{7.6}\selectfont\color{popmuted}\MakeUppercase{Cours signé OnBuch+}\ \textbullet\ {\popsemi\DOCTITRE}}
\rfoot{\tikz[baseline=-0.6ex]\node[rounded corners=8.5pt,fill=popink,minimum height=17pt,minimum width=17pt,inner xsep=5pt,inner sep=0]{\popxbold\fontsize{8}{8}\selectfont\color{popcream}\thepage\,/\,\pageref*{LastPage}};}

% ============================================================
% Titres
% ============================================================
\newcommand{\chaptitre}[2]{%
  \begin{tcolorbox}[enhanced,colback=poporange,colframe=popink,boxrule=2.6pt,arc=11pt,
      drop shadow=popink,left=15pt,right=15pt,top=12pt,bottom=11pt,before skip=3pt,after skip=18pt]
    \poppill{#2}{popink}{popcream}\par\vspace{9pt}
    {\raggedright\hyphenpenalty=10000\exhyphenpenalty=10000\popdisplay\fontsize{22}{24}\selectfont\color{popcream}\MakeUppercase{#1}\par}
  \end{tcolorbox}
}
% Section numérotée : pastille orange avec le numéro + titre Archivo
\newcounter{popsec}
\newcommand{\coursec}[1]{%
  \stepcounter{popsec}\setcounter{popsub}{0}%
  \par\vspace{16pt}\noindent
  \begin{minipage}{\linewidth}
  \tikz[baseline=-0.7ex]\node[draw=popink,line width=1.6pt,fill=poporange,rounded corners=4pt,minimum size=22pt,inner sep=0]{\popdisplay\fontsize{12}{12}\selectfont\color{popcream}\thepopsec};%
  \hspace{8pt}{\popdisplay\fontsize{15}{17}\selectfont\color{popink}\MakeUppercase{#1}}\par
  \vspace{4pt}\noindent{\color{poporange}\rule{36pt}{3.6pt}}
  \end{minipage}\par\nopagebreak\vspace{8pt}
}
\newcounter{popsub}
\newcommand{\courssub}[1]{%
  \stepcounter{popsub}%
  \par\vspace{9pt}\noindent{\popxbold\fontsize{11.5}{13}\selectfont\color{popink}{\color{poporange}\thepopsec.\thepopsub}\hspace{6pt}#1}\par\nopagebreak\vspace{4pt}
}

% ============================================================
% Boîtes pédagogiques — même recette : fond blanc, bordure épaisse colorée,
% ombre dure encre, badge plein. Titre optionnel [#1].
% ============================================================
\tcbset{popbase/.style={breakable,enhanced,colback=white,boxrule=2.3pt,arc=9pt,
  shadow={3pt}{-3pt}{0pt}{popink},left=14pt,right=14pt,top=11pt,bottom=11pt,before skip=11pt,after skip=15pt,
  fontupper=\popsemi\fontsize{10.6}{14.5}\selectfont\color{popink},
  fontlower=\popsemi\fontsize{10.6}{14.5}\selectfont\color{popink}}}
% Coche dessinée (le glyphe ✓ n'existe pas dans les polices du document)
\newcommand{\popcheck}[1]{\tikz[baseline=-0.5ex]\draw[#1,line width=1.6pt,line cap=round,line join=round](-0.13,0.0)--(-0.04,-0.09)--(0.13,0.1);}
\providecommand{\coche}{\popcheck{popgreen}}
\providecommand{\resultat}[1]{\par\smallskip\noindent\fcolorbox{popgreen}{popgreenL}{\parbox{\dimexpr\linewidth-2\fboxsep-2\fboxrule\relax}{\textbf{Résultat.} #1}}\par\smallskip}
\newcommand{\popboxhead}[3]{\poppill{#1}{#2}{white}\ifx\relax#3\relax\else\hspace{6pt}{\popxbold\fontsize{10.6}{12}\selectfont\color{popink}#3}\fi\par\vspace{7pt}}

\newtcolorbox{aretenir}[1][]{popbase,colframe=popgreen,before upper={\popboxhead{À retenir}{popgreen}{#1}}}
\newtcolorbox{exemplebox}[1][]{popbase,colframe=poporange,before upper={\popboxhead{Exemple}{poporange}{#1}}}
\newtcolorbox{definition}[1][]{popbase,colframe=popblue,before upper={\popboxhead{Définition}{popblue}{#1}}}
\newtcolorbox{propriete}[1][]{popbase,colframe=poppurple,before upper={\popboxhead{Propriété}{poppurple}{#1}}}
\newtcolorbox{methode}[1][]{popbase,colframe=popgold,before upper={\popboxhead{Méthode}{popgold}{#1}}}
\newtcolorbox{attention}[1][]{popbase,colframe=poppink,before upper={\popboxhead{Attention}{poppink}{#1}}}
\newtcolorbox{experience}[1][]{popbase,colframe=popblue,colback=popblueL,before upper={\popboxhead{Expérience}{popblue}{#1}}}
% « Le savais-tu ? » : carte violette pleine (contexte, culture, Cameroun)
\newtcolorbox{savaistu}[1][]{popbase,colback=poppurple,colframe=popink,
  fontupper=\popsemi\fontsize{10.4}{14.2}\selectfont\color{white},
  before upper={\poppill{Le savais-tu\,?}{white}{poppurple}\ifx\relax#1\relax\else\hspace{6pt}{\popxbold\color{white}#1}\fi\par\vspace{7pt}}}
% Exercice résolu : énoncé en haut, \tcblower puis la solution sur fond crème
\newcounter{popexo}\newcounter{popexr}
\newtcolorbox{exoresolu}[1][]{popbase,colframe=poporange,colbacklower=popcream,
  segmentation style={popink,line width=1.2pt,dashed},
  before upper={\stepcounter{popexr}\popboxhead{Exercice résolu \thepopexr}{poporange}{#1}},
  before lower={\poppill{Solution}{popink}{popcream}\par\vspace{6pt}}}
% Exercice (non résolu en place) — la correction est dans la partie Corrigés
\newtcolorbox{exercice}[1][]{popbase,colframe=popink,
  before upper={\stepcounter{popexo}\popboxhead{Exercice \thepopexo}{popink}{#1}}}
% Corrigé
\newtcolorbox{corrige}[1][]{popbase,colframe=popgreen,colback=popgreenL,
  before upper={\popboxhead{Corrigé}{popgreen}{#1}}}
% Objectifs / prérequis / bilan
\newtcolorbox{objectifs}{popbase,colframe=popink,colback=poporangeL,
  before upper={\popboxhead{Objectifs de la leçon}{popink}{}}}
\newtcolorbox{prerequis}{popbase,colframe=popink,colback=popblueL,
  before upper={\popboxhead{Prérequis}{popblue}{}}}
\newtcolorbox{bilan}{popbase,colframe=popink,colback=white,boxrule=2.8pt,
  before upper={\popboxhead{Fiche bilan}{poporange}{L'essentiel à connaître pour l'examen}}}
% Figure encadrée avec légende
\newtcolorbox{popfigure}[1][]{enhanced,breakable=false,colback=white,colframe=popline,boxrule=1.4pt,arc=8pt,
  left=10pt,right=10pt,top=10pt,bottom=8pt,before skip=10pt,after skip=12pt,halign=center,
  lower separated=false,halign lower=center,
  fontlower=\popsemi\fontsize{9}{11.5}\selectfont\color{popmuted},
  #1}
\newcommand{\legende}[1]{\par\vspace{6pt}{\popsemi\fontsize{9}{11.5}\selectfont\color{popmuted}#1}}

% ============================================================
% Signature OnBuch+
% ============================================================
\newcommand{\poplogoinline}[1]{{\poplogo\fontsize{#1}{#1}\selectfont\color{popink}OnBuch\color{poporange}+}}
\newcommand{\signatureonbuch}{%
  \begin{tcolorbox}[enhanced,colback=popink,colframe=popink,boxrule=2.2pt,arc=10pt,
      drop shadow=poporange,left=14pt,right=14pt,top=10pt,bottom=10pt,before skip=0pt,after skip=0pt]
    \tikz[baseline=-0.7ex]\node[circle,fill=popgreen,draw=popcream,line width=1.3pt,minimum size=22pt,inner sep=0]{\popcheck{white}};%
    \hspace{9pt}\begin{minipage}[c]{0.62\linewidth}
      {\popxbold\fontsize{12}{13}\selectfont\color{popcream}Signé\ \textbullet\ {\poplogo OnBuch\color{poporange}+}}\par\vspace{2pt}
      {\raggedright\popsemi\fontsize{8}{10.5}\selectfont\color{popline}\MakeUppercase{Équipe pédagogique OnBuch+}\\\MakeUppercase{Conforme au programme officiel MINESEC}\par}
    \end{minipage}\hfill
    {\poplogo\fontsize{22}{22}\selectfont\color{popcream}OnBuch\color{poporange}+}
  \end{tcolorbox}%
}

% ============================================================
% Page de garde
% #1 titre · #2 matière · #3 niveau · #4 module · #5 n° leçon · #6 durée ·
% #7 description / objectifs (texte ou itemize)
% ============================================================
\newcommand{\pagegarde}[7]{%
  \thispagestyle{empty}
  \newgeometry{a4paper,margin=1.7cm,top=1.6cm,bottom=1.4cm}%
  \begin{tikzpicture}[remember picture,overlay]
    \fill[poporange!18] ([xshift=-3.2cm,yshift=-3.6cm]current page.north east) circle (4.6cm);
    \fill[poppurple!14] ([xshift=2.4cm,yshift=7.6cm]current page.south west) circle (3.8cm);
  \end{tikzpicture}%
  {\poplogo\fontsize{30}{30}\selectfont\color{popink}OnBuch\color{poporange}+}\hfill
  \raisebox{0.6ex}{\poppill{Cours signé \textbullet\ Premium}{popink}{popcream}}\par
  \vspace{1pt}\popsquiggle{poppurple}\par
  \vspace{8pt}
  \begin{tcolorbox}[enhanced,colback=poporange,colframe=popink,boxrule=2.8pt,arc=13pt,
      drop shadow=popink,left=18pt,right=18pt,top=12pt,bottom=13pt,after skip=10pt]
    \poppill{#2 \textbullet\ #3}{popink}{popcream}\hspace{5pt}\poppill{#4}{white}{popink}\par\vspace{8pt}
    {\popdisplay\fontsize{14}{14}\selectfont\color{popink}LEÇON #5}\par\vspace{4pt}
    {\raggedright\hyphenpenalty=10000\exhyphenpenalty=10000\popdisplay\fontsize{25}{27}\selectfont\color{popcream}\MakeUppercase{#1}\par}
    \vspace{8pt}
    {\popxbold\fontsize{9.5}{9.5}\selectfont\color{popink}\MakeUppercase{Durée conseillée : #6}}
  \end{tcolorbox}
  \begin{tcolorbox}[enhanced,colback=white,colframe=popink,boxrule=2.2pt,arc=10pt,
      drop shadow=popink,left=16pt,right=16pt,top=10pt,bottom=10pt,after skip=8pt]
    \poppill{Dans cette leçon}{poppurple}{white}\par\vspace{5pt}
    \popsemi\fontsize{9.3}{12.6}\selectfont\color{popink}#7
  \end{tcolorbox}
  \noindent\begin{minipage}[t]{0.487\linewidth}
    \begin{tcolorbox}[enhanced,colback=popgreen,colframe=popink,boxrule=2.2pt,arc=10pt,
        drop shadow=popink,left=11pt,right=11pt,top=10pt,bottom=10pt,height=3.1cm,valign=center]
      \tcbox[colback=white,colframe=popink,boxrule=1.3pt,arc=5pt,boxsep=0pt,nobeforeafter,
        left=3pt,right=3pt,top=3pt,bottom=3pt,valign=center]{\qrcode[height=1.9cm]{https://play.google.com/store/apps/details?id=com.luvvix.onbuch&hl=fr}}%
      \hspace{8pt}\begin{minipage}[c]{0.5\linewidth}
        {\popdisplay\fontsize{11}{12.5}\selectfont\color{white}\MakeUppercase{Télécharge OnBuch}}\par\vspace{4pt}
        {\raggedright\popsemi\fontsize{8.4}{11}\selectfont\color{white}Gratuit \textbullet\ Google Play\\Tuteur IA, annales, concours, résultats\par}
      \end{minipage}
    \end{tcolorbox}
  \end{minipage}\hfill
  \begin{minipage}[t]{0.487\linewidth}
    \begin{tcolorbox}[enhanced,colback=poppurple,colframe=popink,boxrule=2.2pt,arc=10pt,
        drop shadow=popink,left=11pt,right=11pt,top=10pt,bottom=10pt,height=3.1cm,valign=center]
      \tikz[baseline=-0.7ex]\node[circle,fill=white,draw=popink,line width=1.3pt,minimum size=1.6cm,inner sep=0]{\popdisplay\fontsize{24}{24}\selectfont\color{poppurple}+};%
      \hspace{8pt}\begin{minipage}[c]{0.6\linewidth}
        {\popdisplay\fontsize{11}{12.5}\selectfont\color{white}\MakeUppercase{Passe à OnBuch+}}\par\vspace{4pt}
        {\raggedright\popsemi\fontsize{8.4}{11}\selectfont\color{white}Tous les cours signés, Tuteur IA illimité, fiches PDF — onglet OnBuch+ de l'app\par}
      \end{minipage}
    \end{tcolorbox}
  \end{minipage}
  \vspace{8pt}
  \popcontenu
  \vfill
  \signatureonbuch
  \restoregeometry
  \newpage
}

% Rangée « Ce cours contient » (page de garde)
\newcommand{\poptile}[3]{% #1 couleur #2 gros texte #3 légende
  \begin{tcolorbox}[enhanced,colback=#1,colframe=popink,boxrule=2pt,arc=9pt,shadow={2.5pt}{-2.5pt}{0pt}{popink},
      left=6pt,right=6pt,top=9pt,bottom=9pt,halign=center,valign=center,height=1.75cm,width=\linewidth,nobeforeafter]
    {\popdisplay\fontsize{12.5}{13}\selectfont\color{white}\MakeUppercase{#2}}\par\vspace{3pt}
    {\centering\popsemi\fontsize{7.8}{9.5}\selectfont\color{white}#3\par}
  \end{tcolorbox}}
\newcommand{\popcontenu}{%
  \par\noindent{\popxboldls\fontsize{7.5}{7.5}\selectfont\color{popmuted}CE COURS SIGNÉ CONTIENT}\par\vspace{7pt}
  \noindent\begin{minipage}[t]{0.235\linewidth}\poptile{popblue}{Cours}{complet, illustré, pas à pas}\end{minipage}\hfill
  \begin{minipage}[t]{0.235\linewidth}\poptile{poporange}{Méthodes}{et exercices résolus}\end{minipage}\hfill
  \begin{minipage}[t]{0.235\linewidth}\poptile{poppink}{Exercices}{type examen + corrigés}\end{minipage}\hfill
  \begin{minipage}[t]{0.235\linewidth}\poptile{popgreen}{Fiche bilan}{l'essentiel à retenir}\end{minipage}\par}

% Sommaire
\newtcolorbox{sommaire}{enhanced,colback=white,colframe=popink,boxrule=2.2pt,arc=10pt,
  drop shadow=popink,left=18pt,right=18pt,top=13pt,bottom=13pt,before skip=6pt,after skip=20pt,
  before upper={\poppill{Sommaire}{poppurple}{white}\par\vspace{9pt}\popsemi\fontsize{10.6}{14.5}\selectfont\color{popink}}}

% Signature de fin de cours — poussée tout en bas de la dernière page
\newcommand{\findecours}{%
  \par\vfill\nopagebreak
  \begin{center}\popsquiggle{poporange}\end{center}\vspace{4pt}
  \signatureonbuch
}
\AtBeginDocument{\def\llbracket{\mathopen{[\mkern-3.5mu[}}\def\rrbracket{\mathclose{]\mkern-3.5mu]}}} % intervalles d'entiers (glyphe ⟦ absent de la police)
% Glyphes absents de Fira Math : redéfinis à partir de glyphes disponibles
\AtBeginDocument{%
  \def\setminus{\mathbin{\backslash}}\let\smallsetminus\setminus
  \def\vdots{\mathord{\vbox{\baselineskip3.2pt\lineskiplimit0pt\kern4pt\hbox{.}\hbox{.}\hbox{.}}}}%
  \def\ddots{\mathinner{\mkern1mu\raise6.5pt\hbox{.}\mkern2mu\raise3.6pt\hbox{.}\mkern2mu\raise0.7pt\hbox{.}\mkern1mu}}%
  \def\triangle{\mathord{\scalebox{0.9}{$\symup{\Delta}$}}}%
  \def\nabla{\mathord{\raisebox{\depth}{\scalebox{1}[-1]{$\symup{\Delta}$}}}}%
  \def\checkmark{\popcheck{popdark}}%
  \def\star{\mathbin{\ast}}\def\bigstar{\mathord{\ast}}%
  \def\diamond{\mathbin{\scalebox{0.6}{$\Diamond$}}}%
  \def\bigcup{\mathop{\mathchoice{\raisebox{-0.3em}{\scalebox{1.7}{$\cup$}}}{\scalebox{1.25}{$\cup$}}{\cup}{\cup}}\displaylimits}%
  \def\bigcap{\mathop{\mathchoice{\raisebox{-0.3em}{\scalebox{1.7}{$\cap$}}}{\scalebox{1.25}{$\cap$}}{\cap}{\cap}}\displaylimits}%
  \def\lesseqqgtr{\mathrel{\lessgtr}}\def\bigoplus{\mathop{\oplus}\displaylimits}%
}
\providecommand{\ssi}{si et seulement si} % abréviation fréquente
\AtBeginDocument{\providecommand{\wideparen}{\overparen}} % arc de cercle

%% FIGFIX-BEGIN v4 — correctifs globaux des figures (docs/onbuch-generation/tools/figfix)
\usepackage{adjustbox}\usepackage{etoolbox}
% toute figure TikZ/pgfplots plus large que la ligne est réduite à la largeur disponible
\BeforeBeginEnvironment{tikzpicture}{\begin{adjustbox}{max width=\linewidth}}
\AfterEndEnvironment{tikzpicture}{\end{adjustbox}}
% légendes pgfplots lisibles par-dessus les courbes
\pgfplotsset{every axis/.append style={legend style={fill=white,fill opacity=0.92,text opacity=1,draw=black!15,inner sep=3pt}}}
% étiquettes d'arêtes (syntaxe edge["..."]) avec fond blanc
\tikzset{every edge quotes/.append style={fill=white,inner sep=1.2pt}}
%% FIGFIX-END
\begin{document}

\pagegarde{Lutte contre le VIH/Sida}{SVT}{1ère D}{Module 2 : L'Éducation à la Santé}{ND08}{10 h}{Cette leçon étudie le VIH, agent du Sida : sa structure, son cycle d'infection des lymphocytes CD4, ses modes de transmission et les moyens de prévention. Elle décrit l'évolution de l'infection, le dépistage et la prise en charge par les antirétroviraux, afin d'armer l'élève pour la prévention et la lutte contre la stigmatisation.
\vspace{4pt}
\begin{multicols}{2}\small
\begin{itemize}
\item À la fin de cette leçon, je sais décrire la structure du VIH et identifier ses composants essentiels
\item À la fin de cette leçon, je sais expliquer les étapes du cycle d'infection d'un lymphocyte CD4 par le VIH
\item À la fin de cette leçon, je sais citer les modes de transmission du VIH et distinguer les situations à risque des situations sans risque
\item À la fin de cette leçon, je sais proposer les moyens de prévention adaptés à chaque mode de transmission
\item À la fin de cette leçon, je sais décrire les trois phases de l'évolution de l'infection à VIH
\item À la fin de cette leçon, je sais interpréter un graphique d'évolution de la charge virale et des lymphocytes CD4
\end{itemize}\end{multicols}}

\begin{sommaire}
\begin{multicols}{2}
\begin{enumerate}
\item Le VIH : un virus particulier
\item Le cycle d'infection du lymphocyte CD4
\item Transmission et prévention
\item Évolution de l'infection à VIH
\item Dépistage, prise en charge et traitements
\item Activité d'intégration
\item Méthodes et exercices résolus
\item Exercices
\item Corrigés des exercices
\item Fiche bilan
\end{enumerate}
\end{multicols}
\end{sommaire}

\begin{objectifs}
\begin{itemize}
\item À la fin de cette leçon, je sais décrire la structure du VIH et identifier ses composants essentiels
\item À la fin de cette leçon, je sais expliquer les étapes du cycle d'infection d'un lymphocyte CD4 par le VIH
\item À la fin de cette leçon, je sais citer les modes de transmission du VIH et distinguer les situations à risque des situations sans risque
\item À la fin de cette leçon, je sais proposer les moyens de prévention adaptés à chaque mode de transmission
\item À la fin de cette leçon, je sais décrire les trois phases de l'évolution de l'infection à VIH
\item À la fin de cette leçon, je sais interpréter un graphique d'évolution de la charge virale et des lymphocytes CD4
\item À la fin de cette leçon, je sais expliquer le principe du dépistage volontaire et son importance
\item À la fin de cette leçon, je sais argumenter en faveur de l'observance du traitement antirétroviral et lutter contre la stigmatisation
\end{itemize}
\end{objectifs}

\begin{prerequis}
\begin{itemize}
\item Le système immunitaire et les lymphocytes (notions de Seconde)
\item La notion de virus et d'agent pathogène
\item Les réactions immunitaires : anticorps et immunité cellulaire
\item La lecture et l'interprétation de graphiques à deux courbes
\item Les notions de reproduction et d'anatomie de l'appareil génital (début d'année)
\end{itemize}
\end{prerequis}

\newpage

\coursec{Le VIH : un virus particulier}

\textbf{Situation d'accroche.} Dans un centre de santé de Yaoundé, Amina, 19 ans, vient d'apprendre qu'elle est \cle{séropositive} après un dépistage volontaire. Effondrée, elle croit être condamnée à mort. Pourtant, l'infirmier lui explique qu'avec un traitement bien suivi, elle peut vivre longtemps et en bonne santé. Comment un virus invisible peut-il détruire nos défenses ? Pourquoi dit-on qu'on peut être porteur du VIH sans être malade ? Et comment le traitement permet-il de rendre le virus indétectable ?

Pour répondre à ces questions, il faut d'abord connaître l'ennemi : qu'est-ce que le VIH ? De quoi est-il fait ? Pourquoi s'attaque-t-il précisément aux cellules qui commandent notre immunité ? C'est l'objet de cette première partie.

\courssub{VIH et Sida : deux notions à ne pas confondre}

\textbf{Observation.} Dans les campagnes de sensibilisation au Cameroun, on entend souvent : « il a le Sida » pour désigner toute personne séropositive. Or, à l'hôpital, les médecins distinguent soigneusement deux situations : une personne peut être infectée par le virus et se porter bien pendant des années, tandis qu'une autre développe des maladies graves répétées (tuberculose, pneumonies, diarrhées chroniques). Le virus est le même, mais le stade de l'infection est différent.

\textbf{Hypothèse.} L'infection par le virus et la maladie ne sont pas synonymes : la maladie n'apparaîtrait qu'après une longue évolution de l'infection.

\textbf{Données.} Le suivi de cohortes de personnes séropositives non traitées montre qu'entre la contamination et l'apparition des maladies graves, il s'écoule en moyenne 8 à 10 ans, pendant lesquels la personne ne présente aucun symptôme. Ce n'est que lorsque les défenses immunitaires sont détruites que surviennent les \cle{maladies opportunistes}, c'est-à-dire des infections causées par des microbes habituellement inoffensifs, qui « profitent de l'occasion » offerte par l'affaiblissement de l'immunité.

\textbf{Conclusion.} Il faut distinguer rigoureusement le virus (VIH), l'infection (état de séropositivité) et la maladie (Sida).

\begin{definition}[Le VIH]
Le \cle{VIH} (\emph{Virus de l'Immunodéficience Humaine}) est un \cle{rétrovirus} qui infecte préférentiellement les \cle{lymphocytes T4} (ou lymphocytes CD4) et détruit progressivement le système immunitaire. Une personne qui porte le VIH dans son organisme est dite \cle{séropositive}.
\end{definition}

\begin{definition}[Le Sida]
Le \cle{Sida} (\emph{Syndrome d'ImmunoDéficience Acquise}) est le \textbf{stade ultime} de l'infection à VIH. Il est caractérisé par l'effondrement des défenses immunitaires (chute drastique des lymphocytes CD4, en dessous de 200 cellules par mm$^3$ de sang) et par l'apparition de maladies opportunistes (tuberculose, pneumocystose, candidoses, certaines tumeurs comme la maladie de Kaposi).
\end{definition}

On retiendra la logique suivante : le VIH est la \textbf{cause}, l'infection à VIH est l'\textbf{état} de la personne contaminée, et le Sida est la \textbf{conséquence terminale} de l'infection lorsqu'elle n'est pas traitée.

\begin{attention}[Erreur fréquente : confondre VIH et Sida]
On peut être \textbf{séropositif} (porteur du VIH) \textbf{sans être malade du Sida}. Dire « Amina a le Sida » parce qu'elle est séropositive est une erreur scientifique \emph{et} une source de stigmatisation. Avec les traitements actuels, une personne séropositive correctement suivie peut ne jamais atteindre le stade Sida. Au Probatoire, emploie toujours le terme exact : « personne infectée par le VIH » ou « personne séropositive », et réserve « Sida » au stade de maladie déclarée.
\end{attention}

\courssub{La structure du VIH : un virus enveloppé à ARN}

\textbf{Observation.} Au microscope électronique, le VIH apparaît comme une particule sphérique d'environ 100 à 120 nanomètres de diamètre, soit environ 100 fois plus petite qu'une cellule humaine. Comme tous les virus, il est \textbf{acellulaire} : il n'a ni cytoplasme, ni métabolisme propre, et il est incapable de se multiplier seul. C'est un \cle{parasite intracellulaire obligatoire} : pour se reproduire, il doit obligatoirement pénétrer dans une cellule vivante et détourner sa machinerie.

L'analyse biochimique du virus révèle une organisation en couches concentriques, de l'extérieur vers l'intérieur.

\begin{aretenir}[Structure du VIH]
Le VIH est constitué, de l'extérieur vers l'intérieur, de :
\begin{itemize}
\item une \cle{enveloppe} lipidique (dérivée de la membrane de la cellule hôte) dans laquelle sont insérées des \cle{glycoprotéines} : la \textbf{gp120}, qui forme des « clés » de reconnaissance à la surface du virus, et la \textbf{gp41}, qui ancre la gp120 dans l'enveloppe et permet la fusion avec la cellule cible ;
\item une \cle{capside} protéique (coque interne) qui protège le matériel génétique ;
\item \textbf{deux molécules d'\cle{ARN}} (génome viral, en deux exemplaires identiques) ;
\item des \cle{enzymes} virales : la \textbf{transcriptase inverse}, l'\textbf{intégrase} et la \textbf{protéase}, indispensables à la multiplication du virus dans la cellule.
\end{itemize}
\end{aretenir}

\begin{popfigure}
\begin{center}
\begin{tikzpicture}[scale=1.05]
% Enveloppe
\fill[popblueL] (0,0) circle (2.6);
\fill[popcream] (0,0) circle (2.35);
% Capside
\fill[poporangeL] (0,0) circle (1.5);
\fill[white] (0,0) circle (1.25);
% ARN (deux brins ondulés)
\draw[popdark,thick,decorate,decoration={coil,segment length=4pt,amplitude=2.5pt}] (-0.9,0.5) -- (0.9,-0.5);
\draw[popdark,thick,decorate,decoration={coil,segment length=4pt,amplitude=2.5pt}] (-0.9,-0.5) -- (0.9,0.5);
% Transcriptase inverse (petits disques)
\fill[poppurple] (0.1,0.75) circle (0.16);
\fill[poppurple] (-0.2,-0.8) circle (0.16);
% gp120 (boutons de surface) + gp41 (tiges)
\foreach \a in {20,70,120,170,220,270,320}{
  \draw[popgreen,very thick] (\a:2.6) -- (\a:3.0);
  \fill[popgreen] (\a:3.15) circle (0.16);
}
% Légendes pointées
\draw[thin,popmuted] (3.15,1.0) -- (5.3,1.6) node[right,align=left]{\textbf{gp120} : glycoprotéine de\\ reconnaissance (la « clé »)};
\draw[thin,popmuted] (2.75,2.05) -- (4.6,3.0) node[right,align=left]{\textbf{gp41} : glycoprotéine\\ d'ancrage et de fusion};
\draw[thin,popmuted] (-2.3,1.6) -- (-4.6,2.6) node[left,align=right]{\textbf{Enveloppe} lipidique\\ (issue de la cellule hôte)};
\draw[thin,popmuted] (-1.4,-1.35) -- (-4.4,-2.4) node[left,align=right]{\textbf{Capside} protéique};
\draw[thin,popmuted] (0.6,-0.35) -- (4.9,-1.3) node[right,align=left]{\textbf{2 molécules d'ARN}\\ (génome viral)};
\draw[thin,popmuted] (0.1,0.75) -- (2.4,2.4) node[right,align=left]{\textbf{Transcriptase inverse}\\ (enzyme virale)};
\end{tikzpicture}
\end{center}
\legende{Schéma de la structure du VIH. Le virus possède une enveloppe portant les glycoprotéines gp120 et gp41, une capside protéique renfermant deux molécules d'ARN et des enzymes (transcriptase inverse, intégrase, protéase).}
\end{popfigure}

\textbf{Pourquoi ces détails sont-ils importants ?} Chaque élément de structure correspond à une fonction, et donc à une cible thérapeutique : la gp120 permet au virus de reconnaître sa cellule cible ; la transcriptase inverse permet de copier son ARN en ADN ; la protéase découpe les protéines virales lors de l'assemblage des nouveaux virus. Les médicaments antirétroviraux (étudiés dans la section 5) bloquent précisément ces étapes.

\courssub{Un rétrovirus : un génome à ARN qui fonctionne « à l'envers »}

Dans toutes les cellules vivantes, l'information génétique circule dans le sens : ADN $\rightarrow$ ARN $\rightarrow$ protéines. Le VIH, lui, possède un génome constitué d'\textbf{ARN}, et non d'ADN. Pour s'intégrer dans le patrimoine génétique de la cellule hôte (qui est de l'ADN), il doit d'abord fabriquer de l'ADN à partir de son ARN : c'est le sens \emph{inverse} du flux habituel.

\begin{definition}[Rétrovirus]
Un \cle{rétrovirus} est un virus dont le génome est constitué d'ARN et qui possède une enzyme, la \cle{transcriptase inverse}, capable de synthétiser de l'ADN à partir de cet ARN (étape appelée \emph{rétrotranscription}). L'ADN viral ainsi produit est ensuite inséré dans le génome de la cellule hôte grâce à l'\cle{intégrase}.
\end{definition}

Cette particularité a deux conséquences majeures :
\begin{itemize}
\item le virus devient un \textbf{provirus} intégré au cœur même du génome de la cellule : il est quasiment impossible de l'en déloger, ce qui explique qu'on ne guérisse pas de l'infection à VIH, mais qu'on la contrôle ;
\item la transcriptase inverse fait de nombreuses erreurs de copie : le virus \textbf{mute} très facilement, ce qui complique la fabrication d'un vaccin et impose des traitements combinant plusieurs molécules.
\end{itemize}

\courssub{La cellule cible : le lymphocyte T4 (CD4), chef d'orchestre de l'immunité}

\textbf{Situation concrète.} Imagine un orchestre sans chef : chaque musicien joue, mais sans coordination, le résultat est une cacophonie. Dans notre organisme, la réponse immunitaire fonctionne comme un orchestre, et le \cle{lymphocyte T4} en est le chef.

\begin{aretenir}[Rôle du lymphocyte T4 (CD4)]
Le \cle{lymphocyte T4}, appelé aussi lymphocyte \cle{CD4} (car il porte à sa surface la molécule CD4), est un globule blanc qui \textbf{coordonne la réponse immunitaire} : il active les lymphocytes B (producteurs d'anticorps), les lymphocytes T8 (tueurs de cellules infectées) et les macrophages. Sans lui, toutes les autres défenses sont désorganisées.
\end{aretenir}

Le VIH possède à sa surface la glycoprotéine \textbf{gp120}, qui reconnaît spécifiquement la molécule CD4, comme une clé reconnaît sa serrure. Le virus infecte donc préférentiellement les lymphocytes T4 (ainsi que les macrophages et les cellules dendritiques, qui portent aussi CD4). En détruisant progressivement les chefs d'orchestre, le virus désorganise toute la réponse immunitaire : c'est le mécanisme profond de l'immunodéficience.

\begin{popfigure}
\begin{center}
\begin{tikzpicture}[scale=0.95]
% Cellule CD4
\fill[popgreenL] (0,0) circle (2.1);
\fill[popgoldL] (0,0) circle (0.7);
\draw[popdark] (0,0) circle (0.7);
% Récepteurs CD4 (fourchettes en surface)
\foreach \a in {30,90,150,210,270,330}{
  \draw[popdark,very thick] (\a:2.1) -- (\a:2.55);
  \draw[popdark,very thick] (\a:2.55) -- (\a-12:2.85);
  \draw[popdark,very thick] (\a:2.55) -- (\a+12:2.85);
}
% Virus à côté
\begin{scope}[xshift=5.6cm,yshift=1.2cm,scale=0.55]
\fill[popblueL] (0,0) circle (2.6);
\fill[popcream] (0,0) circle (2.35);
\fill[poporangeL] (0,0) circle (1.5);
\fill[white] (0,0) circle (1.25);
\draw[popdark,thick] (-0.8,0.4) -- (0.8,-0.4);
\draw[popdark,thick] (-0.8,-0.4) -- (0.8,0.4);
\foreach \a in {20,80,140,200,260,320}{
  \draw[popgreen,very thick] (\a:2.6) -- (\a:3.0);
  \fill[popgreen] (\a:3.15) circle (0.16);
}
\end{scope}
% Flèche d'interaction
\draw[-{Stealth},popink,very thick] (4.1,0.6) -- (2.6,0.4) node[midway,above,align=center,font=\small]{la gp120 reconnaît\\ le récepteur CD4};
% Légendes
\draw[thin,popmuted] (-1.5,1.6) -- (-3.6,2.4) node[left,align=right]{\textbf{Récepteur CD4}\\ (la « serrure »)};
\draw[thin,popmuted] (0.0,-0.7) -- (-2.6,-2.6) node[left,align=right]{\textbf{Noyau} du lymphocyte\\ (ADN de la cellule)};
\node[align=center] at (0,-3.3) {\textbf{Lymphocyte T4 (CD4)}\\ cellule cible du VIH};
\node[align=center] at (5.6,-2.2) {\textbf{VIH}\\ (enveloppe + gp120 + ARN)};
\end{tikzpicture}
\end{center}
\legende{Reconnaissance entre le VIH et sa cellule cible. La glycoprotéine gp120 du virus (la « clé ») se fixe spécifiquement sur le récepteur CD4 (la « serrure ») à la surface du lymphocyte T4. Cette spécificité explique pourquoi le VIH détruit préférentiellement les cellules qui coordonnent l'immunité.}
\end{popfigure}

\begin{exemplebox}[Pourquoi une simple grippe peut devenir grave chez un patient au stade Sida]
Chez une personne saine, un virus banal est rapidement éliminé : les lymphocytes T4 détectent l'agression, activent les lymphocytes B qui produisent des anticorps et les T8 qui détruisent les cellules infectées. Chez une personne dont les CD4 sont détruits par le VIH, cette coordination n'existe plus : le microbe banal se multiplie sans frein. C'est exactement ce qui se passe avec la tuberculose, première cause de décès des personnes vivant avec le VIH en Afrique subsaharienne, y compris au Cameroun.
\end{exemplebox}

\courssub{Les types de VIH : VIH-1 et VIH-2}

Il existe deux types de virus de l'immunodéficience humaine.

\begin{center}
\begin{tabularx}{\linewidth}{|l|X|X|}\hline
\rowcolor{popblueL} \textbf{Caractère} & \textbf{VIH-1} & \textbf{VIH-2} \\ \hline
Répartition & Mondiale ; \textbf{majoritaire au Cameroun} et en Afrique centrale & Surtout Afrique de l'Ouest \\ \hline
Virulence & Élevée : évolution rapide vers le Sida sans traitement & Plus faible : évolution plus lente \\ \hline
Transmission & Plus efficace & Moins efficace \\ \hline
\end{tabularx}
\end{center}

Au Cameroun, l'écrasante majorité des infections est due au \cle{VIH-1}, qui présente en outre une grande diversité de sous-types (groupes et sous-types génétiques), héritée de l'histoire ancienne de l'épidémie en Afrique centrale. Les tests de dépistage utilisés dans les centres de santé camerounais détectent les deux types.

\begin{savaistu}[D'où vient le VIH ?]
Les travaux des épidémiologistes et des généticiens montrent que le VIH-1 provient d'un virus des singes (le VIS, virus de l'immunodéficience simienne) présent chez les chimpanzés d'Afrique centrale, notamment dans les forêts du sud Cameroun. Le passage à l'Homme se serait produit au début du XX\textsuperscript{e} siècle, probablement lors de contacts avec du sang de gibier. La comparaison des génomes viraux permet de dater et de localiser ce passage : c'est un bel exemple de ce que la science peut reconstituer a posteriori.
\end{savaistu}

\courssub{Un virus fragile dans le milieu extérieur}

\textbf{Observation.} Les enquêtes épidémiologiques sont formelles : aucun cas de contamination n'a jamais été observé par poignée de main, embrassade, partage de repas, piqûre de moustique ou utilisation commune des toilettes. Pourtant, la peur du contact quotidien reste vive et alimente le rejet des personnes séropositives.

\textbf{Interprétation.} Cette absence de transmission s'explique par la structure même du virus : son \textbf{enveloppe lipidique} est très fragile. À l'air libre, la chaleur, le dessèchement et les savons détruisent rapidement cette enveloppe ; privé de ses glycoprotéines gp120, le virus ne peut plus reconnaître ni infecter aucune cellule.

\begin{aretenir}[Fragilité du VIH : conséquences]
Le VIH ne survit que \textbf{quelques minutes à l'air libre} et est détruit par l'eau de Javel, le savon et la chaleur. Il ne se transmet donc \textbf{pas} par les contacts de la vie quotidienne. Sa transmission exige un \textbf{contact direct} entre un liquide biologique contaminé (sang, sécrétions sexuelles, lait maternel) et le sang ou les muqueuses d'une autre personne (ces modes de transmission seront détaillés dans la section 3).
\end{aretenir}

\begin{savaistu}[On ne peut pas attraper le VIH en serrant la main]
Le VIH ne survit que quelques minutes à l'air libre. On ne peut pas le contracter en serrant la main, en partageant un repas, en buvant au même verre, en dormant sous la même moustiquaire, ni par une piqûre de moustique : le moustique ne réinjecte jamais le sang d'un repas précédent, et le virus ne survit pas dans son organisme. Combattre ces idées fausses, c'est combattre la discrimination qui isole les personnes séropositives et les décourage de se faire dépister.
\end{savaistu}

\courssub{Exercice résolu et bilan de la section}

\begin{exoresolu}[QCM de vérification des acquis]
\textbf{Énoncé.} Pour chaque affirmation, indique si elle est vraie ou fausse, en justifiant brièvement.
\begin{enumerate}
\item Une personne séropositive est forcément malade du Sida.
\item Le génome du VIH est constitué de deux molécules d'ADN.
\item La gp120 permet au VIH de reconnaître le récepteur CD4 des lymphocytes T4.
\item Le VIH peut se transmettre par une piqûre de moustique.
\item La transcriptase inverse synthétise de l'ADN à partir de l'ARN viral.
\end{enumerate}
\tcblower
\textbf{Solution détaillée.}
\begin{enumerate}
\item \textbf{Faux.} Être séropositif signifie être porteur du VIH. Le Sida n'est que le stade ultime de l'infection, atteint après des années d'évolution sans traitement. Avec un traitement antirétroviral bien suivi, une personne séropositive peut ne jamais développer le Sida.
\item \textbf{Faux.} Le VIH est un rétrovirus : son génome est constitué de \textbf{deux molécules d'ARN}, et non d'ADN. C'est la transcriptase inverse qui fabriquera de l'ADN à partir de cet ARN dans la cellule hôte.
\item \textbf{Vrai.} La gp120, glycoprotéine de l'enveloppe virale, se fixe spécifiquement sur le récepteur CD4 : c'est l'étape de reconnaissance (la « clé » et la « serrure »), première étape de l'infection.
\item \textbf{Faux.} Le VIH est fragile et ne survit pas dans l'organisme du moustique ; de plus, le moustique ne réinjecte pas le sang prélevé sur un repas précédent. Aucun cas de transmission par moustique n'a jamais été observé.
\item \textbf{Vrai.} C'est précisément la définition de la transcriptase inverse (ou « transcriptase réverse ») : elle réalise la rétrotranscription, c'est-à-dire la synthèse d'ADN à partir d'ARN, étape caractéristique des rétrovirus.
\end{enumerate}
\end{exoresolu}

\begin{aretenir}[Bilan de la section]
\begin{itemize}
\item Le \textbf{VIH} est un rétrovirus enveloppé (gp120, gp41), à génome formé de deux molécules d'ARN, accompagné d'enzymes (transcriptase inverse, intégrase, protéase).
\item Il infecte préférentiellement les \textbf{lymphocytes T4 (CD4)}, chefs d'orchestre de la réponse immunitaire, qu'il détruit progressivement.
\item Le \textbf{Sida} est le stade ultime de l'infection : on peut être séropositif sans être malade.
\item Le \textbf{VIH-1} est le type majoritaire au Cameroun ; le VIH-2, moins virulent, domine en Afrique de l'Ouest.
\item Le virus est \textbf{fragile} dans le milieu extérieur : il ne se transmet pas par les contacts de la vie quotidienne.
\end{itemize}
Maintenant que nous connaissons l'ennemi, la prochaine section expliquera comment il pénètre dans le lymphocyte CD4 et s'y multiplie : le cycle d'infection.
\end{aretenir}

\coursec{Le cycle d'infection du lymphocyte CD4}

Nous avons vu dans la section précédente que le VIH est un \cle{rétrovirus} : son génome est constitué de deux brins d'\cle{ARN} associés à des enzymes, le tout enfermé dans une enveloppe portant des glycoprotéines \cle{gp120}. Mais une question essentielle reste posée : comment ce minuscule virus, incapable de se multiplier seul, parvient-il à détruire tout un système immunitaire ? La réponse tient dans un mécanisme remarquable : le VIH \emph{détourne} la machinerie de nos cellules de défense, les \cle{lymphocytes T4} (ou \cle{CD4}), pour les transformer en véritables « usines à virus ». Comprendre ce cycle, étape par étape, c'est comprendre à la fois pourquoi l'infection est si redoutable et où la médecine peut intervenir pour la bloquer.

\courssub{Observation préliminaire : une cellule qui fabrique le virus qui la tue}

\begin{experience}[Observation au microscope électronique]
Des lymphocytes CD4 prélevés chez des personnes infectées et observés en microscopie électronique à transmission montrent des images frappantes :
\begin{itemize}
\item des particules virales \emph{fixées} à la surface de la membrane cellulaire ;
\item des particules virales \emph{en train de bourgeonner} à la surface de la cellule, comme de petites vésicules qui se détachent ;
\item des cellules dont le noyau contient de l'ADN viral \emph{intégré} au chromosome (mis en évidence par hybridation moléculaire).
\end{itemize}
\textbf{Interprétation :} le virus ne se contente pas de traverser la cellule : il s'y \emph{installe}, s'y \emph{multiplie} et en \emph{ressort}. Le lymphocyte est devenu une usine de production virale.
\end{experience}

\begin{methode}[Décrire un cycle viral : la bonne démarche]
Pour décrire correctement le cycle de multiplication d'un virus, suis toujours le même ordre logique, quelle que soit le virus étudié :
\begin{enumerate}
\item \textbf{Fixation} : le virus reconnaît et s'accroche à sa cellule cible ;
\item \textbf{Entrée} : le virus (ou son génome) pénètre dans la cellule ;
\item \textbf{Réplication} : le génome viral est copié et s'installe dans la cellule ;
\item \textbf{Assemblage} : la cellule fabrique les composants de nouveaux virus ;
\item \textbf{Libération} : les nouveaux virus sortent et vont infecter d'autres cellules.
\end{enumerate}
Cette trame te servira pour le VIH, mais aussi pour tout autre virus étudié au lycée.
\end{methode}

\courssub{Étape 1 : la fixation — une clé qui reconnaît sa serrure}

Tout commence par une rencontre. Le VIH circule dans le sang ou les sécrétions ; à la surface des lymphocytes T4 se trouve une molécule particulière, le \cle{récepteur CD4}, qui sert normalement à la reconnaissance des antigènes. La glycoprotéine \cle{gp120} de l'enveloppe virale possède une forme complémentaire de ce récepteur : elle s'y fixe avec une grande précision, comme une clé dans une serrure.

Mais la fixation sur le CD4 seul ne suffit pas. Elle provoque un changement de conformation de la gp120, qui peut alors se lier à un second récepteur membranaire, le \cle{corécepteur} (le plus souvent \textbf{CCR5} ou \textbf{CXCR4}). C'est cette double fixation — récepteur CD4 \emph{plus} corécepteur — qui déclenche l'étape suivante.

\begin{attention}[Spécificité d'hôte et de cellule]
Cette reconnaissance très précise explique deux faits importants :
\begin{itemize}
\item le VIH n'infecte que les cellules portant le récepteur CD4 (lymphocytes T4, mais aussi macrophages et cellules dendritiques) : c'est un virus \emph{lymphotrope} ;
\item le VIH humain n'infecte pratiquement que l'Homme : on ne peut pas « attraper le Sida » d'un animal ou d'un objet.
\end{itemize}
Ne dis jamais que « le VIH attaque toutes les cellules » : il cible spécifiquement les cellules à CD4.
\end{attention}

\begin{savaistu}[Une mutation qui protège]
Certaines personnes, rares, possèdent une mutation du gène codant le corécepteur CCR5 (mutation dite « delta 32 ») : leurs lymphocytes ne portent pas de CCR5 fonctionnel, et le VIH peine à y entrer. Ces personnes sont naturellement très résistantes à l'infection. C'est en s'inspirant de ce phénomène que des chercheurs tentent de bloquer l'entrée du virus par des médicaments « inhibiteurs d'entrée ».
\end{savaistu}

\courssub{Étape 2 : la pénétration — le virus livre son matériel}

La double fixation rapproche l'enveloppe virale de la membrane cellulaire. Une autre glycoprotéine virale, la \cle{gp41}, provoque alors la \textbf{fusion} des deux membranes : l'enveloppe du virus fusionne avec la membrane du lymphocyte. La capside virale est ainsi déversée dans le \emph{cytoplasme} de la cellule, puis se désassemble : c'est la \textbf{déscapsidation}.

Le contenu livré est précis : \textbf{deux molécules d'ARN viral} (le génome) et \textbf{trois enzymes} essentielles : la \cle{transcriptase inverse}, l'\cle{intégrase} et la \cle{protéase}. Le virus n'apporte rien d'autre : pour tout le reste (énergie, ribosomes, nucléotides), il devra utiliser les ressources de la cellule hôte.

\courssub{Étape 3 : la transcription inverse — l'étape clé}

Nous voici au cœur du problème. Nos cellules fonctionnent selon le « dogme central » de la biologie moléculaire : l'information génétique va de l'\textbf{ADN} vers l'\textbf{ARN}, puis vers les protéines. Or le VIH possède un génome à ARN. Comment cet ARN pourrait-il s'installer durablement dans une cellule dont le génome est de l'ADN ?

C'est là qu'intervient l'enzyme apportée par le virus : la \cle{transcriptase inverse} (ou ADN polymérase ARN-dépendante). Dans le cytoplasme du lymphocyte, elle réalise l'opération \emph{inverse} de la transcription classique :
\begin{enumerate}
\item elle utilise l'ARN viral comme \emph{matrice} et synthétise un premier brin d'ADN complémentaire (on obtient un hybride ARN--ADN) ;
\item elle dégrade le brin d'ARN initial (activité RNase H) ;
\item elle synthétise le second brin d'ADN, complémentaire du premier.
\end{enumerate}
Résultat : une molécule d'\textbf{ADN viral double brin}, copie fidèle (en principe) du génome viral, mais sous une forme compatible avec le génome de la cellule.

\begin{propriete}[La transcriptase inverse : propriété admise]
La \cle{transcriptase inverse} permet au VIH de transformer son ARN en ADN double brin, forme intégrable dans le génome de la cellule hôte. Cette activité enzymatique (synthèse d'ADN sur matrice ARN) \textbf{n'est pas une fonction standard de nos cellules} pour l'expression de leurs gènes : nos cellules ne savent fabriquer de l'ADN qu'à partir d'ADN (réplication) ou de l'ARN qu'à partir d'ADN (transcription). Cette particularité fait de la transcriptase inverse une \textbf{cible thérapeutique privilégiée} : un médicament qui la bloque (comme l'AZT, premier antirétroviral) gêne le virus sans perturber le fonctionnement normal de nos propres cellules.
\end{propriete}

\begin{popfigure}
\begin{center}
\begin{tikzpicture}[font=\small, >=Stealth]
% ARN matrice
\draw[very thick, poporange] (0,3.2) -- (7,3.2);
\foreach \x/\b in {0.4/A, 1.0/U, 1.6/G, 2.2/C, 2.8/A, 3.4/U, 4.0/C, 4.6/G, 5.2/A, 5.8/U, 6.4/C}
  \node[poporange] at (\x,3.55) {\b};
\node[left, poporange] at (0,3.2) {\textbf{ARN viral} ($+$)};
% flèche 1
\draw[->, very thick, popblue] (3.5,2.9) -- (3.5,2.15) node[midway, right, text=popdark, align=left] {\footnotesize transcriptase inverse\\ \footnotesize (1\textsuperscript{er} brin d'ADN)};
% hybride
\draw[very thick, poporange] (0,1.9) -- (7,1.9);
\draw[very thick, popblue] (0,1.3) -- (7,1.3);
\foreach \x/\b in {0.4/T, 1.0/A, 1.6/C, 2.2/G, 2.8/T, 3.4/A, 4.0/G, 4.6/C, 5.2/T, 5.8/A, 6.4/G}
  \node[popblue] at (\x,1.05) {\b};
\node[left, poporange] at (0,1.9) {ARN};
\node[left, popblue] at (0,1.3) {ADN};
\node[popmuted] at (8.3,1.6) {\footnotesize hybride ARN--ADN};
% flèche 2
\draw[->, very thick, popblue] (3.5,1.0) -- (3.5,0.25) node[midway, right, text=popdark, align=left] {\footnotesize dégradation de l'ARN,\\ \footnotesize synthèse du 2\textsuperscript{e} brin};
% ADN double brin
\draw[very thick, popblue] (0,0.0) -- (7,0.0);
\draw[very thick, popblue] (0,-0.6) -- (7,-0.6);
\foreach \x/\b/\c in {0.4/T/A, 1.0/A/T, 1.6/C/G, 2.2/G/C, 2.8/T/A, 3.4/A/T, 4.0/G/C, 4.6/C/G, 5.2/T/A, 5.8/A/T, 6.4/G/C} {
  \node[popblue] at (\x,0.25) {\b};
  \node[popblue] at (\x,-0.35) {\c};
  \draw[popblue!60] (\x,0.0) -- (\x,-0.6);
}
\node[left, popblue] at (0,-0.3) {\textbf{ADN viral}};
\end{tikzpicture}
\end{center}
\legende{La transcription inverse : à partir de l'ARN viral (brin $+$), la transcriptase inverse fabrique un brin d'ADN complémentaire (hybride ARN--ADN), dégrade l'ARN, puis synthétise le second brin d'ADN. On obtient un ADN viral double brin, copie du génome viral, prêt pour l'intégration.}
\end{popfigure}

\begin{savaistu}[Une copieuse qui fait des fautes]
La transcriptase inverse est une enzyme « négligente » : contrairement à nos ADN polymérases, elle ne possède pas de fonction de correction. Elle commet en moyenne \textbf{une erreur par copie du génome}. Comme des milliards de virus sont produits chaque jour chez une personne non traitée, d'innombrables \emph{mutants} apparaissent en permanence. Cette \textbf{grande variabilité du VIH} explique deux difficultés majeures : l'apparition de virus résistants aux médicaments (d'où l'obligation d'associer plusieurs antirétroviraux) et l'impossibilité, à ce jour, de fabriquer un vaccin efficace — le virus « change de visage » en permanence.
\end{savaistu}

\courssub{Étape 4 : l'intégration — le virus s'installe dans le génome}

L'ADN viral double brin migre vers le \textbf{noyau} du lymphocyte et y pénètre. Là, la deuxième enzyme virale entre en scène : l'\cle{intégrase}. Elle découpe l'ADN d'un chromosome de la cellule et y insère l'ADN viral. L'ADN viral intégré prend le nom de \cle{provirus}.

Cette étape est décisive et ses conséquences sont graves :
\begin{itemize}
\item le provirus fait désormais \emph{partie du génome} de la cellule : il est répliqué avec elle à chaque division cellulaire ;
\item il est \textbf{indélogeable} : aucun traitement actuel ne peut l'extraire du chromosome. C'est pourquoi on sait \emph{contrôler} l'infection mais pas encore la \emph{guérir} ;
\item le provirus peut rester \textbf{silencieux} (latent) : la cellule infectée ne fabrique alors aucun virus et paraît normale.
\end{itemize}

\begin{attention}[Le provirus latent : le réservoir caché]
Une cellule portant un provirus peut rester « silencieuse » pendant des \textbf{années} : elle ne produit pas de virus, n'est pas repérée par le système immunitaire et n'est pas atteinte par les traitements (qui ne bloquent que la production de nouveaux virus). Ces cellules forment des \emph{réservoirs} viraux. C'est cette \cle{latence} qui explique la \textbf{phase asymptomatique} de l'infection (que nous étudierons en section 4) et l'obligation de prendre les antirétroviraux \emph{à vie} : si l'on arrête le traitement, les réservoirs relancent la production virale.
\end{attention}

\courssub{Étape 5 : la production virale — la cellule devient une usine}

Lorsque le lymphocyte infecté est activé (par exemple lors d'une infection quelconque), le provirus est « réveillé ». La machinerie cellulaire, qui ne fait pas la différence entre un gène humain et un gène viral, se met alors au service du virus :
\begin{enumerate}
\item l'ARN polymérase \emph{de la cellule} \textbf{transcrit} le provirus en nombreuses molécules d'ARN viral : certaines serviront de génomes aux futurs virus, d'autres de messagers ;
\item les \textbf{ribosomes} de la cellule traduisent ces ARN messagers en \emph{protéines virales} : protéines de capside, enzymes, précurseurs des glycoprotéines d'enveloppe ;
\item ces composants s'assemblent sous la membrane plasmique pour former de nouvelles particules immatures.
\end{enumerate}
Le lymphocyte, dont le rôle normal est de coordonner la défense immunitaire, est ainsi entièrement détourné : c'est devenu une \textbf{usine à virus}.

\courssub{Étape 6 : le bourgeonnement et la libération — la mort du lymphocyte}

Les nouvelles particules virales traversent la membrane plasmique en emportant un morceau de celle-ci : c'est le \cle{bourgeonnement}, qui leur donne leur enveloppe. La troisième enzyme virale, la \cle{protéase}, découpe alors les longues chaînes protéiques immatures en protéines fonctionnelles : le virus devient \textbf{mature} et infectieux.

Chaque cellule infectée peut libérer des \textbf{centaines de nouveaux virions}, qui iront infecter d'autres lymphocytes CD4 : le cycle recommence, de façon exponentielle. Quant au lymphocyte producteur, il est détruit — lyse cellulaire, épuisement ou destruction par les lymphocytes T8 (cytotoxiques) qui reconnaissent les cellules infectées.

\begin{popfigure}
\begin{center}
\begin{tikzpicture}[font=\footnotesize, >=Stealth, node distance=0.45cm]
% Cellule centrale
\draw[fill=popblueL, draw=popblue, thick] (0,0) circle (1.5);
\draw[fill=poppinkL, draw=poppink, thick] (0,0) circle (0.55);
\node at (0,0) {\textbf{noyau}};
\draw[very thick, poppurple] (-0.35,0.12) -- (0.35,0.12);
\node[poppurple] at (0,0.38) {\scriptsize provirus};
\node[popmuted] at (0,-1.9) {\textbf{Lymphocyte CD4}};
% Virus extérieur
\draw[fill=poporangeL, draw=poporange, thick] (-5.2,1.4) circle (0.4);
\node[poporange] at (-5.2,1.4) {\scriptsize VIH};
\draw[poporange, thick] (-5.2,1.85) -- (-5.2,2.05); \draw[poporange, thick] (-5.6,1.2) -- (-5.75,1.1); \draw[poporange, thick] (-4.8,1.2) -- (-4.65,1.1);
% Étapes
\node[draw=popgreen, fill=popgreenL, rounded corners, align=center, text width=2.6cm] (e1) at (-3.4,3.4) {\textbf{1. Fixation}\\ gp120 $\to$ CD4 + corécepteur};
\node[draw=popgreen, fill=popgreenL, rounded corners, align=center, text width=2.6cm] (e2) at (0.4,3.9) {\textbf{2. Entrée}\\ fusion, ARN viral libéré};
\node[draw=popgreen, fill=popgreenL, rounded corners, align=center, text width=2.8cm] (e3) at (3.9,3.0) {\textbf{3. Transcription inverse}\\ ARN $\to$ ADN viral};
\node[draw=popgreen, fill=popgreenL, rounded corners, align=center, text width=2.8cm] (e4) at (4.4,-0.2) {\textbf{4. Intégration}\\ provirus dans le chromosome};
\node[draw=popgreen, fill=popgreenL, rounded corners, align=center, text width=2.8cm] (e5) at (2.6,-3.2) {\textbf{5. Production}\\ ARN et protéines virales};
\node[draw=popgreen, fill=popgreenL, rounded corners, align=center, text width=2.8cm] (e6) at (-1.6,-3.6) {\textbf{6. Bourgeonnement}\\ libération, mort du CD4};
% Flèches du cycle
\draw[->, very thick, popdark] (e1) -- (e2);
\draw[->, very thick, popdark] (e2) -- (e3);
\draw[->, very thick, popdark] (e3) -- (e4);
\draw[->, very thick, popdark] (e4) -- (e5);
\draw[->, very thick, popdark] (e5) -- (e6);
% Liaisons vers cellule
\draw[->, thick, poporange] (-4.8,1.4) -- (-1.6,0.9);
\draw[->, thick, popblue] (e6.west) .. controls (-4.6,-2.6) .. (-4.9,0.9);
% Nouveaux virions
\draw[fill=poporangeL, draw=poporange, thick] (-4.6,-2.2) circle (0.22);
\draw[fill=poporangeL, draw=poporange, thick] (-5.3,-1.7) circle (0.22);
\draw[fill=poporangeL, draw=poporange, thick] (-5.0,-2.9) circle (0.22);
\node[popmuted, align=center] at (-5.0,-3.6) {\scriptsize nouveaux virions\\ \scriptsize (infectent d'autres CD4)};
\end{tikzpicture}
\end{center}
\legende{Les six étapes du cycle d'infection du lymphocyte CD4 par le VIH : fixation, entrée, transcription inverse, intégration du provirus, production de nouveaux virus, bourgeonnement et libération. La cellule infectée finit par être détruite.}
\end{popfigure}

\courssub{Conséquence : la destruction progressive des lymphocytes CD4}

Ce cycle, répété des milliards de fois dans l'organisme, a une conséquence implacable : le nombre de lymphocytes CD4 \textbf{diminue progressivement}. Or les CD4 sont les « chefs d'orchestre » de l'immunité : ce sont eux qui activent les lymphocytes B (production d'anticorps) et les lymphocytes T8 (destruction des cellules infectées). Sans eux, l'organisme devient incapable de se défendre, même contre des microbes habituellement bénins : c'est l'\cle{immunodépression}, qui ouvre la porte aux infections opportunistes caractéristiques du Sida.

\begin{exoresolu}[Compter les étapes du cycle]
\textbf{Énoncé.} Un médicament antirétroviral bloque l'intégrase du VIH. (a) À quelle étape du cycle viral ce médicament agit-il ? (b) Une cellule traitée par ce médicament peut-elle produire de nouveaux virus ? Justifie. (c) Ce médicament peut-il débarrasser l'organisme des provirus déjà intégrés avant le traitement ?
\tcblower
\textbf{Solution.} (a) L'intégrase intervient à l'\textbf{étape 4} : l'intégration de l'ADN viral dans le chromosome de la cellule hôte. Le bloquer empêche la formation du provirus. (b) \textbf{Non} : sans intégration, l'ADN viral ne peut pas être transcrit par la machinerie cellulaire ; la cellule ne devient donc pas une usine à virus et ne produit pas de nouveaux virions. (c) \textbf{Non} : le médicament empêche les \emph{nouvelles} intégrations, mais il est sans effet sur les provirus \emph{déjà} intégrés dans les chromosomes des cellules infectées avant le traitement (réservoirs). C'est pourquoi les traitements actuels contrôlent l'infection sans la guérir.
\end{exoresolu}

\begin{aretenir}[L'essentiel du cycle d'infection]
\begin{itemize}
\item Le VIH se fixe sur le récepteur \cle{CD4} et un corécepteur (CCR5/CXCR4) grâce à sa gp120, puis entre dans le lymphocyte par fusion (gp41).
\item La \cle{transcriptase inverse} transforme l'ARN viral en ADN double brin ; l'\cle{intégrase} insère cet ADN dans le chromosome : c'est le \cle{provirus}.
\item Le provirus peut rester latent des années, puis commander à la cellule la production de nouveaux virus, libérés par \cle{bourgeonnement}.
\item Le lymphocyte producteur est détruit : la répétition du cycle entraîne la chute des CD4 et l'affaiblissement de l'immunité.
\item Chaque enzyme virale (transcriptase inverse, intégrase, protéase) est une cible des traitements antirétroviraux.
\end{itemize}
\end{aretenir}

Maintenant que nous savons \emph{comment} le virus se multiplie dans les lymphocytes, il reste à comprendre \emph{comment} il passe d'une personne à l'autre — et surtout comment l'en empêcher : c'est l'objet de la section suivante, consacrée à la transmission et à la prévention.

\coursec{Transmission et prévention}

Après avoir compris comment le VIH détourne la machinerie du lymphocyte CD4 pour se multiplier (\emph{section précédente}), une question s'impose naturellement : comment ce virus, fragile hors de son hôte, parvient-il à passer d'une personne à une autre ? L'épidémiologie a très tôt identifié que le VIH ne se transmet pas par simple contact social (serrer la main, partager un repas), mais uniquement lors de mises en contact de \textbf{liquides biologiques spécifiques} contenant une charge virale suffisante. Cette section analyse ces portes d'entrée, démêle le vrai du faux sur les situations à risque, et détaille les stratégies de prévention validées scientifiquement et appliquées au Cameroun.

\courssub{Les trois modes de transmission du VIH}

\paragraph{Observation — Situation problème.}
Dans un service de consultations externes d'un hôpital de district au Cameroun, trois patients reçoivent un diagnostic de séropositivité récente :
\begin{itemize}
    \item \textbf{Patiente A}, 24 ans, en couple stable, n'a jamais reçu de transfusion ni utilisé de drogues injectables.
    \item \textbf{Patient B}, 32 ans, soudeur, a partagé du matériel d'injection pour une prise en charge médicale traditionnelle il y a six mois.
    \item \textbf{Nourrisson C}, 6 mois, né d'une mère séropositive non suivie pendant la grossesse, allaité au sein exclusivement.
\end{itemize}
\emph{Question :} Quel liquide biologique a servi de véhicule au virus dans chaque cas ? Quelles mesures auraient pu empêcher ces contaminations ?

\paragraph{Analyse des liquides vecteurs.}
Le VIH est présent en quantité infectieuse (\emph{charge virale} élevée) dans seulement quelques liquides biologiques : le \textbf{sang} (et ses dérivés), le \textbf{sperme} et le \textbf{liquide pré-éjaculatoire}, les \textbf{sécrétions vaginales} et le \textbf{lait maternel}. À l'inverse, la salive, les larmes, la sueur, l'urine, les selles et le vomi ne transmettent \emph{pas} le VIH (sauf contamination visible par du sang), car ils contiennent soit des facteurs antiviraux (salive), soit une concentration virale trop faible, soit aucune cellule cible (lymphocytes CD4).

\begin{propriete}[Liquides biologiques transmissibles]
Le VIH se transmet \textbf{uniquement} par trois liquides biologiques principaux lorsqu'ils pénètrent dans l'organisme d'une personne saine par une porte d'entrée (muqueuse lésée ou intacte, peau effractée, injection directe) :
\begin{enumerate}
    \item Le \textbf{sang} (globules rouges, plasma, plaquettes) et les liquides visiblement sanglants.
    \item Les \textbf{sécrétions sexuelles} : sperme, liquide pré-éjaculatoire, sécrétions vaginales et rectales.
    \item Le \textbf{lait maternel}.
\end{enumerate}
\emph{Toute autre voie (salive, sueur, larmes, piqûre d'insecte, contact cutané intact, air, eau, nourriture) est sans risque de transmission du VIH.}
\end{propriete}

\paragraph{1. Transmission sexuelle (voie majeure au Cameroun).}
C'est le mode de contamination le plus fréquent (plus de 90 \% des cas en Afrique subsaharienne). Elle survient lors de \textbf{rapports sexuels non protégés} (vaginaux, anaux, oraux) avec un partenaire séropositif. Le virus franchit les muqueuses génitales, rectales ou buccales, souvent facilité par la présence d'autres \textbf{infections sexuellement transmissibles (IST)} (ulcérations syphilitiques, herpès, chlamydiae) qui augmentent la perméabilité muqueuse et recrutent des lymphocytes CD4 cibles au portail d'entrée.
\begin{itemize}
    \item \textbf{Rapports vaginaux :} risque élevé pour la femme (surface muqueuse large, contact prolongé avec le sperme) et pour l'homme (contact avec sécrétions vaginales).
    \item \textbf{Rapports anaux :} risque très élevé (muqueuse rectale fine, vascularisée, sujette aux micro-lésions) pour le partenaire réceptif \emph{et} insertif.
    \item \textbf{Rapports oraux :} risque faible mais réel (éjaculation dans la bouche, gingivites, plaies buccales).
\end{itemize}

\paragraph{2. Transmission sanguine (voie parenterale).}
Elle résulte de l'introduction directe de sang contaminé dans la circulation sanguine ou les tissus profonds.
\begin{itemize}
    \item \textbf{Partage de seringues/aiguilles} chez les usagers de drogues injectables (UDI) ou dans le cadre de soins traditionnels non stériles (scarifications, piqûres traditionnelles).
    \item \textbf{Transfusion de sang non testé} ou mal dépisté (risque résiduel en période de fenêtre sérologique). Au Cameroun, le sang est systématiquement dépisté (VIH 1/2, VHB, VHC, Syphilis) avant transfusion.
    \item \textbf{Accidents d'exposition au sang (AES)} : piqûre accidentelle d'aiguille chez le personnel soignant, coupure par instrument souillé.
    \item \textbf{Objets tranchants/perforants souillés} partagés : rasoirs, lames de scarification, matériel de tatouage/piercing non stérilisé, brosse à dents (si saignement gingival).
\end{itemize}

\paragraph{3. Transmission mère-enfant (TME) ou transmission verticale.}
Sans intervention, une femme séropositive transmet le virus à son enfant dans 15 à 45 \% des cas. Elle peut survenir à trois moments :
\begin{enumerate}
    \item \textbf{In utero (grossesse) :} passage transplacentaire (rare, \textasciitilde 5-10 \%), surtout en fin de grossesse si charge virale maternelle élevée ou rupture prématurée des membranes.
    \item \textbf{Péri-partum (accouchement) :} moment le plus critique (\textasciitilde 10-20 \%). Contact direct du sang maternel et des sécrétions génitales avec la peau/muqueuses du nouveau-né (surtout lors du passage dans les voies basses).
    \item \textbf{Post-partum (allaitement) :} transmission par le lait maternel (\textasciitilde 5-20 \% sur 18-24 mois d'allaitement). Le virus est présent sous forme libre et cellulaire (lymphocytes infectés) dans le lait.
\end{enumerate}

\begin{popfigure}
\legende{Tableau synoptique des trois modes de transmission du VIH, liquides vecteurs, situations à risque et prévention principale.}
\begin{center}
\begin{tabularx}{\linewidth}{|l|X|X|X|}
\hline
\rowcolor{popblueL}
\textbf{Mode de transmission} & \textbf{Liquides biologiques vecteurs} & \textbf{Situations à risque typiques} & \textbf{Prévention principale} \\
\hline
\textbf{Sexuelle} & Sperme, liquide pré-éjaculatoire, sécrétions vaginales/rectales & Rapports vaginaux/anaux/oraux non protégés ; partenaires multiples ; IST non traitées & \textbf{Préservatif} (masculin/féminin) à chaque rapport ; Fidélité mutuelle après dépistage ; Dépistage couple ; PrEP (prophylaxie pré-exposition) pour populations clés \\
\hline
\textbf{Sanguine} & Sang (et liquides visiblement sanglants) & Partage seringues/aiguilles (UDI, soins traditionnels) ; Transfusion non testée ; AES personnel soignant ; Partage rasoirs, matériel scarification/tatouage non stérile & \textbf{Matériel stérile à usage unique} (seringues, aiguilles, lames) ; Dépistage systématique donneurs de sang ; Précautions universelles (gants, conteneurs OPCT) ; PPE (prophylaxie post-exposition) après AES \\
\hline
\textbf{Mère-enfant (TME)} & Sang maternel (accouchement), Lait maternel (allaitement) & Grossesse/accouchement/allaitement sans prise en charge PTME ; Charge virale maternelle non contrôlée ; Allaitement mixte (sein + eau/bouillie) & \textbf{PTME complète} : ARV mère (vie entière ou option B+), Accouchement assisté (césarienne si CV > 1000), ARV nourrisson (névirapine/zidovudine 6-12 sem), \textbf{Allaitement exclusif 6 mois} puis sevrage rapide OU lait de substitution si AFASS (Acceptable, Faisable, Abordable, Durable, Sans danger) \\
\hline
\end{tabularx}
\end{center}
\end{popfigure}

\courssub{Situations sans risque et idées reçues : la méthode du « liquide vecteur »}

Face à une situation de la vie courante, la peur ou la méconnaissance conduisent souvent à surestimer ou sous-estimer le risque. La démarche scientifique rigoureuse consiste à identifier le \textbf{liquide biologique} mis en jeu.

\begin{methode}[Évaluer le risque VIH face à une situation de la vie quotidienne]
\textbf{Étape 1 :} Identifier le \textbf{liquide biologique} en contact (sang, sperme, sécrétions vaginales, lait maternel, salive, sueur, larmes, urine, selles).
\textbf{Étape 2 :} Vérifier si ce liquide figure sur la \textbf{liste des liquides transmissibles} (Propriété ci-dessus).
\textbf{Étape 3 :} S'il s'agit d'un liquide transmissible, vérifier l'existence d'une \textbf{porte d'entrée} (muqueuse, peau lésée, injection).
\textbf{Étape 4 :} Conclure :
\begin{itemize}
    \item Liquide non transmissible $\rightarrow$ \textbf{Risque nul} (aucune mesure d'urgence).
    \item Liquide transmissible + porte d'entrée $\rightarrow$ \textbf{Risque réel} $\rightarrow$ Dépistage immédiat + PPE si $< 48-72$ h.
\end{itemize}
\end{methode}

\begin{exemplebox}[Application de la méthode : situations courantes]
\begin{itemize}
    \item \textbf{Embrassade (baiser profond) :} Liquide = salive (non transmissible). Même avec micro-saignements gingivaux, le volume de sang est infime et la salive inhibe le VIH. \textbf{Risque nul}.
    \item \textbf{Partage de couverts/verres/toilettes :} Liquides = salive, urine, sueur (non transmissibles). \textbf{Risque nul}.
    \item \textbf{Piqûre de moustique :} Voir \textbf{Attention} ci-dessous. \textbf{Risque nul}.
    \item \textbf{Morsure humaine :} Liquide = salive (+ sang si peau percée). Si la peau est percée et saignement des deux côtés : \textbf{risque théorique très faible mais documenté} (quelques cas rares). $\rightarrow$ PPE recommandée.
    \item \textbf{Soins à un blessé (premiers secours) :} Liquide = sang. Porte d'entrée = mains du secouriste (coupures, éraflures). \textbf{Risque réel} $\rightarrow$ Port de gants (précautions universelles) + lavage immédiat si contact sanguin sur peau lésée.
\end{itemize}
\end{exemplebox}

\begin{attention}[Erreur fréquente : « Le moustique transmet le VIH »]
\textbf{FAUX.} C'est l'idée reçue la plus tenace. Trois arguments scientifiques l'infirment :
\begin{enumerate}
    \item \textbf{Pas de multiplication :} Le VIH ne se réplique \emph{pas} dans l'intestin ni les glandes salivaires du moustique (contrairement au paludisme, dengue, chikungunya). Il est digéré.
    \item \textbf{Quantité infime :} La proboscis du moustique retient \textasciitilde 0,00004 mL de sang (40 nL), soit 1000 à 10 000 fois moins qu'une seringue. La dose infectieuse n'est pas atteinte.
    \item \textbf{Mécanique :} Le moustique n'injecte pas le sang de la victime précédente ; il injecte sa salive (anticoagulant) et aspire le sang dans un canal distinct.
\end{enumerate}
\textbf{Piège au Probatoire :} Ne jamais écrire « le moustique ne transmet pas le VIH parce qu'il ne pique qu'une fois ». Il pique plusieurs fois, mais la biologie du virus dans l'insecte empêche la transmission.
\end{attention}

\begin{savaistu}[Contexte camerounais : Accès à la prévention]
Au Cameroun, la riposte nationale (CNLS, Ministère de la Santé Publique) garantit :
\begin{itemize}
    \item \textbf{Dépistage VIH gratuit} dans la quasi-totalité des formations sanitaires publiques (hôpitaux, centres de santé, CEDC - Centres de Dépistage Volontaire) et lors de campagnes mobiles (journées mondiales, caravanes).
    \item \textbf{Préservatifs masculins et féminins gratuits} distribués dans les centres de santé, les ONG partenaires (CAMNAFAW, ACMS), les établissements scolaires/universitaires lors de campagnes « Vacances sans Sida ».
    \item \textbf{Prise en charge PTME gratuite} (ARV mère/enfant, biologiques de suivi) dans le cadre de l'élimination de la transmission mère-enfant (eTME).
    \item \textbf{PPE gratuite} pour les AES (personnel de santé) et les victimes de violences sexuelles (PEV - Prophylaxie Post-Exposition Viol).
\end{itemize}
Connaître ces ressources, c'est déjà se protéger et protéger les autres.
\end{savaistu}

\courssub{Les moyens de prévention : la stratégie combinée}

Aucun moyen unique n'offre une protection à 100 \% en situation réelle (rupture de préservatif, oubli de prise ARV, fenêtre sérologique). L'OMS et le Cameroun promeuvent la \textbf{prévention combinée} : association de mesures biomédicales, comportementales et structurelles adaptées au contexte de chaque personne.

\paragraph{1. Prévention de la transmission sexuelle : La règle « ABC » et au-delà.}
\begin{itemize}
    \item \textbf{A — Abstinence} (ou retard de la première relation) : Efficacité 100 \% si respectée. Recommandée pour les adolescents non prêts.
    \item \textbf{B — Be faithful (Fidélité mutuelle)} : Efficacité maximale \textbf{si les deux partenaires sont séronégatifs et le restent}. Exige un \textbf{dépistage de couple} avant d'arrêter le préservatif.
    \item \textbf{C — Condom (Préservatif)} : Seul moyen protégeant \textbf{simultanément} contre VIH, IST et grossesses non désirées.
    \begin{itemize}
        \item \textbf{Masculin (latex/polyuréthane) :} Efficacité \textasciitilde 98 \% en usage parfait, \textasciitilde 85 \% en usage réel. \emph{Règles d'or :} Vérifier date/percée/emballage ; ouvrir avec les doigts (pas dents/ciseaux) ; pincer le réservoir ; dérouler sur pénis en érection ; retirer avant érection complète en tenant la base ; nouer, jeter.
        \item \textbf{Féminin (nitrile) :} Inséré dans le vagin avant le rapport. Protège aussi la vulve. Utilisable par la femme sans négociation directe au moment de l'acte.
    \end{itemize}
    \item \textbf{Dépistage avant mariage / nouvelle union :} Offert gratuitement, permet au couple de connaître son statut et de choisir sa prévention en connaissance de cause (concordance sérologique vs discordance).
    \item \textbf{PrEP (Prophylaxie Pré-Exposition) :} Prise quotidienne de Ténofovir/Emtricitabine (Truvada\textsuperscript{\textregistered} ou génériques) par une personne séronégative à risque élevé (travailleuses du sexe, couples sérodiscordants, HSH, UDI). Efficacité $> 90 \%$ si observance stricte. Disponible au Cameroun dans certains centres pilotes.
    \item \textbf{Traitement comme prévention (TasP / I=I) :} Une personne séropositive sous ARV efficace (charge virale indétectable $< 50$ copies/mL depuis $> 6$ mois) \textbf{ne transmet pas le VIH par voie sexuelle} (Indétectable = Intransmissible). Argument majeur pour l'observance et la lutte contre la stigmatisation.
\end{itemize}

\paragraph{2. Prévention de la transmission sanguine.}
\begin{itemize}
    \item \textbf{Usage unique et stérilité absolue} de tout matériel pénétrant la peau (seringues, aiguilles, cathéters, lames de bistouri, matériel de scarification/tatouage/piercing).
    \item \textbf{Précautions universelles (standard precautions)} pour tout personnel soignant : port de gants pour tout contact potentiel avec du sang/liquides biologiques, gestion sécurisée des déchets piquants-tranchants (conteneurs OPCT jaunes), vaccination anti-HBV obligatoire.
    \item \textbf{Dépistage systématique des donneurs de sang} (tests sérologiques + PCR/NAT pour réduire la fenêtre).
    \item \textbf{Réduction des risques (RdR) pour UDI :} Programmes d'échange de seringues (PEP), substitution aux opioïdes (méthadone, buprénorphine) — en développement au Cameroun.
\end{itemize}

\paragraph{3. Arbre de décision : « Situation à risque ou non ? »}
Cet algorithme guide la décision immédiate face à un contact potentiel.

\begin{popfigure}
\legende{Arbre de décision pour qualifier une situation d'exposition au VIH et orienter la conduite à tenir (PPE = Prophylaxie Post-Exposition).}
\begin{center}
\begin{tikzpicture}[
    node distance=1.2cm and 1.5cm,
    startstop/.style={rectangle, rounded corners, fill=popblue!20, draw=popblue, thick, text width=3.5cm, align=center, font=\small},
    decision/.style={diamond, aspect=2, fill=poporange!20, draw=poporange, thick, text width=3cm, align=center, font=\small, inner sep=2pt},
    process/.style={rectangle, fill=popgreen!15, draw=popgreen, thick, text width=3.5cm, align=center, font=\small},
    io/.style={trapezium, trapezium left angle=70, trapezium right angle=110, fill=poppurple!15, draw=poppurple, thick, text width=3.5cm, align=center, font=\small},
    arrow/.style={-Stealth, thick, draw=popdark},
    yes/.style={arrow},
    no/.style={arrow}
]
% Nodes
\node (start) [startstop] {Contact avec un liquide biologique potentiellement contaminé};
\node (liq) [decision, below of=start] {Le liquide est-il du \textbf{sang}, \textbf{sperme}, \textbf{sécrétions vaginales} ou \textbf{lait maternel} ?};
\node (no_risk) [process, left of=liq, xshift=-3.5cm, align=center]{\textbf{RISQUE NUL}\\(Salive, sueur, larmes, urine, selles, vomi sans sang visible, piqûre moustique, embrassade, objets partagés)\\\textbf{Action :} Rassurer, éduquer, pas de PPE};
\node (entry) [decision, right of=liq, xshift=3.5cm, align=center]{Y a-t-il une \textbf{porte d'entrée} ?\\(Muqueuse, peau lésée, injection, accident aiguille, rapport non protégé)};
\node (low_risk) [process, below of=entry, align=center]{Risque \textbf{THÉORIQUE FAIBLE}\\(ex: sang sur peau intacte, rapport oral sans éjaculation)\\\textbf{Action :} Lavage eau/savon (peau) ou rinçage (muqueuse) ; Dépistage conseil à 6 sem ; PPE \emph{généralement non indiquée}};
\node (high_risk) [io, below of=low_risk, align=center]{\textbf{RISQUE RÉEL AVÉRÉ}\\(Rapport sexuel non protégé, partage seringue, AES profond, transfusion non testée, accouchement non protégé, morsure saignante)\\\textbf{URGENCE :} PPE \textbf{dans les 4 h (idéal) à 48-72 h max} \\ Dépistage immédiat (TROD) + Biologie de base (NFS, Rénale, Hépatites) \\ Suivi sérologique à J0, 6 sem, 3 mois};
\node (ppe) [process, below of=high_risk, yshift=-0.5cm, align=center]{PPE = Traitement ARV d'urgence 28 jours (ex: TDF/3TC + LPV/r ou DTG)\\ Observance stricte = efficacité $\approx$ 80-90\% \\ Gratuite dans les structures publiques (AES, Viol)};

% Arrows
\draw [arrow] (start) -- (liq);
\draw [yes] (liq) -- node[above, xshift=-1.5cm] {\textbf{OUI}} (entry);
\draw [no] (liq) -- node[above] {\textbf{NON}} (no_risk);
\draw [yes] (entry) -- node[above] {\textbf{OUI}} (high_risk);
\draw [no] (entry) -- node[above] {\textbf{NON}} (low_risk);
\draw [arrow] (high_risk) -- (ppe);
\end{tikzpicture}
\end{center}
\end{popfigure}

\courssub{Prévention de la transmission mère-enfant (PTME) : l'objectif « Zéro contamination »}

Au Cameroun, la PTME suit l'option \textbf{B+} (OMS) : \textbf{traitement ARV à vie pour toute femme enceinte/allaitante séropositive}, indépendamment du taux de CD4. C'est la stratégie la plus efficace pour la mère et l'enfant.

\begin{experience}[Démonstration de l'efficacité de la PTME : données de cohorte]
\textbf{Protocole (observationnel) :} Suivi de 1000 couples mère-enfant au Cameroun.
\begin{itemize}
    \item \textbf{Groupe 1 (Sans PTME) :} Mères non traitées, accouchement vaginal, allaitement mixte 18 mois.
    \item \textbf{Groupe 2 (PTME complète Option B+) :} Mères sous ARV dès le diagnostic (idéalement avant grossesse ou 1er trimestre), charge virale indétectable à l'accouchement, accouchement assisté, ARV prophylactique au nourrisson (NVP 6 sem), allaitement exclusif 6 mois puis sevrage rapide.
\end{itemize}
\textbf{Résultats attendus (littérature / données PNLS Cameroun) :}
\begin{center}
\begin{tabular}{|l|c|c|}\hline
\textbf{Indicateur} & \textbf{Groupe 1 (Sans PTME)} & \textbf{Groupe 2 (PTME B+ complète)} \\ \hline
Transmission in utero / intra-partum & 15 -- 20 \% & $< 1$ \% \\ \hline
Transmission post-partum (allaitement 18 m) & 10 -- 15 \% & $< 1$ \% \\ \hline
\textbf{Transmission globale finale} & \textbf{25 -- 35 \%} & \textbf{$< 2$ \% (souvent $< 1$ \%)} \\ \hline
\end{tabular}
\end{center}
\textbf{Interprétation :} La charge virale maternelle indétectable est le déterminant principal. L'allaitement exclusif (sans eau/bouillie) protège la muqueuse intestinale du nourrisson et réduit le risque par rapport à l'allaitement mixte.
\textbf{Conclusion :} La PTME complète transforme une épidémie pédiatrique en événement rare. Le défi camerounais : \textbf{le dépistage précoce (1er trimestre), la rétention en soins et l'observance pendant l'allaitement}.
\end{experience}

\begin{aretenir}[Chiffres clés PTME — Probatoire]
\begin{itemize}
    \item Sans intervention : risque transmission mère-enfant = \textbf{15 à 45 \%}.
    \item Avec PTME complète (ARV mère vie entière + ARV nourrisson + allaitement exclusif 6 mois + sevrage rapide) : risque $< \textbf{2 \%}$ (objectif OMS élimination $< 5 \%$).
    \item \textbf{Allaitement exclusif 6 mois} + sevrage rapide sur 1 mois est recommandé au Cameroun si conditions AFASS non réunies pour le lait de substitution. Le lait de substitution (lait infantile) élimine le risque post-partum mais expose à la diarrhée/malnutrition si eau non potable/hygiène précaire.
\end{itemize}
\end{aretenir}

\courssub{Complément Probatoire : La Prophylaxie Post-Exposition (PPE)}

C'est une \textbf{urgence médicale} : l'administration d'un traitement antirétroviral de 28 jours à une personne séronégative après une exposition à risque majeur, pour empêcher l'installation de l'infection.

\begin{propriete}[Indications, délai et schéma de la PPE]
\begin{itemize}
    \item \textbf{Indications :} Rapport sexuel non protégé (viol, rupture préservatif, partenaire séropositif CV non contrôlée) ; Accident d'exposition au sang (AES : piqûre aiguille profonde, coupure instrument souillé, projection sang sur muqueuse/peau lésée) ; Partage seringue UDI.
    \item \textbf{Délai critique :} \textbf{Idéalement $< 4$ h, maximum 48--72 h} après l'exposition. Au-delà de 72 h, l'efficacité n'est pas prouvée (non indiquée).
    \item \textbf{Bilan initial (J0) :} TROD VIH (si positif $\rightarrow$ pas de PPE, prise en charge ARV vie entière) ; TROD VHB/VHC ; NFS, Créatininémie, Transaminases ; Test de grossesse (femme en âge de procréer).
    \item \textbf{Schéma recommandé (Cameroun/OMS) :} \textbf{3 ARV pendant 28 jours}.
    \begin{itemize}
        \item Adultes/Ados $> 35$ kg : \textbf{TDF 300 mg + 3TC 300 mg (ou FTC 200 mg) 1 cp/jour} \textbf{+ DTG 50 mg 1 cp/jour} (préféré : tolérance, barrière génétique haute, compatible grossesse). Alternative : LPV/r 2 cp 2$\times$/j.
        \item Enfants : Selon poids (sirop/zidovudine + lamivudine + lopinavir/ritonavir ou dolutégravir dispersable).
    \end{itemize}
    \item \textbf{Suivi :} J7 (tolérance), J28 (fin PPE), Sérologie de contrôle à \textbf{6 semaines} et \textbf{3 mois} (fin fenêtre sérologique).
    \item \textbf{Gratuité :} Dans les structures publiques pour AES (personnel) et PEV (violences sexuelles).
\end{itemize}
\end{propriete}

\begin{exoresolu}[Cas clinique : Gestion d'un accident d'exposition au sang (AES)]
\textbf{Énoncé :} Une infirmière de 28 ans, en service de chirurgie, se pique profondément au doigt avec une aiguille creuse venant d'être utilisée sur un patient VIH+ (charge virale 45 000 copies/mL, sous ARV depuis 2 ans mais observance incertaine). L'accident survient à 10 h 30. L'infirmière est séronégative connue (test il y a 3 mois), non enceinte. Que faire ? \textbf{Conduite à tenir étape par étape.}

\textbf{Solution détaillée :}
\begin{enumerate}
    \item \textbf{Geste immédiat (10 h 30) :} Ne pas faire saigner, ne pas sucer. \textbf{Lavage abondant à l'eau et au savon} (ou solution antiseptique type Bétadine\textsuperscript{\textregistered} dermique) pendant 5 min. Rincer. Ne pas injecter d'antiseptique dans la plaie.
    \item \textbf{Déclaration (10 h 35) :} Signaler au médecin chef de service / référent AES. Remplir la fiche de déclaration d'AES (anonyme pour le patient source).
    \item \textbf{Évaluation du risque (10 h 40) :} Aiguille creuse + sang visible + piqûre profonde + patient source VIH+ (CV détectable) = \textbf{Risque ÉLEVÉ (0,3 \% sans PPE)}. Indication formelle de PPE.
    \item \textbf{Bilan J0 (11 h 00) :} TROD infirmière (négatif) ; TROD patient source (confirmation VIH+, idéalement VHB/VHC) ; NFS, Ionogramme, Créatinine, ASAT/ALAT ; Test grossesse (négatif).
    \item \textbf{Initiation PPE (11 h 30 -- bien dans les 4 h) :} Prescription : \textbf{TDF/3TC (1 cp/j) + DTG 50 mg (1 cp/j) pour 28 jours}. Explication : effets secondaires digestifs possibles (nausées, diarrhée) transitoires, prise au repas, ne pas arrêter sans avis. Remise des 28 jours de traitement.
    \item \textbf{Suivi :} J7 : contrôle tolérance clinique/biologique (créatinine). J28 : fin traitement. \textbf{Sérologie contrôle à 6 semaines (mi-juin) et 3 mois (fin juillet)}. Si toutes négatives $\rightarrow$ exclusion contamination.
    \item \textbf{Prévention secondaire :} Analyse des causes (surcharge, matériel inadapté, formation), vaccination anti-HBV si non immunisée.
\end{enumerate}
\textbf{Point clé :} La rapidité (délai $< 4$ h) et l'observance des 28 jours sont les garants de l'efficacité.
\end{exoresolu}

\begin{exercice}[Entraînement : Évaluer le risque et la conduite à tenir]
Pour chacune des situations ci-dessous, précisez : (1) Le liquide en cause, (2) Le niveau de risque (Nul / Faible / Élevé), (3) La conduite immédiate (Rassurer / Lavage + Dépistage différé / PPE Urgente + Bilan).
\begin{enumerate}
    \item Un adolescent partage une canette de soda avec un camarade séropositif (non déclaré) qui a une gingivite saignante.
    \item Une femme enceinte de 8 mois, séropositive connue, sous ARV depuis 5 ans (CV $< 20$ copies/mL), accouche par voie basse. Le nouveau-né reçoit la prophylaxie NVP à la naissance.
    \item Un technicien de laboratoire manipule un tube de sang VIH+ (CV 120 000 copies/mL) sans gants ; il a une micro-coupure fraîche à l'index. Il se lave les mains immédiatement à l'eau de Javel diluée.
    \item Un couple sérodiscordant (Homme VIH+ sous ARV CV indétectable depuis 2 ans, Femme VIH-) a un rapport vaginal non protégé désirant un enfant.
\end{enumerate}
\end{exercice}

\begin{corrige}[Exercice n°1]
\begin{enumerate}
    \item \textbf{Liquide :} Salive (+ traces sang). \textbf{Risque : NUL} (salive non transmissible, volume sanguin infime, inhibition virale). \textbf{Conduite :} Rassurer, éduquer sur les vrais modes de transmission.
    \item \textbf{Liquide :} Sang maternel / sécrétions génitales (accouchement) + Lait maternel (post-partum). \textbf{Risque : TRÈS FAIBLE ($< 1 \%$)} grâce à CV indétectable maternelle (TasP) + PTME (ARV mère + NVP nourrisson + allaitement exclusif). \textbf{Conduite :} Suivi PTME standard (PCR VIH nourrisson à 6 sem, 9 mois, fin allaitement). Pas de PPE pour le nourrisson (prophylaxie NVP 6 sem suffit).
    \item \textbf{Liquide :} Sang. \textbf{Risque : ÉLEVÉ} (sang + peau lésée + CV source élevée). L'eau de Javel sur peau lésée est irritante et non validée comme PPE locale. \textbf{Conduite :} \textbf{PPE URGENTE} (idéalement $< 4$ h) + Bilan J0 + Suivi sérologique 6 sem / 3 mois. Déclaration AES.
    \item \textbf{Liquide :} Sperme. \textbf{Risque : NUL (quasi nul scientifiquement)}. Preuve : Études PARTNER / Opposites Attract : 0 transmission linked dans couples sérodiscordants avec CV indétectable durable ($> 6$ mois). \textbf{Conduite :} Rassurer sur concept I=I (Indétectable = Intransmissible). Dépistage femme routine. Conception naturelle possible sans PPE. Conseil préconceptionnel (acide folique, dépistage IST).
\end{enumerate}
\end{corrige}

\paragraph{Synthèse : Les 5 piliers de la prévention VIH au Cameroun.}
\begin{enumerate}
    \item \textbf{Connaître son statut :} Dépistage volontaire, gratuit, confidentiel (TROD 20 min), recommandé au moins 1 fois/an (population générale) ou 3 mois (populations clés).
    \item \textbf{Se protéger / Protéger l'autre :} Préservatif (gratuit), Fidélité mutuelle testée, PrEP (si éligible), TasP (si VIH+).
    \item \textbf{Éviter l'exposition au sang :} Matériel stérile à usage unique, précautions universelles, RdR pour UDI.
    \item \textbf{Éliminer la transmission mère-enfant :} Dépistage prénatal précoce (1er trimestre), ARV mère vie entière (Option B+), Accouchement sécurisé, Allaitement exclusif 6 mois + sevrage rapide ou substitution si AFASS.
    \item \textbf{Agir vite en cas d'accident :} PPE dans les 4 h (max 72 h) après AES ou viol, gratuite en secteur public.
\end{enumerate}

La lutte contre le VIH/Sida n'est pas seulement médicale, elle est \textbf{sociale} : combattre la stigmatisation, promouvoir l'égalité de genre, scolariser les filles, réduire la pauvreté. Chaque élève de Première, futur citoyen, professionnel de santé, parent ou décideur, porte sa part de cette responsabilité collective. La section suivante détaillera l'évolution naturelle de l'infection sans traitement, pour comprendre pourquoi le dépistage précoce et l'ARV à vie sont vitaux.

\coursec{Évolution de l'infection à VIH}

Après avoir compris \emph{comment} le VIH se transmet et \emph{comment} s'en protéger (section précédente), nous devons maintenant suivre ce qui se passe \emph{à l'intérieur de l'organisme} après la contamination. Pourquoi une personne peut-elle vivre des années sans symptômes tout en étant contagieuse ? Pourquoi le système immunitaire finit-il par s'effondrer ? C'est l'histoire d'une course de fond entre le virus et nos défenses, une course que l'on peut suivre grâce à deux indicateurs biologiques majeurs : la \textbf{charge virale} et le \textbf{taux de lymphocytes CD4}.

\courssub{Phase 1 : L'infection primaire et la séroconversion}

\paragraph{Observation : les premiers jours après la contamination.}
Imagine un jeune adulte contaminé lors d'un rapport sexuel non protégé. Une à quatre semaines plus tard, il consulte pour une fièvre modérée, des ganglions au cou, une éruption cutanée et une grande fatigue. Le médecin, devant ce \emph{syndrome pseudo-grippal} ou \emph{mononucléosique}, ne pense pas immédiatement au VIH car ces signes sont banals et disparaissent spontanément en une à deux semaines. Pourtant, c'est la première bataille perdue par le système immunitaire.

\paragraph{Investigation : que se passe-t-il au niveau cellulaire ?}
\begin{itemize}
    \item \textbf{Hypothèse :} Le virus se réplique activement avant que la réponse immunitaire spécifique ne se mette en place.
    \item \textbf{Données expérimentales (modèle animal et mesures chez l'homme) :} Dès les premiers jours, le VIH infecte massivement les lymphocytes CD4 des tissus lymphoïdes (ganglions, rate, muqueuses intestinales). La \cle{charge virale} (nombre de copies d'ARN viral par mL de sang) atteint un \textbf{pic extrêmement élevé} (souvent $> 10^6$ copies/mL). Simultanément, le taux de \textbf{lymphocytes CD4} chute brutalement (parfois $< 200$ cellules/mm³ transitoirement).
    \item \textbf{Interprétation :} L'organisme ne dispose pas encore d'anticorps ni de lymphocytes T cytotoxiques (CD8) spécifiques efficaces. Le virus a « champ libre ».
    \item \textbf{Réponse immunitaire :} Vers la 2\textsuperscript{e}--4\textsuperscript{e} semaine, le système immunitaire réagit : apparition d'anticorps anti-VIH (séroconversion) et expansion de clones CD8 cytotoxiques. La charge virale chute alors drastiquement (de 2 à 3 log) pour se stabiliser à un \emph{point de set} (niveau d'équilibre), et les CD4 remontent partiellement, sans retrouver leur niveau initial.
\end{itemize}

\begin{definition}[Séroconversion]
La \textbf{séroconversion} est l'apparition dans le sang d'anticorps spécifiques dirigés contre le VIH (anti-gp120, anti-gp41, anti-p24), détectables par les tests de dépistage (ELISA, tests rapides). Elle survient généralement \textbf{3 à 8 semaines} après la contamination (parfois jusqu'à 3 mois). C'est ce marqueur que recherchent les tests de dépistage : avant elle, le test est négatif (faux négatif possible) ; après elle, le test devient positif.
\end{definition}

\begin{attention}[La fenêtre sérologique : un piège majeur]
Pendant la période comprise entre la contamination et la séroconversion, appelée \textbf{fenêtre sérologique}, la personne est \textbf{infectée, très contagieuse (charge virale maximale)}, mais le test de dépistage standard (recherche d'anticorps) est \textbf{négatif}. C'est pourquoi, face à une prise de risque récente, on ne se rassure pas sur un test négatif immédiat : il faut le renouveler après 6 semaines (tests 4\textsuperscript{e} génération Ag/Ac) ou 3 mois (tests rapides anticorps seuls). Au Cameroun, les centres de dépistage (CDV, CPEC) insistent sur ce délai de confirmation.
\end{attention}

\begin{popfigure}
\begin{tikzpicture}[
    node distance=1.2cm and 1.5cm,
    box/.style={draw=popblue, thick, rounded corners, fill=popblue!10, minimum width=2.5cm, minimum height=1cm, align=center},
    arrow/.style={-Stealth, thick, popdark},
    wave/.style={decorate, decoration={snake, amplitude=1pt, segment length=4pt}, thick, poporange}
]
% Timeline
\draw[thick, popdark, ->] (0,0) -- (14,0) node[right] {Temps (semaines)};
\foreach \x/\label in {0/0, 2/2, 4/4, 6/6, 8/8, 10/10, 12/12} {
    \draw (\x, -0.1) -- (\x, 0.1) node[below] {\label};
}

% Viral Load Curve (schematic)
\draw[poporange, very thick, smooth] (0,3) .. controls (1,3.5) and (2.5,4) .. (3,4) node[above, pos=0.9, poporange] {\small Charge virale (ARN VIH)};
\draw[poporange, very thick, smooth] (3,4) .. controls (4,3) and (5,1.5) .. (6,1.5);
\draw[poporange, very thick, smooth] (6,1.5) .. controls (8,1.2) and (10,1.3) .. (12,1.3);

% CD4 Curve (schematic)
\draw[popgreen, very thick, smooth] (0,1) .. controls (1,0.8) and (2.5,0.2) .. (3,0.2) node[below, pos=0.9, popgreen] {\small Taux CD4};
\draw[popgreen, very thick, smooth] (3,0.2) .. controls (4,0.6) and (5,0.9) .. (6,0.9);
\draw[popgreen, very thick, smooth] (6,0.9) .. controls (8,0.85) and (10,0.8) .. (12,0.75);

% Seroconversion marker
\draw[poppurple, dashed, thick] (4.5, -0.5) -- (4.5, 4.2);
\node[poppurple, anchor=south] at (4.5, 4.3) {\textbf{Séroconversion}};
\node[poppurple, anchor=north, align=center] at (4.5, -0.7) {Apparition\\ des Ac anti-VIH};

% Window period bracket
\draw[poppink, thick, decorate, decoration={brace, amplitude=5pt}] (0, -1.2) -- (4.5, -1.2) node[midway, below=5pt, poppink, align=center] {\textbf{Fenêtre sérologique}\\(Test Ac négatif)};

% Phase labels
\node[align=center, popdark] at (1.5, 4.5) {Infection\\ primaire};
\node[align=center, popdark] at (9, 4.5) {Phase\\ asymptomatique};

% Legend boxes
\node[box, anchor=north west, align=center] at (0.5, 3.8){Pic de réplication\\ virale intense};
\node[box, anchor=south west, fill=popgreen!10, draw=popgreen, align=center] at (0.5, 0.5){Chute transitoire\\ des CD4};
\end{tikzpicture}
\legende{Évolution schématique de la charge virale (orange) et du taux de lymphocytes CD4 (vert) lors de l'infection primaire. La séroconversion (ligne violette) marque l'apparition des anticorps et la fin de la fenêtre sérologique (zone rose).}
\end{popfigure}

\begin{savaistu}[Le VIH et l'intestin : le réservoir caché]
Dès l'infection primaire, le VIH détruit massivement les lymphocytes CD4 de la muqueuse intestinale (les \emph{Th17}, gardiens de la barrière intestinale). Cette perte précoce et irréversible favorise la \textbf{translocation microbienne} (passage de bactéries dans le sang), source d'une inflammation chronique qui épuise le système immunitaire bien avant le stade Sida. C'est l'une des raisons pour lesquelles le traitement précoce est crucial.
\end{savaistu}

\courssub{Phase 2 : La phase asymptomatique (ou latence clinique)}

\paragraph{Observation : le calme apparent.}
Après la guérison du syndrome pseudo-grippal, le patient se sent \emph{parfaitement bien}. Il travaille, fait du sport, n'a aucun signe clinique. Cette phase dure en moyenne \textbf{8 à 10 ans sans traitement} (très variable selon les individus, la souche virale, la nutrition, les co-infections comme le paludisme ou les helminthiases fréquentes au Cameroun).

\paragraph{Investigation : une guerre silencieuse.}
\begin{itemize}
    \item \textbf{Hypothèse :} Le virus n'est pas éliminé ; il se réplique activement dans les organes lymphoïdes, détruisant insidieusement le stock de CD4.
    \item \textbf{Données :} La charge virale se stabilise au \emph{point de set} (niveau propre à chaque patient, corrélé au pronostic). Le taux de CD4 diminue \textbf{linéairement et inexorablement} d'environ \textbf{50 à 80 cellules/mm³ par an}.
    \item \textbf{Interprétation :} Il y a un équilibre dynamique précaire : production virale $\approx$ clairance virale ; mort des CD4 $\approx$ production médullaire. Mais le thymus (usine à CD4 naïfs) s'épuise et l'architecture des ganglions (site de rencontre Ag/lymphocytes) se fibrose. Le \textbf{répertoire immunitaire} s'appauvrit : on perd la mémoire contre les antigènes anciens (varicelle, tuberculose latente, vaccins).
    \item \textbf{Conclusion :} Le sujet est \textbf{asymptomatique mais infectieux}. C'est la période où il contamine le plus son entourage par méconnaissance de son statut.
\end{itemize}

\begin{exemplebox}[Cas concret : Paul, 28 ans, Douala]
Paul se contamine en 2015. Infection primaire discrète (fièvre 3 jours, traitée « palu » en automédication). De 2016 à 2023, il ne consulte pas, se porte bien. En 2023, une angine résistante l'amène à l'hôpital : CD4 = 210 cellules/mm³, Charge virale = 450 000 copies/mL. Il est en phase 3 (Sida) sans avoir eu de signes avant-coureurs. \textbf{Leçon :} L'absence de symptômes ne signifie pas absence de maladie ni absence de contagiosité.
\end{exemplebox}

\begin{savaistu}[Durée de la phase asymptomatique au Cameroun]
Sans traitement, la médiane de survie après séroconversion est d'environ \textbf{8 à 10 ans} dans nos contextes tropicaux (co-infections, malnutrition, accès aux soins variables), contre 10--12 ans dans les pays du Nord. La tuberculose, endémique, est souvent la première maladie révélatrice du Sida chez nos patients (co-infection VIH/TB = 30--40\% des cas TB au Cameroun).
\end{savaistu}

\courssub{Phase 3 : Le Sida déclaré (Stade d'immunodépression sévère)}

\paragraph{Définition clinique et biologique.}
On parle de \textbf{Sida (Syndrome d'Immunodéficience Acquise)} lorsque l'immunité cellulaire s'effondre au point de ne plus contrôler des agents pathogènes habituellement inoffensifs ou latents.

\begin{propriete}[Critères de définition du Sida (OMS / MINESEC)]
Le Sida est défini par l'association de :
\begin{enumerate}
    \item \textbf{Critère biologique majeur :} Taux de lymphocytes T CD4 \textbf{$< 200$ cellules/mm³} (ou $< 14\%$ des lymphocytes totaux).
    \item \textbf{ET/OU Critère clinique :} Survenue d'au moins une \textbf{maladie opportuniste (MO)} définissante (stade 3 ou 4 OMS).
\end{enumerate}
\end{propriete}

\begin{center}
\begin{tabularx}{\linewidth}{|l|X|}\hline
\rowcolor{popblueL}
\textbf{Maladie Opportuniste (Exemples fréquents au Cameroun)} & \textbf{Commentaire / Contexte local} \\ \hline
Tuberculose (pulmonaire \& extrapulmonaire) & 1\textsuperscript{re} cause de décès VIH au Cameroun ; diagnostic difficile (bacilloscopie souvent négative). \\ \hline
Méningite à \emph{Cryptococcus neoformans} & Fréquente si CD4 $< 100$ ; léthale sans traitement antifongique (Amphotéricine B + Flucytosine, souvent indisponibles). \\ \hline
Pneumocystose (\emph{Pneumocystis jirovecii}) & Pneumonie interstitielle diffuse, dyspnée majeure ; prophylaxie par Cotrimoxazole (Bactrim\textsuperscript{\textregistered}) recommandée si CD4 $< 350$. \\ \hline
Sarcome de Kaposi & Tumeur vasculaire (HHV-8), nodules cutanés violacés, atteintes viscérales ; fréquent chez l'homme jeune. \\ \hline
Toxoplasmose cérébrale & Abcès cérébraux multiples, signes focaux ; CD4 $< 100$. \\ \hline
Candidose œsophagienne & Dysphagie, odynophagie ; signe d'alerte précoce (CD4 $< 200$). \\ \hline
Diarrhées chroniques (\emph{Isospora}, \emph{Cryptosporidium}, \emph{Microsporidium}) & Amaigrissement majeur (syndrome amaigrissant), déshydratation. \\ \hline
\end{tabularx}
\end{center}

\begin{aretenir}[Seuil des 200 CD4 : la ligne rouge]
Un taux de CD4 \textbf{inférieur à 200 cellules/mm³} marque la bascule vers le Sida. En dessous de ce seuil, l'immunité cellulaire ne peut plus contenir les infections latentes (TB, herpès, zona) ni se défendre contre les champignons et parasites environnementaux. Sans ARV, la survie médiane à ce stade est de \textbf{quelques mois à 1--2 ans}.
\end{aretenir}

\courssub{Notion de charge virale et interprétation croisée Charge Virale / CD4}

\paragraph{Définition.}
La \textbf{charge virale (CV)} est la quantification de l'ARN du VIH dans le plasma sanguin, exprimée en \textbf{copies/mL} (ou log$_{10}$ copies/mL). Elle reflète l'activité de réplication du virus \emph{in vivo}.

\begin{itemize}
    \item \textbf{Infection primaire :} CV très élevée ($10^6$--$10^7$ copies/mL).
    \item \textbf{Phase asymptomatique :} CV stable au \emph{point de set} ($10^3$--$10^5$ copies/mL). Plus le point de set est haut, plus l'évolution vers le Sida est rapide.
    \item \textbf{Sida :} CV souvent très élevée à nouveau ($> 10^5$ copies/mL) car le contrôle immunitaire (CD8) s'effondre avec les CD4.
\end{itemize}

\paragraph{Relation inverse Charge Virale / Taux de CD4.}
C'est la clé de lecture du pronostic et du suivi thérapeutique.

\begin{propriete}[Relation inverse CV / CD4]
\begin{itemize}
    \item \textbf{Charge virale $\uparrow$} $\Longrightarrow$ Destruction massive des CD4 $\Longrightarrow$ \textbf{Taux CD4 $\downarrow$}.
    \item \textbf{Charge virale $\downarrow$} (sous ARV efficaces) $\Longrightarrow$ Arrêt de la destruction + régénération thymique/périphérique $\Longrightarrow$ \textbf{Taux CD4 $\uparrow$}.
\end{itemize}
Cette relation est le fondement du \textbf{suivi biologique} : on vise \textbf{Charge Virale Indétectable ($< 50$ copies/mL) ET CD4 $> 500$ cellules/mm³}.
\end{propriete}

\begin{methode}[Interpréter un graphique d'évolution CV / CD4 dans le temps]
Pour analyser un graphique à deux courbes (CV en log, CD4 en linéaire) représentant l'histoire naturelle ou le suivi sous traitement :
\begin{enumerate}
    \item \textbf{Repérer les axes :} Temps (abscisses) ; Charge virale (ordonnée gauche, échelle logarithmique) ; Taux CD4 (ordonnée droite, échelle linéaire).
    \item \textbf{Identifier les 3 phases :} Pic initial (infection primaire) $\to$ Plateau/descente lente (phase asymptomatique) $\to$ Remontée CV / Chute CD4 (Sida).
    \item \textbf{Comparer les sens de variation :} Vérifier la relation inverse (CV haut $\leftrightarrow$ CD4 bas).
    \item \textbf{Repérer la séroconversion :} Moment où la CV chute pour la 1\textsuperscript{re} fois (réponse immunitaire).
    \item \textbf{Si traitement (ARV) :} Repérer la chute brutale de la CV (2--4 semaines) et la remontée lente des CD4 (mois/années).
    \item \textbf{Conclure :} Statut immunologique (normo / immuno-déprimé / Sida), efficacité thérapeutique, risque de transmission (I=I : Indétectable = Intransmissible).
\end{enumerate}
\end{methode}

\begin{popfigure}
\begin{tikzpicture}
% Axis 1: Viral Load (Left, Log)
\begin{axis}[
    width=\linewidth,
    height=10cm,
    xlabel={Temps (années après contamination)},
    ylabel={Charge virale (copies/mL)},
    ymin=1, ymax=10000000,
    ymode=log,
    log ticks with fixed point,
    ytick={1,10,100,1000,10000,100000,1000000,10000000},
    yticklabels={$10^0$,$10^1$,$10^2$,$10^3$,$10^4$,$10^5$,$10^6$,$10^7$},
    xmin=0, xmax=10,
    xtick={0,1,2,3,4,5,6,7,8,9,10},
    grid=major,
    grid style={popline},
    legend style={at={(0.02,0.98)}, anchor=north west, fill=popcream, draw=popdark, font=\small, cells={anchor=west}},
    axis line style={popdark},
    tick style={popdark},
    label style={popdark},
    axis y line*=left,
    axis x line*=bottom,
]
% Viral Load Curve
\addplot[popcurve, poporange, very thick, domain=0:10, samples=200] 
    ({x}, 
    { 
        (x < 0.2) * (1000000 * exp(-20*x) + 10000) + 
        (x >= 0.2 && x < 0.5) * (5000000 * exp(-10*(x-0.2)) + 50000) +
        (x >= 0.5 && x < 8) * (30000 + 20000*sin(deg(2*x)) + 5000*x) +
        (x >= 8) * (30000 * exp(0.8*(x-8))) 
    }
);
\addlegendentry{Charge Virale}

% Phase shading (on VL axis)
\addplot[poporange!15, fill opacity=0.3, domain=0:0.5, samples=2] {10000000} \closedcycle;
\addplot[popgreen!15, fill opacity=0.3, domain=0.5:8, samples=2] {10000000} \closedcycle;
\addplot[poppink!15, fill opacity=0.3, domain=8:10, samples=2] {10000000} \closedcycle;

% Phase labels
\node[poporange, anchor=south, font=\small\bfseries] at (axis cs:0.25, 8000000) {Phase 1 : Infection primaire};
\node[popgreen, anchor=south, font=\small\bfseries] at (axis cs:4, 8000000) {Phase 2 : Asymptomatique};
\node[poppink, anchor=south, font=\small\bfseries] at (axis cs:9, 8000000) {Phase 3 : Sida};

% Seroconversion arrow
\draw[poppurple, thick, -Stealth] (axis cs:0.3, 100000) -- (axis cs:0.3, 10000);
\node[poppurple, anchor=east, font=\small] at (axis cs:0.28, 50000) {Séroconversion};

% VL Detection limit
\draw[popmuted, dotted, thick] (axis cs:0, 50) -- (axis cs:10, 50) node[right, popmuted, font=\small] {Limite détection (50 cp/mL)};

\end{axis}

% Axis 2: CD4 Count (Right, Linear)
\begin{axis}[
    width=\linewidth,
    height=10cm,
    ylabel={Taux CD4 (cellules/mm³)},
    ymin=0, ymax=900,
    ymode=linear,
    xmin=0, xmax=10,
    xtick=\empty,
    xlabel=\empty,
    grid=none,
    axis y line*=right,
    axis x line=none,
    legend style={at={(0.98,0.98)}, anchor=north east, fill=popcream, draw=popdark, font=\small, cells={anchor=west}},
    axis line style={popdark},
    tick style={popdark},
    label style={popdark},
]
% CD4 Curve
\addplot[popcurve, popgreen, very thick, domain=0:10, samples=200] 
    ({x}, 
    { 
        (x < 0.2) * (800 - 600*x/0.2) + 
        (x >= 0.2 && x < 0.5) * (200 + 400*(x-0.2)/0.3) +
        (x >= 0.5 && x < 8) * (600 - 50*(x-0.5)) +
        (x >= 8) * (225 - 30*(x-8)) 
    }
);
\addlegendentry{Taux CD4}

% CD4 Threshold lines
\draw[popdark, dashed, thick] (axis cs:0, 200) -- (axis cs:10, 200) node[right, popdark, font=\small] {Seuil Sida (200)};
\draw[popblue, dashed, thick] (axis cs:0, 500) -- (axis cs:10, 500) node[right, popblue, font=\small] {Norme basse (500)};

\end{axis}
\end{tikzpicture}
\legende{Évolution conjointe de la charge virale (orange, échelle logarithmique, axe gauche) et du taux de lymphocytes CD4 (vert, échelle linéaire, axe droit) au cours de l'histoire naturelle de l'infection à VIH (sans traitement). Les trois phases sont clairement individualisées. Notez la relation inverse : quand la charge virale remonte (fin phase 2, phase 3), les CD4 s'effondrent.}
\end{popfigure}

\courssub{Impact de la prise en charge antirétrovirale (ARV) sur l'évolution}

\paragraph{Observation : changer le cours de l'histoire.}
Depuis l'accès aux ARV (au Cameroun, gratuité depuis 2007, protocole « Test and Treat » depuis 2016), le schéma ci-dessus n'est plus une fatalité. Le traitement bloque la réplication virale.

\paragraph{Mécanisme d'action des ARV (rappel).}
Les associations trithérapies (ex: Ténofovir + Lamivudine + Dolutégravir = TLD, 1\textsuperscript{re} ligne au Cameroun) ciblent la transcriptase inverse, l'intégrase ou la protéase. Elles empêchent la production de nouveaux virions infectieux.

\paragraph{Évolution sous traitement efficace (Observance $\ge$ 95\%).}
\begin{enumerate}
    \item \textbf{Semaines 1--4 :} Chute exponentielle de la charge virale (2 à 3 log). Objectif : \textbf{Indétectabilité ($< 50$ copies/mL) à 6 mois}.
    \item \textbf{Mois 1--12 :} Remontée des CD4 (phase de redistribution puis production thymique). Gain moyen : $+100$ à $+150$ cellules/mm³ la 1\textsuperscript{re} année.
    \item \textbf{Années suivantes :} Normalisation progressive du taux CD4 (souvent $> 500$), restauration du répertoire immunitaire, disparition des maladies opportunistes, espérance de vie \textbf{quasi normale}.
\end{enumerate}

\begin{aretenir}[I = I : Indétectable = Intransmissible]
Une personne sous ARV efficace, avec charge virale indétectable durablement ($> 6$ mois), \textbf{ne transmet pas le VIH par voie sexuelle}. C'est une avancée scientifique majeure (études PARTNER, HPTN 052) qui déstigmatise et motive l'observance. Au Cameroun, c'est le message clé des campagnes « \emph{Undetectable = Untransmittable} » portées par le CNLS et les associations de PVVIH.
\end{aretenir}

\begin{attention}[Pièges de l'observance et échec virologique]
\begin{itemize}
    \item \textbf{Oubli de prises :} Le VIH mute vite. Une réplication résiduelle sous pression sélective $\rightarrow$ \textbf{résistances} (ex: mutation M184V pour la Lamivudine, K103N pour l'Efavirenz).
    \item \textbf{Échec virologique :} CV $> 1000$ copies/mL à 6 mois ou rebond confirmé. $\rightarrow$ Nécessite un \textbf{test de résistance} (génotypage) et un changement de ligne (2\textsuperscript{e} ligne : PI/r ou Dolutégravir si 1\textsuperscript{re} ligne NNRTI).
    \item \textbf{Reconstitution immunitaire inflammatoire (IRIS) :} Au démarrage des ARV (souvent CD4 $< 100$), la remontée brutale des CD8 peut déclencher une inflammation paradoxale contre des antigènes microbiens (TB, Cryptocoque). Ce n'est pas un échec du traitement, mais un signe de reprise immunitaire à prendre en charge (corticoïdes parfois).
\end{itemize}
\end{attention}

\courssub{Exercice résolu : Interprétation d'un suivi biologique}

\begin{exoresolu}[Suivi de Mireille, 32 ans, diagnostiquée VIH+ en 2022]
\textbf{Énoncé :} Mireille découvre sa séropositivité lors d'une consultation prénatale (CPN) en janvier 2022. Elle démarre immédiatement la trithérapie TLD (Ténofovir/Lamivudine/Dolutégravir). Voici son suivi biologique :

\begin{center}
\begin{tabular}{|c|c|c|c|}\hline
\textbf{Date} & \textbf{Charge Virale (copies/mL)} & \textbf{CD4 (cellules/mm³)} & \textbf{Observance déclarée} \\ \hline
Janv 2022 (J0) & 185 000 & 320 & -- \\ \hline
Mars 2022 (M2) & 450 & 410 & Bonne \\ \hline
Juil 2022 (M6) & \textbf{$< 20$ (Indétectable)} & 480 & Bonne \\ \hline
Janv 2023 (M12) & \textbf{$< 20$} & 580 & Bonne \\ \hline
Juil 2023 (M18) & \textbf{1 200} & 550 & \textbf{Irrégulière (oubli fréquent)} \\ \hline
\end{tabular}
\end{center}

\textbf{Questions :}
\begin{enumerate}
    \item Situer le stade clinique au diagnostic (janv 2022) et justifier.
    \item Analyser l'évolution entre J0 et M6. Le traitement est-il efficace ?
    \item Interpréter les résultats de M18. Quelle hypothèse formuler ? Quelle conduite tenir ?
\end{enumerate}

\tcblower
\textbf{Solution détaillée :}

\begin{enumerate}
    \item \textbf{Stade au diagnostic (Janv 2022) :} \textbf{Phase asymptomatique (Stade 2 OMS)}.
        \begin{itemize}
            \item CD4 = 320 cellules/mm³ $\rightarrow$ $> 200$, pas de Sida biologiquement.
            \item Pas de signe clinique majeur évoqué (découverte à la CPN).
            \item Charge virale élevée (185 000 cp/mL) typique du point de set en phase asymptomatique.
        \end{itemize}
    \item \textbf{Évolution J0 $\rightarrow$ M6 : Succès virologique et immunologique.}
        \begin{itemize}
            \item \textbf{Virologique :} CV passe de 185 000 à $< 20$ (indétectable) en 6 mois. Chute $> 4$ log. Objectif atteint (CV $< 50$ cp/mL à M6). Cela prouve l'efficacité de la trithérapie et une bonne observance.
            \item \textbf{Immunologique :} CD4 passent de 320 à 480 ($+160$ en 6 mois). Remontée rapide typique de la phase de redistribution (lymphocytes mémoires revenant du compartiment lymphoïde vers le sang). Le seuil de 500 est presque atteint.
            \item \textbf{Conclusion :} Réponse thérapeutique optimale. Mireille est \textbf{indétectable = intransmissible (I=I)} pour son conjoint et son enfant à naître (prophylaxie mère-enfant optimale).
        \end{itemize}
    \item \textbf{Interprétation M18 : Échec virologique suspecté sur non-observance.}
        \begin{itemize}
            \item \textbf{Donnée :} CV rebondit à 1 200 cp/mL (confirmation nécessaire par 2\textsuperscript{e} dosage à 1 mois), CD4 stables (550).
            \item \textbf{Hypothèse :} \textbf{Non-observance} (oubli fréquent déclaré) $\rightarrow$ Pression sélective insuffisante $\rightarrow$ Réplication virale résiduelle $\rightarrow$ Risque de sélection de mutants résistants (surtout sur Dolutégravir : barrière génétique haute mais pas infinie ; sur 3TC : mutation M184V rapide).
            \item \textbf{Conduite à tenir (Algorithme OMS/Cameroun) :}
                \begin{enumerate}
                    \item \textbf{Conseil d'observance renforcé (Adhérence) :} Éducation thérapeutique, boîtier pilulier, soutien pair/psychosocial. C'est la 1\textsuperscript{re} action.
                    \item \textbf{Contrôle CV de confirmation à 1 mois (M19) :} Si CV redescend $< 50$ $\rightarrow$ « Blip » (rebond transitoire) ou reprise observance $\rightarrow$ Maintenir 1\textsuperscript{re} ligne.
                    \item \textbf{Si CV confirmé $> 1000$ cp/mL à M19 :} \textbf{Échec virologique avéré.} Demander \textbf{génotypage de résistance} (si disponible) et passer en \textbf{2\textsuperscript{e} ligne} (ex: Ténofovir + Lamivudine + Atazanavir/ritonavir ou Darunavir/ritonavir si résistance aux INI).
                \end{enumerate}
        \end{itemize}
\end{enumerate}
\end{exoresolu}

\begin{exercice}[Auto-évaluation : L'histoire naturelle et le traitement]
\textbf{1.} Un patient a un taux de CD4 à 150 cellules/mm³ et une charge virale à 500 000 copies/mL. Il n'a jamais été traité. Précisez son stade clinique OMS et citez 3 maladies opportunistes à rechercher en priorité dans notre contexte camerounais.
\textbf{2.} Expliquez pourquoi la charge virale est exprimée en échelle logarithmique (log$_{10}$) sur les graphiques de suivi.
\textbf{3.} Un patient sous ARV depuis 2 ans a une CV indétectable et des CD4 à 650. Il arrête son traitement de son propre chef pendant 3 semaines. Que va-t-il se passer sur le plan virologique et immunologique ? Quel risque pour son entourage ?
\textbf{4.} À l'aide du graphique de la figure ci-dessus (histoire naturelle), estimez la durée moyenne de la phase asymptomatique et le taux de déclin annuel moyen des CD4 durant cette phase (en cellules/mm³/an).
\end{exercice}

\begin{corrige}[Exercice Auto-évaluation]
\textbf{1.} Stade \textbf{3 (Sida)} car CD4 $< 200$. MO prioritaires au Cameroun : \textbf{Tuberculose} (recherche BK au crachats + GeneXpert), \textbf{Méningite à Cryptocoques} (recherche Ag cryptococcique sang/LCS si CD4 $< 100$ ou signes méningés), \textbf{Pneumocystose} (radio poumon + recherche P. jirovecii si dyspnée), \textbf{Candidose œsophagienne} (si dysphagie).
\textbf{2.} La charge virale varie sur plusieurs ordres de grandeur ($10^2$ à $10^7$). L'échelle log permet de visualiser les variations relatives (ex: chute de 2 log = division par 100) sur un même graphique ; en linéaire, le pic initial écrabouillerait la phase basse.
\textbf{3.} \textbf{Rebond viral rapide :} CV redevient détectable en quelques jours (réservoirs latents + réplication résiduelle), remonte vers le point de set initial en 2--4 semaines. \textbf{CD4 :} Commencent à rebaisser (perte du contrôle viral). \textbf{Risque transmission :} Réapparaît dès que CV $> 200$--$1000$ cp/mL. Perte du bénéfice I=I. Reprise ARV d'urgence + conseil observance.
\textbf{4.} Lecture graphique : Phase 2 $\approx$ de 0.5 à 8 ans $\rightarrow$ \textbf{7.5 ans}. CD4 : départ $\approx$ 600 (à 0.5 an) $\rightarrow$ arrivée $\approx$ 225 (à 8 ans). Déclin = $(600 - 225) / 7.5 = 375 / 7.5 = \textbf{50 cellules/mm³/an}$. (Compatible avec la littérature : 50--80/an).
\end{corrige}

\coursec{Dépistage, prise en charge et traitements}

Après avoir suivi l'évolution naturelle de l'infection à VIH — de la primo-infection souvent silencieuse à la phase de Sida déclaré marquée par l'effondrement des lymphocytes CD4 et l'apparition d'infections opportunistes —, une question s'impose avec urgence : \textbf{comment intervenir pour briser cette évolution fatale ?} La réponse repose sur deux piliers indissociables : le \textbf{dépistage} pour connaître son statut sérologique le plus tôt possible, et la \textbf{prise en charge thérapeutique} par les antirétroviraux (ARV) pour bloquer la réplication virale. Au Cameroun, comme partout dans le monde, l'accès universel à ces deux piliers a transformé le VIH d'une maladie mortelle en une pathologie chronique contrôlée.

\courssub{Le dépistage du VIH : principe, déroulement et enjeux}

\paragraph{Observation — Situation concrète.}
Un jeune adulte se rend dans un Centre de Dépistage Volontaire (CDV) à Yaoundé après une prise de risque. Il est anxieux : « \emph{Et si j'étais positif ? Combien de temps avant de le savoir ? Est-ce que tout le monde le saura ? } » Ces questions reflètent les freins réels au dépistage : la peur du résultat, la méconnaissance des délais (fenêtre sérologique) et la crainte de la stigmatisation.

\paragraph{Principe biologique : la recherche des anticorps anti-VIH (et de l'antigène p24).}
Le VIH ne se détecte pas directement par observation au microscope dans le sang (trop petit, trop peu nombreux hors culture). Le dépistage de routine repose sur la \textbf{sérologie} : on recherche les \textbf{anticorps} produits par le système immunitaire du patient \emph{contre} le VIH (anti-gp120, anti-gp41, anti-p24). Les tests de \textbf{4\up{ème} génération (combinés Ag/Ab)}, aujourd'hui standard au Cameroun, détectent \textbf{simultanément l'antigène p24 (protéine de capside, précoce) et les anticorps anti-VIH-1/2}, raccourcissant la fenêtre sérologique.

\begin{definition}[Fenêtre sérologique]
Période comprise entre la contamination et l'apparition des marqueurs viraux (antigène p24 et/ou anticorps anti-VIH) détectables par les tests de dépistage. Avec les tests de 4\up{ème} génération, elle dure en moyenne \textbf{2 à 6 semaines}. Pendant cette fenêtre, le test est \textbf{négatif alors que la personne est contaminée et contaminante} (charge virale très élevée, pic de virémie).
\end{definition}

\paragraph{Démarche d'investigation : les étapes du diagnostic sérologique.}
\begin{enumerate}
    \item \textbf{Test de dépistage (TROD ou ELISA 4\up{ème} génération) :} Recherche combinée antigène p24 + anticorps anti-VIH-1/2.
    \begin{itemize}
        \item \textbf{TROD (Test Rapide d'Orientation Diagnostique) :} Résultat en 15-30 min sur sang total (bout de doigt) ou fluide oral. Idéal pour le dépistage de proximité, en milieu communautaire ou en urgence (accouchement). Sensibilité légèrement inférieure aux ELISA sur prélèvement récent.
        \item \textbf{ELISA (Enzyme Linked Immunosorbent Assay) 4\up{ème} gén. :} Technique de laboratoire sur sérum ou plasma, haute sensibilité/spécificité, adaptée au dépistage de masse (dons de sang, PTNE, Consultations Prénatales).
    \end{itemize}
    \item \textbf{Test de confirmation (Immunoblot / Western Blot ou 2\up{ème} ELISA à principe différent) :} Obligatoire si le test de dépistage est \textbf{réactif (positif)}. L'immunoblot sépare les protéines du VIH par électrophorèse et visualise les anticorps dirigés contre \emph{chaque} protéine virale (gp160, gp120, gp41, p24, p17, p31, p51, p66). Le test est \textbf{confirmé positif} seulement si des bandes spécifiques sont présentes selon les critères OMS (ex: au moins 2 bandes \emph{env} [gp160/120/41] OU 1 bande \emph{env} + 1 bande \emph{gag} [p24/p17] OU 1 bande \emph{env} + 1 bande \emph{pol} [p31/p51/p66]).
    \item \textbf{Biologie moléculaire (PCR-ADN/ARN VIH) :} En cas de résultat indéterminé (Immunoblot incomplet) ou pour le \textbf{dépistage précoce du nourrisson} né de mère séropositive (les anticorps maternels traversent le placenta et faussent la sérologie jusqu'à 18 mois). La PCR recherche le \textbf{génome viral (ARN ou ADN proviral)} directement. C'est le standard pour le diagnostic infantile précoce (DIP) au Cameroun (prélèvement à 6 semaines, 9 mois, fin allaitement).
\end{enumerate}

\begin{popfigure}
\begin{tikzpicture}[
    node distance=1.0cm and 1.2cm,
    box/.style={draw=popblue, thick, rounded corners, fill=popblueL, text width=3.0cm, align=center, minimum height=1.6cm},
    arrow/.style={-Stealth, thick, popdark},
    decision/.style={diamond, draw=poporange, thick, fill=poporangeL, text width=2.5cm, align=center, aspect=2, inner sep=2pt}
]
\node[box, align=center] (dep){Dépistage\\ (TROD ou ELISA 4\up{ème} gén.)};
\node[decision, below=of dep] (res) {Résultat ?};
\node[box, left=2.2cm of res, align=center] (neg){Négatif\\ $\rightarrow$ Non infecté*\\(*si hors fenêtre sérologique)};
\node[box, right=2.2cm of res, align=center] (pos){Réactif (Positif)\\ $\rightarrow$ Confirmation};
\node[box, below=of pos, align=center] (conf){Confirmation\\ (Immunoblot / WB)};
\node[decision, below=of conf] (confres) {Résultat Confirmation ?};
\node[box, left=2.2cm of confres, align=center] (neg2){Négatif / Indéterminé\\ $\rightarrow$ Contrôle à J+4-6 sem.\\ ou PCR VIH};
\node[box, right=2.2cm of confres, align=center] (pos2){Positif\\ $\rightarrow$ \textbf{Séropositivité confirmée}};

\draw[arrow] (dep) -- (res);
\draw[arrow] (res) -- node[above, sloped, font=\small] {Négatif} (neg);
\draw[arrow] (res) -- node[above, sloped, font=\small] {Réactif} (pos);
\draw[arrow] (pos) -- (conf);
\draw[arrow] (conf) -- (confres);
\draw[arrow] (confres) -- node[above, sloped, font=\small] {Négatif/Indéterminé} (neg2);
\draw[arrow] (confres) -- node[above, sloped, font=\small] {Positif} (pos2);
\end{tikzpicture}
\legende{Algorithme simplifié du diagnostic sérologique du VIH (tests 4\up{ème} génération). Tout test de dépistage réactif doit être confirmé par un test spécifique (Immunoblot/Western Blot ou 2\up{ème} ELISA) avant d'annoncer une séropositivité.}
\end{popfigure}

\paragraph{Le dépistage : Volontaire, Gratuit, Confidentiel (VGC).}
Au Cameroun, la politique nationale garantit les \textbf{3 C} (Loi n°2021/013) :
\begin{itemize}
    \item \textbf{Consentement :} Nul ne peut être dépisté à son insu (sauf exceptions légales rares : dons de sang, recherche anonyme). Le conseil pré-test est obligatoire pour expliquer le sens du test, la fenêtre sérologique et les implications du résultat.
    \item \textbf{Confidentialité / Secret médical :} Le résultat ne peut être divulgué qu'au patient. Le personnel de santé est tenu au secret professionnel.
    \item \textbf{Gratuité :} Le test et le conseil sont gratuits dans les structures publiques et agréées (CDV, CPEC, hôpitaux, formations sanitaires de base).
\end{itemize}

\begin{methode}[Argumenter en faveur du dépistage volontaire]
Pour convaincre un pair ou rédiger une copie d'examen, structure ton argumentation en \textbf{trois piliers} :
\begin{enumerate}
    \item \textbf{Intérêt individuel (Santé propre) :} Connaître son statut permet une \textbf{prise en charge précoce}. Démarrer les ARV au stade asymptomatique (CD4 $> 500$/mm$^3$) préserve le capital immunitaire, évite les infections opportunistes et assure une espérance de vie quasi normale.
    \item \textbf{Intérêt collectif (Santé publique) :} Une personne dépistée positive et traitée voit sa charge virale devenir indétectable $\rightarrow$ elle \textbf{ne transmet plus le virus par voie sexuelle} (I=I : Indétectable = Intransmissible). Le dépistage de masse réduit donc la circulation du virus dans la population (stratégie \emph{TasP} : Treatment as Prevention).
    \item \textbf{Intérêt familial (PTME) :} Pour une femme enceinte, le dépistage (systématique en CPN au Cameroun) ouvre l'accès à la \textbf{Prévention de la Transmission Mère-Enfant (PTME)} : ARV mère (TLD idéalement) + AZT nourrisson (4-6 sem.) + accouchement sécurisé + conseil alimentation (allaitement maternel exclusif sous ARV maternels efficaces) $\rightarrow$ taux de transmission $< 2\%$ (vs $25-40\%$ sans intervention).
\end{enumerate}
\end{methode}

\courssub{Les antirétroviraux (ARV) : bloquer le cycle viral à plusieurs étages}

\paragraph{Observation — Le défi thérapeutique.}
Le VIH est un rétrovirus : son génome ARN est rétrotranscrit en ADN par la \textbf{transcriptase inverse}, puis intégré dans l'ADN de la cellule hôte par l'\textbf{intégrase}. Les protéines virales sont découpées par la \textbf{protéase}. Cibler une seule enzyme (monothérapie, ex: AZT seul dans les années 90) sélectionne rapidement des mutants résistants. La solution : la \textbf{trithérapie} (association de 3 molécules de classes différentes), standard depuis 1996, aujourd'hui souvent en \textbf{comprimé unique (monocomprimé)} pour faciliter l'observance.

\begin{propriete}[Efficacité de la trithérapie combinée (admise)]
L'association de \textbf{trois antirétroviraux} ciblant \textbf{au moins deux étapes distinctes} du cycle viral (recommandation OMS 1\up{ère} ligne : 2 INTI + 1 II, ex: Ténofovir + Lamivudine + Dolutégravir = TLD) permet :
\begin{itemize}
    \item Une \textbf{suppression virologique durable} : charge virale plasmatique $< 50$ copies/mL (seuil de détection standard) = \textbf{charge virale indétectable}.
    \item L'\textbf{arrêt de la destruction des lymphocytes CD4} et leur remontée progressive (reconstitution immunitaire).
    \item La \textbf{prévention de l'émergence de résistances} : la probabilité qu'un virus mute simultanément sur les 3 cibles est infinitésimale ($10^{-15}$ à $10^{-18}$).
\end{itemize}
\textit{Note : Cette propriété est admise au Probatoire. On ne demande pas de démontrer la cinétique virale, mais de l'interpréter à partir de graphiques ou tableaux.}
\end{propriete}

\paragraph{Les cibles thérapeutiques majeures (Classes d'ARV).}
\begin{center}
\begin{tabularx}{\linewidth}{|l|X|X|}\hline
\rowcolor{popblueL}
\textbf{Classe (Abréviation)} & \textbf{Cible virale / Étape bloquée} & \textbf{Exemples de molécules (Cameroun, 1\up{ère}/2\up{ème} ligne)} \\ \hline
\textbf{INTI} (Inhibiteurs Nucléosidiques de la Transcriptase Inverse) & \textbf{Transcriptase Inverse} : analogues de nucléosides, terminateurs de chaîne lors de la synthèse de l'ADN viral. & Ténofovir (TDF), Lamivudine (3TC), Emtricitabine (FTC), Zidovudine (AZT), Abacavir (ABC) \\ \hline
\textbf{INNTI} (Inhibiteurs Non Nucléosidiques de la TI) & \textbf{Transcriptase Inverse} : fixation allostérique $\rightarrow$ changement conformation $\rightarrow$ inactivation. & Névirapine (NVP), Éfavirenz (EFV), Doravirine (DOR), Rilpivirine (RPV) \\ \hline
\textbf{IP} (Inhibiteurs de la Protéase) & \textbf{Protéase} : empêchent le clivage des polyprotéines Gag-Pol $\rightarrow$ virions immatures, non infectieux. (Utilisés boostés par ritonavir /r ou cobicistat /c). & Lopinavir/r (LPV/r), Atazanavir/r (ATV/r), Darunavir/r (DRV/r) \\ \hline
\textbf{II} (Inhibiteurs de l'Intégrase) & \textbf{Intégrase} : bloquent l'intégration de l'ADN viral dans le génome hôte $\rightarrow$ ADN viral épisomal dégradé. \textbf{Barrière génétique haute (DTG, BIC).} & Dolutegravir (DTG), Bictegravir (BIC), Raltégravir (RAL) \\ \hline
\textbf{IE} (Inhibiteurs d'Entrée / Fixation / Fusion) & \textbf{Co-récepteur CCR5} (Maraviroc) ou \textbf{Fusion gp41} (Enfuvirtide). & Maraviroc (MVC), Enfuvirtide (T-20, injectable) \\ \hline
\textbf{IC} (Inhibiteurs de Capside) & \textbf{Capside (p24)} : interfèrent avec l'assemblage, le désassemblage et le transport nucléaire. Action prolongée. & Lenacapavir (LEN, injection 2x/an) \\ \hline
\end{tabularx}
\end{center}

\begin{popfigure}
\begin{tikzpicture}[
    scale=0.85,
    every node/.style={font=\small},
    virus/.style={draw=popdark, thick, fill=popcream, rounded corners=3pt, minimum height=1.4cm, text width=2.8cm, align=center},
    arrow/.style={-Stealth, thick, popdark},
    target/.style={draw=popgreen, thick, fill=popgreenL, circle, minimum size=1.6cm},
    inhib/.style={draw=popblue, thick, fill=popblueL, rectangle, rounded corners=2pt, minimum width=2.5cm, minimum height=0.8cm, align=center}
]
% Cycle viral simplifié (haut)
\node[virus, align=center] (entree) at (0,4){1. Entrée\\ Fixation CD4/CCR5 $\rightarrow$ Fusion};
\node[virus, align=center] (rt) at (5.5,4){2. Transcription inverse\\ ARN $\xrightarrow{\text{TI}}$ ADN};
\node[virus, align=center] (int) at (11,4){3. Intégration\\ ADN viral $\xrightarrow{\text{Intégrase}}$ ADN hôte};
\node[virus, align=center] (trans) at (16.5,4){4. Transcription/Traduction\\ ARNm $\rightarrow$ Polyprotéines};
\node[virus, align=center] (mat) at (22,4){5. Assemblage/Maturation\\ Protéase $\rightarrow$ Virions matures};
\node[virus] (sortie) at (27.5,4) {6. Bourgeonnement/Sortie};

\draw[arrow] (entree) -- (rt);
\draw[arrow] (rt) -- (int);
\draw[arrow] (int) -- (trans);
\draw[arrow] (trans) -- (mat);
\draw[arrow] (mat) -- (sortie);

% Cibles ARV (bas)
\node[target, align=center] (c1) at (2.75,1.2){IE\\ (Maraviroc,\\ Enfuvirtide)};
\node[target, align=center] (c2) at (8.25,1.2){INTI / INNTI\\ (TDF, 3TC, FTC,\\ EFV, DOR, NVP)};
\node[target, align=center] (c3) at (13.75,1.2){II\\ (Dolutegravir,\\ Bictegravir, RAL)};
\node[target, align=center] (c4) at (19.25,1.2){IP\\ (LPV/r, DRV/r,\\ ATV/r)};
\node[target, align=center] (c5) at (24.75,1.2){IC\\ (Lenacapavir)};

% Flèches liaison
\draw[popblue, dashed, thick] (entree.south) -- (c1.north) node[midway, right, font=\tiny] {Blocage entrée};
\draw[popblue, dashed, thick] (rt.south) -- (c2.north) node[midway, right, font=\tiny] {Blocage TI};
\draw[popblue, dashed, thick] (int.south) -- (c3.north) node[midway, right, font=\tiny] {Blocage Intégration};
\draw[popblue, dashed, thick] (mat.south) -- (c4.north) node[midway, right, font=\tiny] {Blocage Maturation};
\draw[popblue, dashed, thick] (sortie.south) -- (c5.north) node[midway, right, font=\tiny] {Blocage Capside};

% Légende
\node[align=center, font=\footnotesize, text width=2.5cm] at (-3, 1.2) {\textbf{Classes d'ARV}\\\textit{(Point d'action)}};
\node[align=center, font=\footnotesize, text width=2.5cm] at (-3, 4) {\textbf{Cycle de réplication du VIH}\\\textit{(dans le lymphocyte CD4)}};
\end{tikzpicture}
\legende{Points d'action des principales classes d'antirétroviraux (ARV) sur le cycle de réplication du VIH. La 1\up{ère} ligne OMS (TLD) combine 2 INTI (TDF+3TC) + 1 II (DTG) en un seul comprimé/jour. *Lenacapavir : nouvelle classe (inhibiteur de capside), injection sous-cutanée tous les 6 mois, pour échec thérapeutique ou prévention (PrEP).}
\end{popfigure}

\paragraph{Évolution sous traitement : les marqueurs de succès.}
L'efficacité du traitement se suit par deux examens biologiques clés, réalisés tous les 3 à 6 mois :
\begin{enumerate}
    \item \textbf{Charge virale (CV) VIH-ARN (copies/mL) :} Doit devenir \textbf{indétectable ($< 50$ copies/mL)} dès 3-6 mois. C'est le marqueur \textbf{le plus précoce et le plus sensible} de l'efficacité. Seuil OMS de suppression virologique (santé publique) : $< 1000$ copies/mL.
    \item \textbf{Numération des lymphocytes T CD4 (cellules/mm$^3$) :} Remontée progressive (gain moyen $50-150$ cellules/an les premières années). Objectif : $> 500$ cellules/mm$^3$ (normalisation).
\end{enumerate}

\begin{popfigure}
\begin{tikzpicture}
\begin{axis}[
    width=\linewidth, height=8cm,
    xlabel={Temps (années)}, ylabel={CD4 (cellules/mm$^3$)},
    xmin=0, xmax=10, ymin=0, ymax=1100,
    grid=major, grid style={popline},
    legend style={at={(0.02,0.98)}, anchor=north west, fill=popcream, draw=popdark, font=\small},
]
% Sans traitement
\addplot[popcurve, popred, domain=0:10, samples=100] {800 - 60*x - 10*x^2};
\addlegendentry{Sans ARV : chute progressive vers le Sida}

% Avec traitement (début à t=2 ans) - Continuité assurée
\addplot[popcurve, popgreen, domain=0:2, samples=50] {800 - 60*x - 10*x^2};
\addplot[popcurve, popgreen, domain=2:10, samples=100] {640 + 80*(x-2) - 5*(x-2)^2}; % 640 = valeur à t=2
\addlegendentry{Avec ARV (début à 2 ans) : remontée durable}

% Seuils
\addplot[dashed, poporange, thick, domain=0:10] {200};
\node[poporange, right, font=\small] at (axis cs:10,200) {Seuil Sida (200 CD4)};
\addplot[dashed, popblue, thick, domain=0:10] {500};
\node[popblue, right, font=\small] at (axis cs:10,500) {Zone normale ($>500$)};
\end{axis}
\end{tikzpicture}
\legende{Évolution typique de la numération des lymphocytes CD4 avec et sans traitement antirétroviral (ARV). Sans traitement, la chute inexorable franchit le seuil de 200 CD4 (stade Sida). Avec une trithérapie efficace (ici démarrée à 2 ans, CD4 $\approx$ 640), la charge virale devient indétectable et les CD4 remontent vers la normale. Plus le traitement est commencé tôt (CD4 hauts), meilleure est la reconstitution immunitaire.}
\end{popfigure}

\begin{exemplebox}[« Indétectable = Intransmissible » (I = I) : une révolution préventive]
\textbf{Données scientifiques solides (Études PARTNER 1 \& 2, Opposites Attract, HPTN 052) :} Aucune transmission sexuelle du VIH n'a été observée au sein de couples sérodifférents (un partenaire séropositif sous ARV efficace, l'autre séronégatif) lorsque le partenaire séropositif maintenait une \textbf{charge virale indétectable ($< 200$ copies/mL) depuis au moins 6 mois} et une \textbf{observance parfaite}.

\textbf{Exemple chiffré :} Un couple camerounais à Douala. L'homme est séropositif, sous Dolutégravir/TDF/3TC (TLD) depuis 3 ans, CV $< 20$ copies/mL, CD4 $= 680$. Sa femme est séronégative. Ils ont des rapports sexuels non protégés (désir d'enfant). \textbf{Risque de transmission sexuelle : ZÉRO (négligeable, non mesurable).} La femme peut tomber enceinte naturellement sans risque d'infection pour elle ni pour l'enfant (si le père maintient l'indétectabilité).

\textbf{Attention :} I=I ne s'applique \textbf{pas} à la transmission mère-enfant pendant l'allaitement (le lait maternel contient des cellules infectées même si CV plasmatique indétectable) ni au partage de seringues. Le préservatif reste recommandé pour se protéger des autres IST et grossesses non désirées.
\end{exemplebox}

\begin{attention}[Pièges majeurs à éviter — Probatoire et vie réelle]
\begin{itemize}
    \item \textbf{« Les ARV guérissent le VIH. » \textbf{FAUX.}} Les ARV bloquent la réplication mais \textbf{n'éliminent pas les réservoirs viraux} (lymphocytes CD4 mémoire à longue durée de vie, système nerveux central, ganglions, tractus génital). L'arrêt du traitement entraîne \textbf{systématiquement} une reprise virale (rebond) en 2-4 semaines.
    \item \textbf{« Je me sens bien, j'arrête mes médicaments. » \textbf{DANGEREUX.}} L'interruption (même quelques jours) favorise la sélection de \textbf{mutants résistants}. La résistance rend la trithérapie inefficace $\rightarrow$ échec virologique $\rightarrow$ passage aux lignes de traitement de secours (plus chères, plus d'effets secondaires, moins disponibles).
    \item \textbf{Confondre « Charge virale indétectable » et « Séronégativité ».} Le test sérologique (anticorps) \textbf{reste positif à vie} même sous ARV efficaces. Seule la PCR (charge virale) devient négative.
    \item \textbf{Pensée magique : « Il existe un vaccin ou un remède traditionnel qui guérit. »} \textbf{Aucun vaccin préventif ni traitement curatif validé} n'existe à ce jour (recherches en cours : vaccins à ARNm, anticorps broadly neutralizing, stratégies « shock and kill » ou « block and lock » pour vider/contrôler les réservoirs).
\end{itemize}
\end{attention}

\courssub{Observance, résistance et vie avec le VIH : au-delà de la biologie}

\paragraph{L'observance thérapeutique : condition \emph{sine qua non} du succès.}
L'observance (adherence) désigne le fait de \textbf{prendre correctement son traitement : bonne dose, bon moment, tous les jours, à vie}.
\begin{itemize}
    \item \textbf{Seuil critique :} $\ge 95\%$ de prises effectives (pour les anciens ARV comme EFV, NVP, IP/r). Les nouvelles molécules (Dolutégravir, Bictegravir) ont une \textbf{barrière génétique à la résistance élevée} et une longue demi-vie, tolérant mieux quelques oublis occasionnels, mais la règle reste : \textbf{zéro oubli visé}.
    \item \textbf{Conséquences de la non-observance :} Réplication virale en présence de concentrations sous-optimales d'ARV $\rightarrow$ mutations $\rightarrow$ \textbf{résistance} $\rightarrow$ échec thérapeutique $\rightarrow$ transmission de souches résistantes (résistance primaire).
    \item \textbf{Outils d'aide :} Boîtes à pilules (hebdomadaires), alarmes téléphone, applications mobiles, éducation thérapeutique du patient (ETP), groupes de pairs, soutien psychosocial.
\end{itemize}

\begin{exoresolu}[Interprétation de suivi biologique et décision thérapeutique]
\textbf{Énoncé :} Marie, 32 ans, séropositive depuis 5 ans, sous TDF/3TC/DTG (1\up{ère} ligne OMS, TLD) depuis 3 ans. Ses derniers bilans :
\begin{center}
\begin{tabular}{|c|c|c|c|}\hline
\textbf{Date} & \textbf{Charge Virale (copies/mL)} & \textbf{CD4 (cell/mm$^3$)} & \textbf{Observance déclarée} \\ \hline
J-18 mois & $< 20$ (Indétectable) & 580 & Bonne \\ \hline
J-12 mois & $< 20$ & 620 & Bonne \\ \hline
J-6 mois & \textbf{1 250} & 550 & « Parfois j'oublie le week-end » \\ \hline
J (aujourd'hui) & \textbf{45 000} & 420 & « J'ai arrêté 2 semaines le mois dernier » \\ \hline
\end{tabular}
\end{center}
\textbf{Questions :}
\begin{enumerate}
    \item Interpréter l'évolution de la charge virale et des CD4.
    \item Expliquer le mécanisme biologique reliant l'observance à cette évolution.
    \item Quelle décision médicale préconiser ? Justifier.
\end{enumerate}

\textbf{Solution détaillée :}
\begin{enumerate}
    \item \textbf{Évolution :} De J-18 à J-12 mois : succès virologique (CV indétectable $< 50$) et immunologique (CD4 en hausse). À J-6 mois : \textbf{reprise virale / virémie de bas niveau (LLV)} : CV $= 1\,250$ copies/mL ($> 1\,000$ = seuil d'alerte OMS, suspect d'échec virologique). À J : \textbf{échec virologique confirmé} (CV $= 45\,000$ copies/mL $> 1\,000$ confirmé) avec début de chute des CD4 ($420$).
    \item \textbf{Mécanisme :} Les oublis répétés (week-ends) puis l'arrêt de 2 semaines ont créé des « fenêtres » de concentration plasmatique d'ARV \textbf{sous le seuil inhibiteur (Cmin $<$ IC50/IC90)}. Le VIH a pu se répliquer en présence de concentrations sous-optimales de DTG/TDF/3TC $\rightarrow$ \textbf{sélection de mutants résistants} (ex: mutation M184V/I pour 3TC/FTC conférant résistance croisée TDF partielle ; mutation R263K/G140S pour DTG). Ces variants résistants ont pris le dessus sur le virus sauvage sensible.
    \item \textbf{Décision :} \textbf{Ne pas changer de traitement aveuglément.} Réaliser un \textbf{génotypage de résistance} (séquençage du gène \emph{pol} : RT, Protéase, Intégrase) sur le prélèvement J. En attendant les résultats (2-3 semaines) : \textbf{conseil d'observance intensif} (éducation thérapeutique, boîtes à pilules, soutien psychosocial, recherche cause : effets secondaires, stigmatisation, rupture de stock, santé mentale). 
    \begin{itemize}
        \item Si résistances majeures confirmées (ex: M184V + R263K) $\rightarrow$ passage en \textbf{2\up{ème} ligne} (ex: AZT/3TC + LPV/r ou DRV/r + DTG selon profil de résistance et disponibilité).
        \item Si pas de mutation majeure (échec purement lié à l'observance, virus sauvage) $\rightarrow$ \textbf{maintien de la 1\up{ère} ligne (TLD)} avec accompagnement renforcé (le DTG a une barrière génétique haute, le virus sauvage sensible peut réapparaître et redevenir dominant si l'observance redevient parfaite).
    \end{itemize}
\end{enumerate}
\end{exoresolu}

\paragraph{Lutte contre la stigmatisation et la discrimination.}
La séropositivité n'est \textbf{pas une condamnation à mort}, ni une faute morale. C'est une \textbf{maladie chronique} comme le diabète ou l'hypertension, qui se soigne.
\begin{itemize}
    \item \textbf{Droits :} Accès aux soins, au travail, à l'éducation, à la vie familiale (mariage, enfants) sans discrimination (Loi n°2021/013 au Cameroun relative à la protection des personnes vivant avec le VIH/Sida).
    \item \textbf{Rôle de l'entourage :} Soutien moral, rappel des rendez-vous, aide à l'observance (« buddy »), non-divulgation du statut sans consentement.
    \item \textbf{Groupes de parole / Associations :} (ex: CAMNAFAW, RECAJ+, Réseau des Associations de PVVIH du Cameroun) : brisent l'isolement, défendent les droits, forment des pairs-éducateurs.
\end{itemize}

\begin{savaistu}[Accès aux ARV au Cameroun : une histoire de solidarité]
Depuis \textbf{2000}, le Cameroun s'engage dans l'accès universel aux ARV.
\begin{itemize}
    \item \textbf{2000-2007 :} Accès via l'Initiative « Accès aux soins » (MSF, ESTHER, Fonds Mondial) puis gratuité progressive dans le secteur public.
    \item \textbf{2007 :} Décret du Premier Ministre rendant les ARV \textbf{gratuits} dans toutes les structures publiques et agréées (CPEC - Centres de Prise en Charge).
    \item \textbf{Aujourd'hui :} La \textbf{1\up{ère} ligne recommandée par l'OMS (TLD : Ténofovir + Lamivudine + Dolutégravir)} est disponible gratuitement en monocomprimé (1 cp/jour). Peu d'effets secondaires, barrière génétique haute, efficacité supérieure.
    \item \textbf{Défi actuel :} Assurer la \textbf{disponibilité continue} (lutte contre les ruptures de stock), la \textbf{prise en charge des enfants} (formulations pédiatriques dispersibles, pDTG), et l'accès à la \textbf{2\up{ème}/3\up{ème} ligne} (plus coûteuses, ex: DRV/r, DTG en 2\up{ème} ligne, Lénacapavir) pour les échecs thérapeutiques.
\end{itemize}
\end{savaistu}

\begin{exercice}[Entraînement Probatoire — Dépistage et Traitement]
\textbf{Document :} Résultats de sérologie VIH de trois patients (A, B, C) testés par ELISA 4\up{ème} génération puis Immunoblot (IB).
\begin{center}
\begin{tabular}{|c|c|c|c|c|}\hline
\textbf{Patient} & \textbf{ELISA Ag/Ab} & \textbf{Bandes IB détectées} & \textbf{Interprétation IB} & \textbf{Statut final} \\ \hline
A & Positif & gp160, gp120, gp41, p24, p17 & Positif & ? \\ \hline
B & Positif & gp120 seulement & Indéterminé & ? \\ \hline
C & Négatif & Non réalisé & - & ? \\ \hline
\end{tabular}
\end{center}
\textbf{Questions :}
\begin{enumerate}
    \item Compléter le tableau (Statut final) en justifiant pour chaque patient.
    \item Pour le patient B, proposer deux hypothèses expliquant ce profil et la conduite à tenir.
    \item Un patient sous ARV efficace depuis 2 ans a une charge virale indétectable. Son test ELISA reste positif. Expliquer pourquoi.
    \item Citer trois arguments pour encourager un ami à se faire dépister volontairement (Méthode vue en cours).
\end{enumerate}
\end{exercice}

\begin{corrige}[Exercice n°1]
\begin{enumerate}
    \item \textbf{Patient A :} \textbf{Séropositif confirmé.} Bandes spécifiques \emph{env} (gp160/120/41) + \emph{gag} (p24/p17) présentes $\rightarrow$ critère de positivité Immunoblot rempli (au moins 2 bandes \emph{env} ou 1 \emph{env} + 1 \emph{gag}).
    \textbf{Patient B :} \textbf{Statut indéterminé (nécessite contrôle).} Une seule bande (gp120) $\rightarrow$ insuffisant pour la positivité. Peut être : infection récente (séroconversion en cours, anticorps pas encore matures contre p24/gp41), infection par VIH-2 (réactivité croisée partielle sur gp120), ou faux positif ELISA (réaction croisée non spécifique, maladies auto-immunes, grossesse, vaccins récents).
    \textbf{Patient C :} \textbf{Séronégatif (si hors fenêtre sérologique).} ELISA Ag/Ab négatif $\rightarrow$ pas d'antigène p24 ni d'anticorps détectés. \textbf{Attention :} Si prise de risque $< 6$ semaines, refaire test à J+6 semaines (fenêtre sérologique résiduelle possible même en 4\up{ème} gén.).
    \item \textbf{Hypothèses pour B :} (1) \textbf{Primo-infection / Séroconversion :} Les anti-gp120 apparaissent souvent en premier, anti-p24/gp41 plus tard. (2) \textbf{Faux positif ELISA / Réaction croisée :} Anticorps hétérophiles ou dirigés contre d'autres antigènes réagissant faiblement sur la bande gp120 de l'IB.
    \textbf{Conduite à tenir :} \textbf{Contrôle à J+4 à 6 semaines} (nouveau ELISA Ag/Ab + IB) \textbf{ET/OU réalisation d'une PCR-ARN VIH (charge virale)} pour diagnostic direct précoce. Conseil prévention (préservatifs, PrEP si risque récurrent) en attendant.
    \item Les ARV bloquent la réplication virale (charge virale indétectable) mais \textbf{ne suppriment pas les lymphocytes B mémoires / plasmocytes à longue durée de vie} qui ont déjà produit les anticorps anti-VIH. Ces anticorps circulent encore dans le sang pendant des années (demi-vie longue des IgG, survie des plasmocytes). La sérologie reste donc positive à vie. Seule la charge virale (PCR) reflète l'activité virale actuelle.
    \item \textbf{1. Intérêt individuel :} Traitement précoce = vie longue et en bonne santé (CD4 préservés, pas d'infections opportunistes). \textbf{2. Intérêt collectif :} Indétectable = Intransmissible (I=I), on protège ses partenaires et on freine l'épidémie (TasP). \textbf{3. Intérêt familial (PTME) :} Si femme enceinte, dépistage = accès PTME = bébé né sain (taux transmission $< 2\%$).
\end{enumerate}
\end{corrige}

\newpage

\coursec{Activité d'intégration}

\coursec{Leçon ND08 : Lutte contre le VIH/Sida}

\courssub{Objectifs de l'activité}
Cette activité d'intégration vise à mobiliser l'ensemble des savoirs et savoir-faire de la leçon \og Lutte contre le VIH/Sida \fg à travers une mise en situation professionnelle réaliste. Tu devras :
\begin{itemize}
    \item Analyser et interpréter des données biologiques longitudinales (charge virale, lymphocytes T CD4) pour identifier les phases de l'infection non traitée et estimer le passage au stade Sida.
    \item Classer des situations de la vie courante camerounaise selon leur risque de transmission du VIH, en justifiant par le liquide biologique vecteur.
    \item Construire un argumentaire scientifique, éthique et rassurant pour convaincre une patiente enceinte séropositive de l'intérêt vital du traitement antirétroviral (PTME, observance, concept I=I).
\end{itemize}
Cette activité développe les compétences \og Interpréter des données \fg, \og Argumenter \fg et \og Communiquer \fg dans un contexte de santé publique camerounais.

\courssub{Situation}
Tu es stagiaire \og Conseiller en santé communautaire \fg au Centre de Santé Intégré (CSI) de \textbf{New Bell}, à Douala. L'infirmière chef de poste te confie trois dossiers urgents à traiter avant la réunion de concertation médicale de 14 h.

\textbf{Dossier 1 : Suivi biologique de M. Ekobo, 32 ans, non traité.}
M. Ekobo a découvert sa séropositivité il y a 10 ans lors d'un dépistage volontaire. Par peur du stigmate et par méconnaissance, il a refusé toute prise en charge antirétrovirale. Le biologiste du centre te remet le graphique ci-dessous retraçant l'évolution de sa charge virale (copies ARN VIH/mL, échelle logarithmique) et de son taux de lymphocytes T CD4 (cellules/µL) sur 10 ans.

\begin{popfigure}
\begin{tikzpicture}
\begin{semilogyaxis}[
    width=\linewidth,
    height=10cm,
    xlabel={Temps (années après infection)},
    ylabel={Charge virale (copies/mL)},
    ymin=10, ymax=10000000,
    xmin=0, xmax=10.5,
    grid=major,
    grid style={dashed, gray!30},
    legend pos=north west,
    legend cell align=left,
    tick label style={font=\small},
    label style={font=\small},
    xtick={0,1,2,3,4,5,6,7,8,9,10},
    ytick={10,100,1000,10000,100000,1000000,10000000},
    yticklabels={$10^1$,$10^2$,$10^3$,$10^4$,$10^5$,$10^6$,$10^7$},
]
% Charge virale
\addplot[popcurve, very thick, mark=o, mark size=2.5, color=popred] coordinates {
    (0.1, 5000000) (0.25, 8000000) (0.5, 2000000) (1, 50000) (2, 30000) (3, 40000) (4, 50000) (5, 80000) (6, 120000) (7, 200000) (8, 500000) (9, 1500000) (10, 4000000)
};
\addlegendentry{Charge virale (copies/mL)}

% Phase annotations
\node[align=center, font=\tiny, popred] at (axis cs:0.2, 6000000) {Infection\\(aiguë)};
\node[align=center, font=\tiny, poporange] at (axis cs:2, 40000) {Phase\\(asymptomatique)};
\node[align=center, font=\tiny, popgreen] at (axis cs:9, 2000000) {Stade\\Sida};
\draw[dashed, popmuted] (axis cs:8.5, 10) -- (axis cs:8.5, 10000000);
\node[rotate=90, anchor=south, font=\tiny, popmuted] at (axis cs:8.5, 50000) {Seuil Sida (CD4 < 200)};
\end{semilogyaxis}

\begin{axis}[
    width=\linewidth,
    height=10cm,
    xlabel={Temps (années après infection)},
    ylabel={Lymphocytes T CD4 (cellules/$\mu$L)},
    ymin=0, ymax=1200,
    xmin=0, xmax=10.5,
    axis y line*=right,
    axis x line=none,
    legend pos=south west,
    legend cell align=left,
    tick label style={font=\small},
    label style={font=\small},
    xtick={0,1,2,3,4,5,6,7,8,9,10},
    ytick={0,200,400,600,800,1000,1200},
]
% CD4
\addplot[popcurve, very thick, mark=square, mark size=2.5, color=popblue] coordinates {
    (0, 1000) (0.1, 900) (0.25, 500) (0.5, 600) (1, 700) (2, 680) (3, 650) (4, 600) (5, 520) (6, 450) (7, 380) (8, 300) (9, 220) (10, 150)
};
\addlegendentry{Lymphocytes T CD4 ($\mu$L$^{-1}$)}
\end{axis}
\end{tikzpicture}
\legende{Évolution de la charge virale (courbe rouge, échelle log) et des lymphocytes T CD4 (courbe bleue, échelle linéaire) chez M. Ekobo, non traité, sur 10 ans. Le trait vertical pointillé marque le franchissement du seuil de 200 CD4/µL.}
\end{popfigure}

\textbf{Dossier 2 : Éducation pour la santé — Atelier \og Idées reçues et risques réels \fg.}
Dans le cadre d'une causerie éducative au marché de \textbf{Nkololoun}, tu dois classer les 10 situations suivantes, fréquemment évoquées par la population, en \textbf{À risque} ou \textbf{Sans risque} de transmission du VIH. Pour chaque situation \og À risque \fg, tu dois identifier le \textbf{liquide biologique} vecteur du virus.

\begin{enumerate}
    \item Partage d'un rasoir ou d'une tondeuse non stérilisée au salon de coiffure \og Chez Papa \fg.
    \item Piqûre de moustique (anophèle) après avoir piqué une personne séropositive.
    \item Rapport sexuel vaginal non protégé avec un partenaire de statut sérologique inconnu.
    \item Partage du même plat de \textit{ndolé} et du même verre d'eau lors d'un repas familial.
    \item Scarifications traditionnelles (cicatrisation rituelle) réalisées avec le même couteau non stérilisé sur plusieurs adolescents.
    \item Baiser sur la joue ou embrassade amicale.
    \item Allaitement maternel exclusif par une mère séropositive non traitée.
    \item Transfusion sanguine avec une poche de sang non dépistée (contexte d'urgence vitale en zone rurale).
    \item Utilisation des toilettes publiques (siège, chasse d'eau) après une personne séropositive.
    \item Accouchement par voie basse sans prophylaxie antirétrovirale (PTME) pour une mère séropositive.
\end{enumerate}

\textbf{Dossier 3 : Conseil pour l'observance — Cas de Mme Ngo Bitang.}
Mme Ngo Bitang, 24 ans, enceinte de 6 mois (26 SA), vient d'apprendre sa séropositivité lors de la consultation prénatale (CPN) au CSI. Elle est effondrée et refuse de commencer le traitement antirétroviral (ARV) à vie (Option B+), craignant les effets secondaires, le regard de son mari et de sa belle-famille, et pensant que \og de toute façon, l'enfant sera contaminé \fg. Elle préfère \og prier et prendre des plantes \fg. L'infirmière te demande de rédiger un \textbf{texte argumentatif court (8 à 10 lignes maximum)} à lui lire, pour la convaincre de démarrer le traitement immédiatement et de l'observer strictement.

\courssub{Consignes}
\textbf{Partie A — Analyse du dossier clinique (M. Ekobo)}
\begin{enumerate}
    \item Sur le graphique, identifie et nomme les \textbf{trois phases successives} de l'infection à VIH non traitée. Indique pour chacune la plage temporelle approximative (en années).
    \item Décris l'évolution de la charge virale et des lymphocytes T CD4 lors de la \textbf{phase aiguë (infection primaire)}. Explique le mécanisme immunologique responsable de la chute brutale des CD4 puis de leur remontée partielle.
    \item Caractérise la \textbf{phase asymptomatique} (ou phase de latence clinique) sur le plan virologique (charge virale) et immunologique (CD4). Pourquoi le terme \og latence \fg est-il impropre sur le plan virologique ?
    \item À partir de quelle année estimates-tu que M. Ekobo entre dans le \textbf{stade Sida déclaré} ? Justifie ta réponse en te basant sur \textbf{deux critères biologiques} visibles sur le graphique.
    \item Quelle aurait été l'évolution attendue des deux courbes si M. Ekobo avait débuté un traitement ARV efficace dès l'année 1 (charge virale indétectable $< 50$ copies/mL, reconstitution des CD4) ? Dessine schématiquement ces deux courbes hypothétiques sur une copie du graphique.
\end{enumerate}

\textbf{Partie B — Classification des risques (Atelier marché Nkololoun)}
\begin{enumerate}
    \setcounter{enumi}{5}
    \item Complète le tableau ci-dessous en classant chaque situation (1 à 10) et en justifiant pour les situations à risque.
\end{enumerate}

\begin{center}
\begin{tabularx}{\linewidth}{|c|X|c|X|}
\hline
\rowcolor{popblueL}
\textbf{N°} & \textbf{Situation} & \textbf{Classement} & \textbf{Liquide biologique vecteur (si à risque)} \\
\hline
1 & Partage rasoir/tondeuse non stérilisé & & \\
\hline
2 & Piqûre de moustique & & \\
\hline
3 & Rapport sexuel non protégé & & \\
\hline
4 & Repas partagé (ndolé, verre) & & \\
\hline
5 & Scarifications couteau commun & & \\
\hline
6 & Baiser sur la joue & & \\
\hline
7 & Allaitement maternel (mère non traitée) & & \\
\hline
8 & Transfusion sang non dépisté & & \\
\hline
9 & Toilettes publiques & & \\
\hline
10 & Accouchement sans PTME & & \\
\hline
\end{tabularx}
\end{center}

\textbf{Partie C — Argumentaire pour l'observance (Mme Ngo Bitang)}
\begin{enumerate}
    \setcounter{enumi}{6}
    \item Rédige le texte argumentatif (8 à 10 lignes) destiné à Mme Ngo Bitang. Ton texte doit impérativement mobiliser les notions suivantes : \textbf{PTME (Prévention de la Transmission Mère-Enfant)}, \textbf{Observance thérapeutique}, \textbf{Charge virale indétectable = Intransmissibilité (I=I)}, \textbf{Risque de transmission verticale sans traitement}, \textbf{Bénéfice pour la santé maternelle}.
\end{enumerate}

\begin{aretenir}[Ressources à mobiliser]
\begin{itemize}
    \item \textbf{Structure du VIH :} Rétrovirus à ARN, enveloppe (gp120/gp41), capside, transcriptase inverse, intégrase, protéase. Cible : lymphocytes T CD4+ (récepteur CD4 + co-récepteurs CCR5/CXCR4).
    \item \textbf{Cycle d'infection :} Fixation $\rightarrow$ Fusion $\rightarrow$ Transcription inverse (ARN $\rightarrow$ ADN) $\rightarrow$ Intégration (provirus) $\rightarrow$ Transcription/Traduction $\rightarrow$ Assemblage $\rightarrow$ Bourgeonnement (maturation par protéase).
    \item \textbf{Critères biologiques du Sida :} CD4 $< 200$ cellules/µL \textbf{OU} survenue d'une infection opportuniste définissante (ex: pneumocystose, toxoplasmose cérébrale, tuberculose disséminée).
    \item \textbf{Transmission :} Quatre liquides biologiques principaux à forte concentration virale : \textbf{Sang}, \textbf{Sperme/Liquide pré-éjaculatoire}, \textbf{Sécrétions vaginales}, \textbf{Lait maternel}. \emph{Salive, sueur, larmes, urine, selles : non contaminants (concentration virale négligeable, inhibiteurs).}
    \item \textbf{PTME (Option B+ au Cameroun) :} Traitement ARV à vie pour toute femme enceinte/allaitante séropositive, débuté le plus tôt possible (idéalement $\le$ 14 SA). Objectif : Charge virale indétectable ($< 50$ copies/mL) à l'accouchement et pendant l'allaitement.
    \item \textbf{Concept I=I (Indétectable = Intransmissible) :} Personne sous ARV efficace, charge virale indétectable durablement ($> 6$ mois) $\rightarrow$ Risque nul de transmission sexuelle (études PARTNER, HPTN 052).
    \item \textbf{Observance :} Prise quotidienne, horaire fixe, sans omission. Résistance virale si observance $< 95\%$.
\end{itemize}
\end{aretenir}

\begin{exoresolu}{Activité d'intégration : Lutte contre le VIH/Sida}
\textbf{Énoncé :} Voir les trois dossiers et les consignes des Parties A, B, C ci-dessus (analyse graphique M. Ekobo, classification risques Nkololoun, argumentaire Mme Ngo Bitang).
\tcblower
\textbf{PARTIE A — Analyse du dossier clinique (M. Ekobo)}

\textbf{1. Identification des trois phases de l'infection non traitée}
L'analyse du graphique permet de distinguer trois phases successives :
\begin{itemize}
    \item \textbf{Phase 1 : Infection primaire (aiguë) / Séroconversion} — \textbf{De 0 à $\approx$ 0,5--1 an.} Pic de charge virale ($> 10^6$ copies/mL), chute brutale des CD4 (jusqu'à $\approx 500$ cellules/µL), puis remontée partielle des CD4 et baisse de la charge virale (séroconversion, apparition des anticorps anti-VIH).
    \item \textbf{Phase 2 : Phase asymptomatique (latence clinique)} — \textbf{De $\approx$ 1 an à $\approx$ 8--9 ans.} Charge virale stabilisée à un \og point de consigne \fg (set point) bas ($10^4$--$10^5$ copies/mL). Déclin lent, progressif et continu des lymphocytes T CD4 (perte $\approx 50$--$80$ cellules/µL/an). Le patient ne présente pas de symptômes cliniques majeurs.
    \item \textbf{Phase 3 : Stade Sida (immunodépression sévère)} — \textbf{À partir de $\approx$ 8,5--9 ans.} Rupture de l'équilibre : charge virale repart à la hausse ($> 10^5$--$10^6$ copies/mL), effondrement des CD4 sous le seuil critique de $200$ cellules/µL. Apparition d'infections opportunistes.
\end{itemize}

\textbf{2. Mécanismes de la phase aiguë (0--6 mois)}
\begin{itemize}
    \item \textbf{Charge virale :} Explosion exponentielle. Le virus infecte massivement les lymphocytes T CD4+ activés (surtout au niveau du tissu lymphoïde intestinal -- GALT), produisant jusqu'à $10^7$--$10^8$ copies/mL. L'absence de réponse immunitaire adaptative spécifique permet cette réplication incontrôlée.
    \item \textbf{CD4 :} Chute brutale (lyse virale directe, mort cellulaire par syncytia, élimination par CD8 cytotoxiques, apoptose des cellules non infectées \og bystander \fg). Minimum atteint vers 2--4 semaines ($\approx 500$ cellules/µL sur le graphique).
    \item \textbf{Remontée partielle des CD4 :} Mise en place de la \textbf{réponse immune adaptative} : expansion clonale de lymphocytes T CD8+ cytotoxiques spécifiques (reconnaissance peptides VIH/HLA-I) et production d'anticorps neutralisants (séroconversion). Cette réponse contrôle partiellement la virémie (baisse de 2--3 log), permettant une reconstitution partielle du stock de CD4 (niveau \og set point \fg $\approx 700$ cellules/µL), mais n'élimine pas le réservoir viral.
\end{itemize}

\textbf{3. Caractérisation de la phase asymptomatique (An 1 -- An 8)}
\begin{itemize}
    \item \textbf{Virologie :} Charge virale stable au \og point de consigne \fg (ici $\approx 3 \times 10^4$ à $10^5$ copies/mL). \textbf{Le terme \og latence \fg est impropre :} le virus réplique \textbf{activement et en permanence} ($10^9$--$10^{10}$ virions produits/jour), avec un turnover viral très rapide (demi-vie du virion $\approx 6$ h, demi-vie du lymphocyte infecté $\approx 1$--2 jours). L'équilibre dynamique est maintenu par la réponse immune.
    \item \textbf{Immunologie :} Érosion lente et inexorable du pool de lymphocytes T CD4+ naïfs et à mémoire centrale (thymus insuffisant, activation immune chronique, épuisement des cellules souches hématopoïétiques, fibrose des organes lymphoïdes). Perte moyenne $\approx 60$ cellules/µL/an ici (de 700 à 220 en 8 ans).
\end{itemize}

\textbf{4. Estimation du passage au stade Sida}
\textbf{Vers l'année 8,5 -- 9 (entre la 8ème et la 9ème année).}
\textbf{Justification par deux critères biologiques concomitants :}
\begin{enumerate}
    \item \textbf{Critère immunologique majeur :} Le taux de lymphocytes T CD4+ franchit durablement le seuil de \textbf{200 cellules/µL} (valeur lue : 220 à l'an 9, 150 à l'an 10). C'est le critère définitionnel principal du Sida (stade C3 CDC / stade 4 OMS) en l'absence d'infection opportuniste déclarée.
    \item \textbf{Critère virologique de rupture :} La charge virale, stable depuis 8 ans, \textbf{remonte significativement} (de $1,2 \times 10^5$ à l'an 7 à $4 \times 10^6$ à l'an 10). Cette \og rupture virologique \fg traduit la perte de contrôle immunitaire (défaillance des CD8+, épuisement des cellules T, diversité virale échappant à la réponse immune).
\end{enumerate}
\emph{Note :} Cliniquement, ce seuil de 200 CD4 expose M. Ekobo à un risque élevé d'infections opportunistes (pneumocystose, toxoplasmose, tuberculose, candidose œsophagienne...).

\textbf{5. Évolution hypothétique sous ARV efficace (début année 1)}
\begin{center}
\begin{tikzpicture}
\begin{semilogyaxis}[
    width=0.9\linewidth, height=7cm,
    xlabel={Temps (années)}, ylabel={Charge virale (copies/mL)},
    xmin=0, xmax=10, ymin=10, ymax=10000000,
    grid=major, grid style={dashed, gray!30},
    legend pos=north west, tick label style={font=\small},
    xtick={0,1,2,3,4,5,6,7,8,9,10},
    ytick={10,100,1000,10000,100000,1000000,10000000},
    yticklabels={$10^1$,$10^2$,$10^3$,$10^4$,$10^5$,$10^6$,$10^7$},
]
% VL sans traitement (rappel pointillés)
\addplot[popred, dashed, thick] coordinates {(0.1,5000000) (0.25,8000000) (0.5,2000000) (1,50000) (2,30000) (3,40000) (4,50000) (5,80000) (6,120000) (7,200000) (8,500000) (9,1500000) (10,4000000)};
\addlegendentry{VL sans traitement (rappel)}

% VL AVEC ARV (début an 1)
\addplot[popgreen, very thick, mark=o] coordinates {
    (0.1, 5000000) (0.25, 8000000) (0.5, 2000000) (1, 50000) 
    (1.1, 5000) (1.5, 200) (2, 40) (3, 20) (4, 20) (5, 20) (6, 20) (7, 20) (8, 20) (9, 20) (10, 20)
};
\addlegendentry{VL sous ARV (cible $<50$)}
\draw[dashed, popgreen] (axis cs:1, 50) -- (axis cs:10, 50);
\node[popgreen, font=\tiny, anchor=west] at (axis cs:1.2, 60) {Seuil indétectable ($<50$)};
\end{semilogyaxis}

\begin{axis}[
    width=0.9\linewidth, height=7cm,
    xlabel={Temps (années)}, ylabel={Lymphocytes T CD4 (cellules/$\mu$L)},
    xmin=0, xmax=10, ymin=0, ymax=1200,
    axis y line*=right, axis x line=none,
    legend pos=south west, tick label style={font=\small},
    xtick={0,1,2,3,4,5,6,7,8,9,10},
    ytick={0,200,400,600,800,1000,1200},
]
% CD4 sans traitement (rappel)
\addplot[popblue, dashed, thick] coordinates {(0,1000) (0.1,900) (0.25,500) (0.5,600) (1,700) (2,680) (3,650) (4,600) (5,520) (6,450) (7,380) (8,300) (9,220) (10,150)};
\addlegendentry{CD4 sans traitement (rappel)}

% CD4 AVEC ARV
\addplot[poppurple, very thick, mark=square] coordinates {
    (0, 1000) (0.1, 900) (0.25, 500) (0.5, 600) (1, 700)
    (1.5, 750) (2, 800) (3, 850) (4, 900) (5, 950) (6, 980) (7, 1000) (8, 1020) (9, 1030) (10, 1040)
};
\addlegendentry{CD4 sous ARV (reconstitution)}
\draw[dashed, poppurple] (axis cs:0, 200) -- (axis cs:10, 200);
\node[poppurple, font=\tiny, anchor=west] at (axis cs:1, 220) {Seuil Sida (200) -- loin au-dessus};
\end{axis}
\end{tikzpicture}
\end{center}
\legende{Évolution théorique sous ARV efficace débuté à l'an 1. Charge virale : chute rapide (2--3 log en 1 mois) puis indétectable ($<50$ copies/mL) durablement. CD4 : reconstitution progressive (gain $\approx 50$--$100$ cellules/an les 2 premières années) jusqu'à normalisation. Le patient ne développe pas le Sida.}

\textbf{PARTIE B — Classification des risques (Atelier Nkololoun)}

\begin{center}
\begin{tabularx}{\linewidth}{|c|X|c|X|}
\hline
\rowcolor{popblueL}
\textbf{N°} & \textbf{Situation} & \textbf{Classement} & \textbf{Liquide biologique vecteur (si à risque)} \\
\hline
1 & Partage rasoir/tondeuse non stérilisé & \textbf{À RISQUE} & \textbf{Sang} (micro-coupures cutanées, lésions du cuir chevelu). Concentration virale élevée dans le sang. \\
\hline
2 & Piqûre de moustique & \textbf{SANS RISQUE} & Le VIH ne se réplique pas dans le moustique (pas de récepteur CD4). La proboscis n'injecte pas le sang du repas précédent (salive seulement). Dose virale insuffisante. \\
\hline
3 & Rapport sexuel vaginal non protégé & \textbf{À RISQUE} & \textbf{Sperme / Liquide pré-éjaculatoire} (homme $\rightarrow$ femme) ET \textbf{Sécrétions vaginales / Sang menstruel} (femme $\rightarrow$ homme). Muqueuses perméables, forte charge virale. \\
\hline
4 & Repas partagé (ndolé, verre) & \textbf{SANS RISQUE} & Salive : concentration virale très faible ($< 1$ copie/mL), présence d'inhibiteurs (SLPI, anticorps, pH). Pas de porte d'entrée (muqueuse buccale intacte). Pas de transmission documentée par partage couverts/verre. \\
\hline
5 & Scarifications couteau commun & \textbf{À RISQUE} & \textbf{Sang}. Pratique à haut risque : effraction cutanée profonde, instrument contaminé par sang frais, entrée directe dans la circulation. Fréquent en milieu traditionnel/rural. \\
\hline
6 & Baiser sur la joue & \textbf{SANS RISQUE} & Contact cutané intact. Pas d'échange de liquides biologiques contaminants. Salive non contaminante. \\
\hline
7 & Allaitement maternel (mère non traitée) & \textbf{À RISQUE} & \textbf{Lait maternel}. Contient VIH libre et cellules infectées (lymphocytes, macrophages). Risque cumulatif $\approx 15$--$20\%$ sur 18--24 mois sans ARV. Porte d'entrée : muqueuse digestive du nourrisson (immature). \\
\hline
8 & Transfusion sang non dépisté & \textbf{À RISQUE} & \textbf{Sang total / Produits sanguins}. Concentration virale maximale. Entrée directe intraveineuse. Risque de transmission $\approx 95\%$. (Au Cameroun, le sang est systématiquement dépisté : risque résiduel fenêtre sérologique). \\
\hline
9 & Toilettes publiques & \textbf{SANS RISQUE} & Urines/selles : concentration virale nulle ou négligeable. Virus fragile dans l'environnement (séchage, pH, température). Pas de porte d'entrée par voie cutanée intacte. \\
\hline
10 & Accouchement sans PTME & \textbf{À RISQUE} & \textbf{Sang maternel / Liquides génitaux / Liquide amniotique}. Contact prolongé muqueuses/cutanées lésées du nouveau-né lors du passage. Risque $\approx 15$--$30\%$ sans intervention. \\
\hline
\end{tabularx}
\end{center}

\textbf{PARTIE C — Argumentaire pour Mme Ngo Bitang (Texte lu par le conseiller)}

\begin{quote}
\textit{« Ma chère sœur, écoute-moi bien. Ton bébé peut naître séronégatif, c'est une certitude scientifique aujourd'hui. Si tu commences \textbf{immédiatement} ton traitement ARV (Option B+) et que tu le prends \textbf{tous les jours sans oublier} (c'est ça, l'\textbf{observance}), ta charge virale va devenir \textbf{indétectable} dans ton sang et ton lait.}

\textit{Grâce au principe \textbf{I=I (Indétectable = Intransmissible)}, tu ne transmettras \textbf{PAS} le virus à ton enfant, ni pendant la grossesse, ni à l'accouchement, ni en l'allaitant. C'est le cœur de la \textbf{PTME} : te soigner, c'est le protéger.}

\textit{Sans traitement, le risque pour lui est de 1 sur 3 à 1 sur 2. Avec le traitement, il est \textbf{zéro}. Les médicaments sont gratuits, efficaces, et les effets secondaires se gèrent. Ton mari et ta belle-famille n'ont pas à savoir ton statut si tu ne le souhaites pas : le secret médical est ton droit. Choisis la vie pour toi et pour ton bébé. L'infirmière et moi sommes là pour t'accompagner, chaque jour. »}
\end{quote}

\textbf{Analyse de l'argumentaire (vérification des consignes) :}
\begin{itemize}
    \item \textbf{Longueur :} 9 lignes (format A4, police standard) — \textbf{Respecté}.
    \item \textbf{PTME :} Cité explicitement (\og Option B+ \fg, \og cœur de la PTME \fg).
    \item \textbf{Observance :} Cité et définie (\og tous les jours sans oublier \fg).
    \item \textbf{I=I :} Cité et expliqué (\og charge virale indétectable... ne transmettras PAS \fg).
    \item \textbf{Risque sans traitement :} Chiffré (\og 1 sur 3 à 1 sur 2 \fg).
    \item \textbf{Bénéfice maternel :} Implicite (\og te soigner \fg, \og choisir la vie pour toi \fg).
    \item \textbf{Ton :} Empathique (\og Ma chère sœur \fg), rassurant, autorisant (droit au secret), mobilisateur (\og Choisis la vie \fg).
\end{itemize}
\end{exoresolu}

\begin{aretenir}[Bilan et points clés validés]
Cette activité d'intégration a permis de mettre en évidence, par la pratique, l'articulation indispensable entre \textbf{biologie fondamentale} (cinétique virale/immune, structure virale, cycle réplicatif), \textbf{épidémiologie de terrain} (identification des liquides contaminants et des pratiques à risque réelles au Cameroun : scarifications, rasoirs, accouchement, allaitement) et \textbf{compétences psychosociales} (communication thérapeutique, lutte contre la stigmatisation, promotion de l'observance).

Points clés validés :
\begin{itemize}
    \item La \textbf{dynamique biphasique} de l'infection non traitée : réplication virale massive continue masquée par un contrôle immunitaire partiel et transitoire, aboutissant inéluctablement à l'effondrement immunitaire (CD4 $< 200$) sans traitement.
    \item La \textbf{spécificité des liquides biologiques} vecteurs (Sang, Sperme, Sécr. vaginales, Lait) vs les liquides non vecteurs (Salive, Sueur, Larmes, Urines, Selles) — base scientifique de la prévention combinée (préservatifs, matériel stérile, PTME, PrEP/PEP).
    \item La \textbf{réversibilité du pronostic} par les ARV : l'atteinte de l'indétectabilité virale restaure l'espérance de vie normale et annule le risque de transmission (I=I), transformant l'infection en pathologie chronique contrôlée.
    \item L'\textbf{observance stricte} ($> 95\%$) comme condition \textit{sine qua non} de l'efficacité durable et de la prévention de l'émergence de résistances (barrière génétique).
\end{itemize}
En tant que futurs professionnels de santé ou citoyens éclairés, vous êtes désormais capables d'\textbf{interpréter un suivi biologique}, de \textbf{déconstruire les idées reçues} locales et de \textbf{porter un message d'espoir fondé sur la preuve scientifique} auprès des populations vulnérables. C'est toute la finalité de l'Éducation à la Santé dans notre programme.
\end{aretenir}

\newpage

\coursec{Méthodes et exercices résolus}

\coursec{Lutte contre le VIH/Sida}

\courssub{Le VIH : un virus particulier}

\begin{definition}[Virus de l'Immunodéficience Humaine (VIH)]
Le \textbf{VIH} est un \textbf{rétrovirus} (famille des \emph{Retroviridae}, genre \emph{Lentivirus}) à ARN simple brin, enveloppé. Deux types existent : \textbf{VIH-1} (pandémique, virulent, groupe M dominant au Cameroun) et \textbf{VIH-2} (endémique Afrique de l'Ouest, moins virulent, résistance naturelle aux INNRTI). Le VIH cible préférentiellement les \textbf{lymphocytes T CD4$^+$} (coordinateurs de l'immunité adaptative), les macrophages et les cellules dendritiques.
\end{definition}

\begin{popfigure}
\begin{tikzpicture}[scale=0.9, every node/.style={font=\small}]
% Envelope
\draw[popblue, thick, fill=popblue!10] (0,0) ellipse (2.5cm and 1.5cm);
% Glycoproteins
\foreach \angle/\r in {30/2.5, 60/2.5, 90/2.3, 120/2.5, 150/2.5, 210/2.5, 240/2.5, 270/2.3, 300/2.5, 330/2.5} {
    \draw[poporange, thick] (\angle:\r) -- (\angle:\r+0.4) node[above right, poporange] {\tiny gp120};
    \draw[popdark, thick] (\angle:\r+0.1) -- (\angle:\r+0.3) node[below left, popdark] {\tiny gp41};
}
% Matrix (p17)
\draw[popgreen, thick, fill=popgreen!10] (0,0) ellipse (2.2cm and 1.2cm);
% Capsid (p24) - conical
\draw[poppurple, thick, fill=poppurple!15] (-1.2,0.2) -- (0,0.8) -- (1.2,0.2) -- (0,-0.8) -- cycle;
% RNA + Enzymes inside capsid
\node[poppink, align=center] at (0,0) {\textbf{2 ARN génomiques}\\(+)\\\textbf{Transcriptase Inverse (RT)}\\\textbf{Intégrase (IN)}\\\textbf{Protéase (PR)}};
% Legends
\draw[->, popblue] (2.8,1.2) -- (3.5,1.5) node[right] {Enveloppe lipidique (origine cellulaire)};
\draw[->, poporange] (2.8,0.5) -- (3.5,0.8) node[right] {Glycoprotéines \textbf{gp120/gp41} (clé entrée)};
\draw[->, popgreen] (-3.5,-0.5) -- (-4.0,-0.8) node[left] {Matrice \textbf{p17} (stabilité)};
\draw[->, poppurple] (-3.5,0.5) -- (-4.0,0.8) node[left] {Capside \textbf{p24} (protection ARN)};
\draw[->, poppink] (3.0,-0.8) -- (3.5,-1.1) node[right] {Enzymes virales : RT, IN, PR};
\end{tikzpicture}
\legende{\textbf{Structure du virion VIH-1 (vue schématique en coupe).} L'enveloppe porte les gp120 (liaison CD4/co-récepteur) et gp41 (fusion). La capside p24 protège les 2 ARN identiques et les 3 enzymes essentielles : Transcriptase Inverse (RT), Intégrase (IN), Protéase (PR).}
\end{popfigure}

\begin{propriete}[Génome et enzymes cibles thérapeutiques]
Le génome (≈ 9,7 kb) code pour les gènes structuraux (\emph{gag, pol, env}) et régulateurs/accessoires (\emph{tat, rev, nef, vif, vpr, vpu}).
\begin{itemize}
    \item \textbf{Gag} $\rightarrow$ Pr55 $\xrightarrow{\text{PR}}$ p24 (capside), p17 (matrice), p7 (nucléocapside).
    \item \textbf{Pol} $\rightarrow$ Pr160 $\xrightarrow{\text{PR}}$ \textbf{RT} (ARN$\rightarrow$ADN), \textbf{IN} (intégration), \textbf{PR} (maturation).
    \item \textbf{Env} $\rightarrow$ gp160 $\xrightarrow{\text{clivage cellulaire}}$ \textbf{gp120} (liaison) + \textbf{gp41} (fusion).
\end{itemize}
\cle{Cibles ARV :} \textbf{RT} (INRTI : TDF, 3TC, AZT ; INNRTI : EFV, NVP, DOR), \textbf{IN} (II : DTG, RAL, BIC), \textbf{PR} (IP : ATV/r, LPV/r, DRV/r), \textbf{gp120} (Inhibiteurs fixation : Fostemsavir), \textbf{gp41} (Inhibiteurs fusion : Enfuvirtide), \textbf{CCR5} (Antagonistes : Maraviroc).
\end{propriete}

\courssub{Le cycle d'infection du lymphocyte CD4}

\begin{experience}[Cycle réplicatif du VIH : étapes et cibles thérapeutiques]
\textbf{Observation :} Le VIH ne se réplique que dans les cellules exprimant le récepteur CD4 et un co-récepteur (CCR5 ou CXCR4).
\textbf{Étapes (cible ARV en rouge) :}
\begin{enumerate}
    \item \textbf{Fixation (Attachement) :} \textbf{gp120} se lie au \textbf{CD4} $\rightarrow$ changement conformationnel $\rightarrow$ liaison au \textbf{co-récepteur} (CCR5 : souches R5 « précoces » ; CXCR4 : souches X4 « tardives », plus cytopathiques). \textcolor{popred}{\textbf{Cible :} Maraviroc (CCR5), Fostemsavir (gp120).}
    \item \textbf{Fusion :} \textbf{gp41} médie la fusion enveloppe/membrane cellulaire $\rightarrow$ entrée capside dans le cytoplasme. \textcolor{popred}{\textbf{Cible :} Enfuvirtide (T-20, gp41).}
    \item \textbf{Désencapsidation :} Libération de l'ARN viral et des enzymes (RT, IN, PR).
    \item \textbf{Transcription inverse :} \textbf{RT} synthétise l'ADN double brin (ADN proviral) à partir de l'ARN viral. \emph{Erreur fréquente : RT n'a pas d'activité proofreading $\rightarrow$ taux mutation élevé ($3 \times 10^{-5}$/base/cycle).} \textcolor{popred}{\textbf{Cible :} INRTI (TDF, 3TC, AZT - analogues nucléosidiques) ; INNRTI (EFV, NVP, DOR - site allostérique).}
    \item \textbf{Transport nucléaire \& Intégration :} Complexe de pré-intégration (PIC) traverse le pore nucléaire (actif en phase mitotique ou via importines). \textbf{IN} insère l'ADN proviral dans l'ADN hôte (préférence gènes actifs). \emph{Irréversible = réservoir latent.} \textcolor{popred}{\textbf{Cible :} II (DTG, RAL, BIC - blocage transfert brin).}
    \item \textbf{Transcription / Translation :} Promoteur viral (LTR 5') $\xrightarrow{\text{ARN Pol II hôte + Tat}}$ ARN messagers (épissés ou non) $\rightarrow$ Traduction protéines virales (Gag, Pol, Env, accessoires). \emph{Rev exporte les ARN non/peu épissés.}
    \item \textbf{Assemblage :} Pr55 Gag + ARN génomique + enzymes $\rightarrow$ bourgeonnement à la membrane plasmique (domaines lipidiques riches en cholestérol).
    \item \textbf{Maturation :} \textbf{PR} clive Pr55 et Pr160 $\rightarrow$ protéines fonctionnelles (p24, RT, IN, PR, gp120/gp41). Virion \textbf{infectieux}. \textcolor{popred}{\textbf{Cible :} IP (ATV/r, LPV/r, DRV/r - inhibent PR).}
\end{enumerate}
\legende{Chaque étape est une cible potentielle. La trithérapie standard (ex: TDF+3TC+DTG) combine 2 INRTI + 1 II (ou 1 IP/r, ou 1 INNRTI) pour bloquer la réplication à plusieurs stades et prévenir la résistance.}
\end{experience}

\courssub{Transmission et prévention}

\begin{definition}[Fluides contaminants et portes d'entrée]
Seuls 4 fluides biologiques contiennent une concentration virale suffisante pour transmettre l'infection : \textbf{Sang}, \textbf{Sperme}, \textbf{Sécrétions vaginales}, \textbf{Lait maternel}.
\emph{Non contaminants (sauf contamination sanguine visible) :} Salive, sueur, larmes, urine, selles, vomissures.
\textbf{Portes d'entrée :} Muqueuses (génitales, rectales, buccales, nasales, conjonctivales), Peau lésée (plaie, ulcération, injection), Voie intraveineuse directe, Barrière placentaire (in utero), Voie digestive du nourrisson (allaitement).
\end{definition}

\begin{propriete}[Modes de transmission (Épidémiologie Cameroun - EDS 2018 / ONUSIDA 2023)]
\begin{itemize}
    \item \textbf{Sexuel (≈ 80-85\% des cas) :} Rapports vaginaux/anaux/oraux non protégés. Risque par acte : Réceptif anal ($>$1\%) $>$ Insertif anal / Réceptif vaginal (0,1-1\%) $>$ Insertif vaginal $>$ Oral (faible, réceptif avec éjaculation). Facteurs augmentant le risque : IST ulcéreuses (syphilis, herpès), charge virale élevée, phase aiguë, non-circoncision (homme HSH).
    \item \textbf{Sanguin :} Partage seringues/aiguilles (UDI), matériel scarification/tatouage/piercing non stérilisé, accident exposition sang (AES : piqûre aiguille $\approx$ 0,3\%), transfusion (rare, dépistage systématique).
    \item \textbf{Mère-Enfant (Verticale) :} \emph{In utero} (transplacentaire, rare $<$ 5\% si pas d'ARV), \emph{Per partum} (accouchement, contact sang/sécrétions, principal sans ARV), \emph{Post partum} (allaitement, risque cumulé 5-20\% sans ARV).
\end{itemize}
\cle{Fausses voies :} Contacts quotidiens (serrer main, baiser social, partager repas/toilettes/piscine), piqûres moustiques (pas de réplication dans le moustique, pas d'inoculation sanguine).
\end{propriete}

\begin{methode}[Analyser une situation à risque et identifier le mode de transmission du VIH]
\textbf{Objectif :} Expliquer les modes de transmission du VIH à partir d'un cas concret.
\begin{enumerate}
    \item \textbf{Identifier le fluide biologique mis en jeu :} Sang, sperme, sécrétions vaginales, lait maternel (fluides contaminants). Salive, sueur, larmes, urine, selles (non contaminants sauf contamination sanguine visible).
    \item \textbf{Repérer la porte d'entrée :} Muqueuses (génitales, rectales, buccales), peau lésée, voie intraveineuse, voie transplacentaire/pernatale.
    \item \textbf{Qualifier le mode de transmission :}
    \begin{itemize}
        \item \textbf{Sexuel :} Rapports non protégés (vaginaux, anaux, oraux \emph{receptif} avec éjaculation).
        \item \textbf{Sanguin :} Partage de seringues/aiguilles, matériel de scarification/tatouage non stérilisé, transfusion non dépistée (rare aujourd'hui), accident d'exposition au sang (AES).
        \item \textbf{Mère-enfant (verticale) :} \emph{In utero} (transplacentaire), \emph{per partum} (au passage, contact sang/sécrétions), \emph{post partum} (allaitement).
    \end{itemize}
    \item \textbf{Écarter les fausses voies :} Contacts quotidiens (serrer la main, embrasser, partager repas/toilettes, piqûres de moustiques).
    \item \textbf{Rédiger la réponse :} « Le mode de transmission est \textbf{...} car le fluide \textbf{...} contenant le virus a franchi la barrière \textbf{...} lors de \textbf{...} ».
\end{enumerate}
\textbf{Reflexe Probatoire :} Ne jamais écrire « transmission par la salive » sans préciser « contaminée par le sang ». Préciser toujours le fluide \textbf{et} la porte d'entrée.
\end{methode}

\begin{propriete}[La Prévention Combinée (OMS / ONUSIDA / MINESEC Cameroun)]
Aucune méthode seule n'offre 100\% de protection (sauf abstinence totale). La stratégie efficace associe trois piliers :
\begin{center}
\begin{tabularx}{\linewidth}{|l|X|X|}\hline
\textbf{Pilier} & \textbf{Interventions clés} & \textbf{Cible / Commentaire} \\ \hline
\textbf{Comportementale} & Abstinence, fidélité mutuelle dépistée, réduction partenaires, \textbf{Préservatif masculin/féminin (usage correct + systématique)} & Base pour tous. Protège VIH + IST + Grossesse. Norme ISO 4074 / CE. \\ \hline
\textbf{Biomédicale} & \textbf{PrEP} (TDF/3TC ou TDF/FTC quotidiens, séronégatifs risque élevé). \textbf{PEP} (ARV urgence $<$ 48h, 28 j). \textbf{TasP / I=I} (ARV partenaire $+$ $\rightarrow$ CV indétectable). \textbf{PMTCT} (Option B+ : ARV mère vie + AZT/3TC enfant 6 sem + accouchement sécurisé + conseil alimentation). \textbf{Circoncision médicale volontaire} (réduit acquisition homme 60\%). & PrEP : observance stricte (7j rectal, 20j vaginal). PEP : délai critique. TasP : CV $<$ 50 copies/mL durable. PMTCT : taux transmission $<$ 1\% au Cameroun. \\ \hline
\textbf{Structurelle} & Éducation sexuelle complète (ESC), lutte stigmatisation/discrimination, dépistage élargi (CDV, auto-test, PICT), programmes réduction risques (RdR : échange seringues, substitution opiacés), accès universel soins (gratuité ARV/PrEP/PEP au Cameroun). & Crée l'environnement favorable. Clé pour populations clés (HSH, TS, UDI, prisonniers, jeunes filles). \\ \hline
\end{tabularx}
\end{center}
\textbf{Formule clé :} \textbf{Prévention Combinée = Comportementale + Biomédicale + Structurelle}.
\end{propriete}

\courssub{Évolution de l'infection à VIH}

\begin{definition}[Histoire naturelle de l'infection VIH non traitée]
L'infection évolue en 3 phases cliniques successives, marquées par la dynamique antagoniste \textbf{Charge Virale (CV) vs Lymphocytes T CD4}.
\begin{center}
\begin{tabular}{|c|c|c|c|}\hline
\textbf{Phase} & \textbf{Délai / Durée} & \textbf{Charge Virale (copies/mL)} & \textbf{Lymphocytes T CD4 (cell/mm$^3$)} \\ \hline
\textbf{1. Primo-infection / Séroconversion} & 2--6 sem. post-contamination & \textbf{Pic extrême} ($> 10^6$, souvent $10^7$) & \textbf{Chute brutale} (souvent $< 500$), puis remontée partielle \\ \hline
\textbf{2. Phase asymptomatique (Latence clinique)} & \textbf{Moyenne 8--10 ans} (extrêmes 2--15+ ans) & \textbf{Plateau : « Point de consigne viral »} ($10^4$--$10^5$) & \textbf{Déclin lent, linéaire} ($\approx$ 50--80 / an) \\ \hline
\textbf{3. Phase Sida (Symptomatique)} & Sans ARV : décès en 1--2 ans & \textbf{Remontée exponentielle} ($> 10^5$--$10^6$) & \textbf{Effondrement} ($< 200$) $\rightarrow$ Infections Opportunistes (IO) / Cancers \\ \hline
\end{tabular}
\end{center}
\cle{Séroconversion :} Apparition des anticorps anti-VIH (détectables tests rapides 3--12 sem., tests Ag/Ab 2--6 sem.). Syndrome rétroviral aigu (fièvre, adénopathies, éruption, pharyngite) dans 50-90\% cas, souvent méconnu.
\end{definition}

\begin{popfigure}
\begin{tikzpicture}
\begin{axis}[
    width=\linewidth, height=9cm,
    xlabel={Temps (années) sans traitement}, ylabel={Valeurs},
    xmin=0, xmax=12,
    ymin=0, ymax=1200,
    legend pos=north east,
    grid=major,
    ytick={0,200,350,500,700,1000},
    xtick={0,1,2,3,4,5,6,7,8,9,10,11,12},
]
% CD4 curve (descending)
\addplot[popcurve, mark=none, color=poporange, thick, domain=0:10, samples=100] {800 - 60*x + 10*sin(deg(x*2))};
\addlegendentry{CD4 (cellules/mm$^3$) -- Échelle gauche}
% Viral Load curve (log scale simulated on linear axis for simplicity, or use second axis)
% We'll use a second axis for log VL
\end{axis}
\begin{axis}[
    width=\linewidth, height=9cm,
    xlabel={Temps (années)}, ylabel={Charge Virale (copies/mL)},
    ymode=log, ymin=10, ymax=1e7,
    xmin=0, xmax=12,
    xtick={0,1,2,3,4,5,6,7,8,9,10,11,12},
    ytick={10,100,1000,10000,100000,1000000,10000000},
    yticklabels={$10^1$,$10^2$,$10^3$,$10^4$,$10^5$,$10^6$,$10^7$},
    axis y line*=right, axis x line=none,
    legend pos=south east,
]
% VL curve
\addplot[popcurve, mark=none, color=popblue, thick, domain=0:0.2, samples=50] {1e7 * exp(-20*x) + 50000}; % Acute peak drop
\addplot[popcurve, mark=none, color=popblue, thick, domain=0.2:8, samples=100] {50000 + 10000*sin(deg(x*3))}; % Set point plateau
\addplot[popcurve, mark=none, color=popblue, thick, domain=8:12, samples=50] {50000 * exp(1.5*(x-8))}; % AIDS rise
\addlegendentry{Charge Virale (log) -- Échelle droite}
% Threshold lines
\draw[dashed, popred] (axis cs:0,200) -- (axis cs:12,200) node[above right, pos=0.9, popred, font=\tiny] {Seuil Sida (CD4=200)};
\draw[dashed, popblue] (axis cs:0,1000) -- (axis cs:12,1000) node[above left, pos=0.9, popblue, font=\tiny] {Seuil échec virologique (1000)};
\draw[dashed, popblue] (axis cs:0,50) -- (axis cs:12,50) node[below left, pos=0.9, popblue, font=\tiny] {Seuil détection (50)};
\end{axis}
\end{tikzpicture}
\legende{\textbf{Évolution naturelle (sans ARV) de la Charge Virale (bleu, échelle log) et des CD4 (orange).} Phase 1 : Pic CV + chute CD4 (séroconversion). Phase 2 : Plateau CV (point de consigne) + déclin lent CD4. Phase 3 : Remontée CV + effondrement CD4 $<$ 200 (Sida).}
\end{popfigure}

\begin{aretenir}[Stadification OMS (utilisée au Cameroun pour décision ARV avant "Test \& Treat" universel)]
\begin{itemize}
    \item \textbf{Stade 1 :} Asymptomatique / Adénopathies généralisées persistantes (AGP). CD4 $>$ 500.
    \item \textbf{Stade 2 :} Perte poids $<$ 10\%, infections cutanées/muqueuses récidivantes (herpès, zona, candida), infections bactériennes respiratoires. CD4 350--500.
    \item \textbf{Stade 3 :} Perte poids $>$ 10\%, fièvre/diarrhée chronique $>$ 1 mois, tuberculose pulmonaire, infections bactériennes graves, candidose orale/œsophagienne. CD4 200--350.
    \item \textbf{Stade 4 (Sida) :} Infections opportunistes sévères (Pneumocystose, Toxoplasmose cérébrale, Cryptococcose, CMV, MAC), Cancers (Sarcome Kaposi, Lymphome), Encéphalopathie VIH. CD4 $< 200$.
\end{itemize}
\textbf{Actuel (depuis 2016 au Cameroun) :} \textbf{Test \& Treat (Dépister et Traiter)} $\rightarrow$ ARV initiés \textbf{immédiatement} quel que soit le stade ou le taux de CD4.
\end{aretenir}

\courssub{Dépistage, prise en charge et traitements}

\begin{methode}[Algorithme national de dépistage VIH au Cameroun (Stratégie 3 tests)]
\textbf{Objectif :} Garantir une spécificité finale $> 99,9\%$ (VPP $\approx$ 100\%) avant d'annoncer la séropositivité.
\begin{enumerate}
    \item \textbf{Test 1 (Dépistage / Screening) :} Test Rapide (TR) \textbf{Haute Sensibilité} (ex: \textbf{Determine HIV-1/2}, SD Bioline HIV-1/2 3.0). But : ne rater aucun vrai positif.
    \begin{itemize}
        \item Négatif $\rightarrow$ \textbf{Séronégatif} (sauf fenêtre sérologique $\rightarrow$ retrôtest).
        \item Positif (Réactif) $\rightarrow$ Passer au Test 2.
    \end{itemize}
    \item \textbf{Test 2 (Confirmation 1) :} TR \textbf{Haute Spécificité}, \textbf{principe/antigènes différents} du T1 (ex: \textbf{Uni-Gold Recombigen HIV}, \textbf{Stat-Pak}).
    \begin{itemize}
        \item Positif $\rightarrow$ \textbf{Concordance T1+/T2+ = Séropositivité confirmée}. Prise en charge immédiate.
        \item Négatif $\rightarrow$ \textbf{Discordance T1+/T2-}. Passer au Test 3 (Arbitrage).
    \end{itemize}
    \item \textbf{Test 3 (Arbitrage / Tie-breaker) :} TR ou ELISA (ex: \textbf{Bioline HIV 1/2 3.0}, \textbf{ELISA Ag/Ab}).
    \begin{itemize}
        \item Positif $\rightarrow$ T1+/T3+ = \textbf{Séropositif confirmé} (2/3 positifs).
        \item Négatif $\rightarrow$ T2-/T3- = \textbf{Séronégatif} (2/3 négatifs). T1 était un faux positif.
    \end{itemize}
    \item \textbf{Cas particulier :} Indéterminé (bande floue) $\rightarrow$ Nouveau prélèvement à J+14.
\end{enumerate}
\cle{Terminologie :} On dit « \textbf{Réactif au test de dépistage} » (T1), « \textbf{Confirmé par test de confirmation} » (T2/T3). Jamais « Test positif = Séropositif » après un seul test.
\end{methode}

\begin{definition}[Fenêtre sérologique (Période d'éclipse) et Tests de génération]
Période post-contamination où l'infection est présente mais \textbf{non détectable} par le test utilisé.
\begin{center}
\begin{tabular}{|c|c|c|c|}\hline
\textbf{Génération Test} & \textbf{Cible} & \textbf{Fenêtre moyenne} & \textbf{Exemples} \\ \hline
\textbf{1$^{\text{re}}$ / 2$^{\text{e}}$ (Anciennes)} & Anticorps (Ab) seulement & 6--12 semaines (jusqu'à 6 mois) & Obsolètes \\ \hline
\textbf{3$^{\text{e}}$ génération} & Anticorps (Ab) IgM + IgG (sensibilité $\uparrow$) & \textbf{3--4 semaines} (22--25 j) & Tests rapides Ab (Determine, Uni-Gold, OraQuick salivaire) \\ \hline
\textbf{4$^{\text{e}}$ génération (Combo Ag/Ab)} & \textbf{Antigène p24 + Anticorps (Ab)} & \textbf{2--3 semaines} (15--20 j pour p24) & \textbf{Alere Determine HIV-1/2 Ag/Ab Combo}, ELISA Ag/Ab (labo) \\ \hline
\textbf{Charge Virale (ARN VIH / PCR)} & ARN viral (recherche directe) & \textbf{7--10 jours} & Diagnostic précoce nourrisson (PMTCT), confirmation fenêtre, bilan pré-thérapeutique \\ \hline
\end{tabular}
\end{center}
\textbf{Conduite à tenir face à un test négatif récent :} « Test négatif \textbf{mais exposition $<$ 6 sem. (test Ag/Ab) ou $<$ 3 mois (test rapide Ab)} $\rightarrow$ \textbf{Fenêtre sérologique}. Recommander \textbf{retrôtest à 6 semaines (test Ag/Ab préféré) ou 3 mois}. Prévention (préservatif/PrEP) pendant le délai. »
\end{definition}

\begin{propriete}[Valeurs Prédictives et Prévalence : pourquoi la confirmation est vitale]
\begin{align*}
\text{VPP} &= \frac{\text{VP}}{\text{VP}+\text{FP}} = \frac{\text{Se} \times \text{Prév}}{\text{Se} \times \text{Prév} + (1-\text{Sp}) \times (1-\text{Prév})} \\
\text{VPN} &= \frac{\text{VN}}{\text{VN}+\text{FN}} = \frac{\text{Sp} \times (1-\text{Prév})}{\text{Sp} \times (1-\text{Prév}) + (1-\text{Se}) \times \text{Prév}}
\end{align*}
\textit{Se = Sensibilité, Sp = Spécificité, Prév = Prévalence.}
\textbf{Règle fondamentale :}
\begin{itemize}
    \item \textbf{Faible prévalence (ex: population générale 3\%)} $\rightarrow$ \textbf{VPP baisse} (beaucoup de faux positifs relatifs parmi les positifs) $\rightarrow$ \textbf{Confirmation par T2/T3 IMPÉRATIVE}.
    \item \textbf{Haute prévalence (ex: population clé 20\%)} $\rightarrow$ \textbf{VPN baisse} (risque faux négatifs) $\rightarrow$ Vigilance fenêtre sérologique, test Ag/Ab ou PCR.
\end{itemize}
\end{propriete}

\begin{experience}[Prise en charge : De l'annonce à la suppression virale (Parcours patient Cameroun)]
\textbf{1. Annonce post-test positif (Conseil post-test) :} Accueil, écoute, annonce claire (« Vous êtes séropositif »), évaluation psychosociale, information sur ARV (efficacité, effets secondaires, observance), dépistage partenaire(s) / enfants (index testing), orientation vers Unité de Prise en Charge (PEC).
\textbf{2. Bilan pré-thérapeutique (BPT) :} CV (baseline), CD4, Créatininémie (rein $\rightarrow$ TDF), Hémogramme (AZT), Glycémie/Lipides (IP), HBsAg / Anti-HCV (co-infections), Grossesse (femme âge procréer), TB (recherche BK / GeneXpert).
\textbf{3. Initiation ARV (1$^{\text{re}}$ ligne recommandée OMS/Cameroun 2023+) :} \textbf{TDF 300 mg + 3TC 300 mg (ou FTC 200 mg) + DTG 50 mg} (1 cp/jour, prise unique, avec ou sans repas). \emph{Alternative si CI TDG (femme désir grossesse non contraceptée sans assurance folates) : TDF/3TC/EFV 600 mg (prise soir, jeun).}
\textbf{4. Suivi \& Monitoring :}
\begin{itemize}
    \item \textbf{M3, M6 :} CV (objectif $< 50$ copies/mL à M6), CD4, Créat (si TDF), Hémato (si AZT), Poids, Observance, Effets indésirables.
    \item \textbf{Tous les 6 mois ensuite :} CV (si $< 50$ stable $\rightarrow$ espacer à 12 mois possible), CD4 (si $> 500$ stable $\rightarrow$ arrêter CD4 possible).
    \item \textbf{Échec virologique :} CV $> 1000$ copies/mL à 2 mesures consécutives (intervalle 3 mois) APRÈS conseil observance renforcé $\rightarrow$ Génotypage $\rightarrow$ Switch 2$^{\text{e}}$ ligne.
\end{itemize}
\end{experience}

\begin{definition}[Traitement Antirétroviral (ARV) : Classes, Mécanismes, Lignes au Cameroun]
\begin{center}
\begin{tabular}{|c|l|l|l|}\hline
\textbf{Classe (Abrév.)} & \textbf{Mécanisme / Cible} & \textbf{Molécules clés (Cameroun)} & \textbf{Effets indésirables majeurs / Surveillance} \\ \hline
\textbf{INRTI} (NRTI) & Analogues nucléos(t)idiques $\rightarrow$ Terminaison chaîne ADN (RT) & \textbf{TDF} (Ténofovir), \textbf{3TC} (Lamivudine), \textbf{FTC}, \textbf{AZT} (Zidovudine), \textbf{ABC} & TDF : Rénal (créat, phosphaturie), Osseux (ostéoporose). AZT : Anémie, Neutropénie (hémato). ABC : Hypersensibilité (HLA-B*5701). \\ \hline
\textbf{INNRTI} (NNRTI) & Liaison site allostérique RT $\rightarrow$ Inhibition non compétitive & \textbf{EFV} (Efavirenz), \textbf{NVP} (Névirapine), \textbf{DOR} (Doravirine), \textbf{ETR} (Etravirine - 2$^{\text{e}}$ gén) & EFV : Neuropsychiatrique (rêves, dépression, suicidaire), Tératogène (1$^{\text{er}}$ trim). NVP : Hépatite, Rash (lead-in 14j). \\ \hline
\textbf{II} (INSTI) & Blocage transfert brin (Intégrase) $\rightarrow$ Pas d'intégration & \textbf{DTG} (Dolutégravir - 1$^{\text{re}}$ ligne), \textbf{RAL} (Raltegravir), \textbf{BIC} (Bictegravir), \textbf{CAB} (Cabotégravir - injectable) & DTG : Prise de poids, Insomnie, Hyperglycémie, \textbf{Risque TNN (défaut tube neural) si conception sous DTG} (supplémentation folates 5mg/j). \\ \hline
\textbf{IP} (PI) & Inhibition Protéase $\rightarrow$ Virions immatures non infectieux & \textbf{ATV/r} (Atazanavir/ritonavir), \textbf{LPV/r} (Lopinavir/ritonavir), \textbf{DRV/r} (Darunavir/ritonavir) & \textbf{Boosté par Ritonavir (r) ou Cobicistat (c)} (inhibent CYP3A4 $\rightarrow$ interactions médicamenteuses majeures). Lipides, Glycémie, Néphrolithiase (ATV), Diarrhée (LPV). \\ \hline
\textbf{Autres} & Entrée (CCR5 : MVC), Fusion (ENF), Attachement (FTR), Capside (LEN - injectable 6 mois) & Maraviroc, Enfuvirtide, Fostemsavir, Lenacapavir & Réservés échecs multi-lignes / intolérances. Coût élevé. \\ \hline
\end{tabular}
\end{center}
\textbf{1$^{\text{re}}$ ligne standard (Cameroun 2024) :} \textbf{TDF + 3TC (ou FTC) + DTG} (TLD). \textbf{2$^{\text{e}}$ ligne standard (échec 1$^{\text{re}}$ ligne INRTI+INNRTI) :} \textbf{ATV/r (ou LPV/r) + AZT + 3TC}. \textbf{3$^{\text{e}}$ ligne :} Comité d'experts, molécules de sauvetage (DRV/r, DTG si non utilisé, ETR, MVC, LEN).
\end{definition}

\begin{propriete}[TasP (Treatment as Prevention) / I=I (Indétectable = Intransmissible) : Preuve scientifique]
\begin{itemize}
    \item \textbf{Essais majeurs :} HPTN 052 (couples sérodiscordants, 2011), PARTNER 1 \& 2 (couples gays/hétéros, 2016/2018), Opposites Attract.
    \item \textbf{Résultat :} \textbf{Zéro transmission sexuelle liée} chez les couples où le partenaire $+$ avait CV $< 200$ copies/mL (seuil OMS) / $< 50$ copies/mL (seuil labo standard) \textbf{durablement} (généralement $> 6$ mois) et sous observance.
    \item \textbf{Message clé :} \textbf{I = I}. Une personne vivant avec le VIH, sous ARV efficace, charge virale indétectable, \textbf{ne transmet pas le VIH par voie sexuelle}. C'est un argument majeur pour le dépistage, l'observance, la déstigmatisation et le projet parental naturel.
\end{itemize}
\end{propriete}

\begin{definition}[PrEP (Prophylaxie Pré-Exposition) et PEP (Post-Exposition)]
\begin{itemize}
    \item \textbf{PrEP :} ARV pris \textbf{avant} exposition par séronégatif à risque substantiel.
        \begin{itemize}
            \item \textbf{Schémas :} \textbf{Quotidien} (1 cp TDF/3TC ou TDF/FTC / jour) : \textbf{Tous profils} (HSH, Hétéro, UDI, Trans). Protection optimale : 7j (rectal), 20j (vaginal/cervical). \textbf{À la demande (2-1-1)} : 2 cp 2-24h avant, 1 cp 24h après, 1 cp 24h après $\rightarrow$ \textbf{Uniquement validé pour HSH} (études IPERGAY, PREVENIR).
            \item \textbf{Suivi :} Sérologie VIH trimestrielle (détecter séroconversion $\rightarrow$ éviter mono/bi-thérapie), Créatininémie (rein), Dépistage IST, Conseil observance.
        \end{itemize}
    \item \textbf{PEP :} ARV d'urgence \textbf{après} exposition (AES pro, viol, rapport non protégé partenaire $+$ CV détectable/inconnue, partage seringue).
        \begin{itemize}
            \item \textbf{Délai critique :} Idéalement $< 4$h, \textbf{maximum 48h} (après $\rightarrow$ inefficace).
            \item \textbf{Durée :} \textbf{28 jours} complets.
            \item \textbf{Régime recommandé :} TDF/3TC/DTG (ou TDF/3TC/ATV/r ou LPV/r).
            \item \textbf{Suivi :} Sérologie J0, M1, M3 (6 sem), M6. Effets indésirables (nausées, maux de tête fréquents 1$^{\text{re}}$ sem).
        \end{itemize}
\end{itemize}
\end{definition}

\begin{definition}[PMTCT (Prévention Transmission Mère-Enfant) : Option B+ au Cameroun]
\begin{enumerate}
    \item \textbf{Dépistage systématique} 1$^{\text{re}}$ CPN (Consultation Prénatale). Test partenaire (index testing).
    \item \textbf{Femme $+$ :} \textbf{ARV à vie (Option B+)} immédiats (TDF/3TC/DTG 1cp/j) quel que soit CD4, stade, terme grossesse.
    \item \textbf{Accouchement :} Voie basse si CV $< 1000$ (ou indétectable). Césarienne programmée si CV $> 1000$ ou inconnue au travail.
    \item \textbf{Nouveau-né :} \textbf{Prophylaxie ARV} : AZT sirop (4 mg/kg/12h) + 3TC sirop (2 mg/kg/12h) pendant \textbf{6 semaines} (si mère ARV $>$ 4 sem + CV indétectable) ou \textbf{12 semaines} (si mère ARV $<$ 4 sem, CV détectable, ou allaitement). Vitamine D / Fer / Cotrimoxazole (si exposé).
    \item \textbf{Alimentation :} \textbf{Allaitement exclusif 6 mois} (recommandation OMS/Cameroun même si mère $+$ sous ARV efficace) $\rightarrow$ Introduction aliment complémentaire + sevrage progressif 12-24 mois. \emph{Alternative : Lait artificiel SI conditions AFASS (Acceptable, Faisable, Abordable, Durable, Sûr) réunies.}
    \item \textbf{Dépistage enfant :} PCR ADN VIH (PoC ou Labo) à \textbf{6 semaines} (ou naissance si haut risque), puis 6 sem. après sevrage, puis test rapide 18 mois (confirmation sérologique).
\end{enumerate}
\cle{Objectif Cameroun :} Taux transmission mère-enfant $< 5\%$ (cible élimination $< 2\%$).
\end{definition}

\begin{propriete}[Résistance aux ARV : Mécanisme, Conséquences, Surveillance]
\begin{itemize}
    \item \textbf{Mécanisme moléculaire :} RT sans proofreading $\rightarrow$ $10^{-5}$ erreurs/base/cycle. Réplication $\approx 10^{10}$ virions/jour $\rightarrow$ \textbf{Tous les mutants simples existent quotidiennement}. Sous pression ARV (concentration sous-CMI par inobservance), mutants résistants sélectionnés $\rightarrow$ deviennent majoritaires (quasi-espèces).
    \item \textbf{Barrière génétique :} Nombre de mutations nécessaires pour résistance haut niveau. \textbf{DTG / IP boostés : Barrière haute} (3+ mutations). \textbf{INNRTI (EFV/NVP) / 3TC/FTC : Barrière basse} (1 mutation : K103N, Y181C / M184V).
    \item \textbf{Résistance croisée :} Mutation conférant résistance à plusieurs molécules d'une même classe. Ex: K103N $\rightarrow$ Résistance \textbf{totale EFV + NVP} (1$^{\text{re}}$ gén NNRTI). M184V $\rightarrow$ Résistance \textbf{3TC + FTC} (mais hypersensibilise à TDF/AZT).
    \item \textbf{Résistance transmise (Primaire) :} Infection par souche déjà résistante. Prévalence $\approx$ 5-10\% au Cameroun $\rightarrow$ Justifie test de résistance (génotypage) avant initiation si possible, ou choix 1$^{\text{re}}$ ligne à haute barrière (DTG).
    \item \textbf{Génotypage (Test de résistance) :} Séquençage gènes \emph{pol} (RT, PR, IN). Indiqué : Échec virologique confirmé, Infection aiguë/récente, Femme enceinte naïve ARV (optionnel), Enfant exposé.
\end{itemize}
\end{propriete}

\begin{methode}[Interpréter l'évolution de la Charge Virale (CV) et des CD4 sous ARV]
\textbf{Objectif :} Lire courbes/tableaux pour identifier stade clinique et efficacité thérapeutique.
\begin{enumerate}
    \item \textbf{Axes \& Unités :} CV (copies/mL, \textbf{échelle log} : $10^2$ à $10^6$). Seuil détection 20-50. Seuil échec OMS : \textbf{1000 copies/mL}. CD4 (cell/mm$^3$). Norme 500-1500. Seuils : $<500$ (modéré), $<350$ (sévère), $<200$ (Sida).
    \item \textbf{Infection naturelle (non traitée) :} Pic aigu $\rightarrow$ Point de consigne (plateau CV) $\rightarrow$ Déclin lent CD4 ($\approx$ -50 à -80/an) $\rightarrow$ Remontée CV + Effondrement CD4 (Sida).
    \item \textbf{Sous ARV efficace (Succès virologique) :} \textbf{Décroissance biphasique CV :} Phase 1 (J0-14) : chute rapide (blocage nouvelles infections, 1/2 vie virions $\approx$ 6h). Phase 2 (Sem 2-24) : pente lente (mort cellules infectées vie longue, réservoirs). \textbf{Objectif : CV $< 50$ à M6.} CD4 : Remontée (gain 50-150 1$^{\text{re}}$ an, puis lente).
    \item \textbf{Échec virologique :} CV $> 1000$ copies/mL à 2 mesures (M6 et M9 ou M12) APRÈS observance optimisée. $\rightarrow$ Génotypage $\rightarrow$ Switch.
    \item \textbf{Échec immunologique :} CV indétectable MAIS CD4 $< 350$ ou gain $< 50-100$/an $\rightarrow$ Recherche cause (co-infections, inflammation, âge, toxicité, non-adhérence partielle).
    \item \textbf{Blip (Poussée virale transitoire) :} CV faible ($50$--$1000$) isolée, retour spontané $< 50$ $\rightarrow$ \textbf{Pas d'échec}, pas de changement ARV. Surveiller.
\end{enumerate}
\textbf{Piège :} Ne pas confondre \textbf{log copies/mL} (ex: 4 log = 10 000 copies) et copies/mL brutes. Une baisse de 1 log = division par 10 de la CV.
\end{methode}

\begin{methode}[Argumenter l'observance thérapeutique : Enjeux individuel et collectif]
\textbf{Observance =} Prise médicaments \textbf{bon moment, bonne dose, sans oubli, sans interruption}, conditions respectées. Objectif \textbf{$> 95\%$} (DTG plus tolérant mais viser 100\%).
\begin{enumerate}
    \item \textbf{Efficacité virologique (Mécanisme) :} ARV bloquent réplication $\rightarrow$ CV indétectable. Oubli $\rightarrow$ Concentration $<$ CMI $\rightarrow$ Réplication $\rightarrow$ Mutations RT $\rightarrow$ Sélection résistants $\rightarrow$ Échec virologique (CV $> 1000$).
    \item \textbf{Conséquences cliniques patient :} Chute CD4, Infections Opportunistes, Décès. \textbf{Résistance croisée} (ex: NNRTI 1$^{\text{re}}$ gén $\rightarrow$ perte classe entière) $\rightarrow$ Passage 2$^{\text{e}}$/3$^{\text{e}}$ ligne (plus chères, + effets secondaires, + pilules).
    \item \textbf{Santé publique (TasP / I=I) :} Observance $\rightarrow$ CV indétectable $\rightarrow$ \textbf{Zéro transmission sexuelle}. Non-observance $\rightarrow$ Transmission + \textbf{Transmission souches résistantes} (résistance primaire chez nouveau infecté $\rightarrow$ limite options entrée).
    \item \textbf{Soutien pratique :} Boîte pilulier, alarmes, applications, Groupes d'Aide Paroles (GAP), Associations PVVIH, Prise en charge effets indésirables (ne pas cacher), Dépistage dépression/addictions.
\end{enumerate}
\textbf{Phrase choc :} « \textbf{Prendre ses ARV tous les jours, c'est se protéger soi-même, protéger l'autre, et préserver ses traitements futurs.} »
\end{methode}

\courssub{Méthodes et exercices résolus}

\begin{methode}[Analyser une situation à risque et identifier le mode de transmission du VIH]
\textbf{Objectif :} Expliquer les modes de transmission du VIH à partir d'un cas concret.
\begin{enumerate}
    \item \textbf{Identifier le fluide biologique mis en jeu :} Sang, sperme, sécrétions vaginales, lait maternel (fluides contaminants). Salive, sueur, larmes, urine, selles (non contaminants sauf contamination sanguine visible).
    \item \textbf{Repérer la porte d'entrée :} Muqueuses (génitales, rectales, buccales), peau lésée, voie intraveineuse, voie transplacentaire/pernatale.
    \item \textbf{Qualifier le mode de transmission :}
    \begin{itemize}
        \item \textbf{Sexuel :} Rapports non protégés (vaginaux, anaux, oraux \emph{receptif} avec éjaculation).
        \item \textbf{Sanguin :} Partage de seringues/aiguilles, matériel de scarification/tatouage non stérilisé, transfusion non dépistée (rare aujourd'hui), accident d'exposition au sang (AES).
        \item \textbf{Mère-enfant (verticale) :} \emph{In utero} (transplacentaire), \emph{per partum} (au passage, contact sang/sécrétions), \emph{post partum} (allaitement).
    \end{itemize}
    \item \textbf{Écarter les fausses voies :} Contacts quotidiens (serrer la main, embrasser, partager repas/toilettes, piqûres de moustiques).
    \item \textbf{Rédiger la réponse :} « Le mode de transmission est \textbf{...} car le fluide \textbf{...} contenant le virus a franchi la barrière \textbf{...} lors de \textbf{...} ».
\end{enumerate}
\textbf{Reflexe Probatoire :} Ne jamais écrire « transmission par la salive » sans préciser « contaminée par le sang ». Préciser toujours le fluide \textbf{et} la porte d'entrée.
\end{methode}

\begin{exoresolu}[Analyse de situations à risque au Cameroun]
\textbf{Énoncé :} Trois patients consultent au Centre de Dépistage Volontaire (CDV) de l'Hôpital Central de Yaoundé.
\begin{itemize}
    \item \textbf{Patient A :} Homme de 28 ans, déclare avoir eu un rapport sexuel vaginal non protégé il y a 3 semaines avec une partenaire de statut sérologique inconnu.
    \item \textbf{Patient B :} Femme de 35 ans, séropositive connue sous ARV depuis 2 ans (charge virale indétectable), enceinte de 6 mois, inquiète pour son enfant.
    \item \textbf{Patient C :} Jeune de 20 ans, usager de drogues injectables, partage sa seringue avec deux camarades. Il présente une plaie suintante au bras.
\end{itemize}
Pour chaque patient, identifier le mode de transmission principal encouru, le fluide contaminant et la porte d'entrée. Justifier pourquoi le patient B a un risque de transmission mère-enfant \textbf{faible} mais non nul.
\tcblower
\textbf{Solution détaillée :}
\begin{itemize}
    \item \textbf{Patient A :}
    \begin{itemize}
        \item \textbf{Mode :} Transmission \textbf{sexuelle} (voie hétérosexuelle).
        \item \textbf{Fluide :} \textbf{Sperme} (ou sécrétions vaginales de la partenaire).
        \item \textbf{Porte d'entrée :} \textbf{Muqueuse génitale} (pénienne) : micro-lésions souvent invisibles lors du rapport.
        \item \textit{Justification :} Le VIH est présent dans le sperme et les sécrétions vaginales. L'absence de préservatif permet le contact direct fluide/muqueuse.
    \end{itemize}
    \item \textbf{Patient B :}
    \begin{itemize}
        \item \textbf{Mode :} Transmission \textbf{mère-enfant (verticale)}.
        \item \textbf{Fluides concernés :} \textbf{Sang maternel} (accouchement), \textbf{Lait maternel} (allaitement). \emph{In utero} : passage transplacentaire rare si barrière placentaire intacte.
        \item \textbf{Portes d'entrée :} Muqueuses fœtales (accouchement), tube digestif du nourrisson (allaitement).
        \item \textbf{Risque faible justifié :} La patiente est sous ARV avec \textbf{charge virale indétectable} ($< 50$ copies/mL). Le principe \textbf{TasP (Treatment as Prevention) / I=I (Indétectable = Intransmissible)} s'applique : le risque de transmission \emph{in utero} et \emph{per partum} est quasi nul ($< 1\%$). Cependant, le risque \textbf{post-partum par l'allaitement} persiste si l'allaitement est pratiqué (recommandation OMS/Cameroun : allaitement exclusif 6 mois \textbf{sous couverture ARV maternelle et prophylaxie infantile} pour réduire le risque à $< 1\%$).
    \end{itemize}
    \item \textbf{Patient C :}
    \begin{itemize}
        \item \textbf{Mode :} Transmission \textbf{sanguine} (partage de matériel d'injection).
        \item \textbf{Fluide :} \textbf{Sang} contaminé résidant dans l'aiguille/seringue (volume résiduel « dead space »).
        \item \textbf{Porte d'entrée :} \textbf{Voie intraveineuse directe} (injection) \textbf{ET} peau lésée (plaie suintante = porte d'entrée supplémentaire pour contact sang-sang).
        \item \textit{Justification :} L'injection directe inocule le virus dans la circulation sanguine. C'est la voie la plus efficace (probabilité de transmission $\approx$ 90\% par acte de partage).
    \end{itemize}
\end{itemize}
\textbf{Conclusion :} L'identification précise du fluide et de la porte d'entrée permet de classer le risque et d'orienter la prophylaxie post-exposition (PEP) pour A et C, et le suivi PMTCT (Prévention de la Transmission Mère-Enfant) pour B.
\end{exoresolu}

\begin{methode}[Choisir et argumenter une stratégie de prévention adaptée (Combinaison préventive)]
\textbf{Objectif :} Proposer des moyens de prévention primaires et secondaires adaptés au profil de risque.
\begin{enumerate}
    \item \textbf{Évaluer le niveau de risque :} Comportement sexuel (nombre partenaires, type rapports), usage de drogues injectables, contexte professionnel (AES), statut sérologique du partenaire.
    \item \textbf{Mobiliser la « Combinaison préventive » (OMS/ONUSIDA) :}
    \begin{itemize}
        \item \textbf{Comportementale :} Abstinence, fidélité mutuelle dépistée, réduction du nombre de partenaires, \textbf{usage correct et systématique du préservatif (masculin/féminin)}.
        \item \textbf{Biomédicale :} \textbf{PrEP (Prophylaxie Pré-Exposition)} : ARV quotidiens (TDF/3TC ou TDF/FTC) pour séronégatifs à risque élevé. \textbf{PEP (Prophylaxie Post-Exposition)} : ARV d'urgence $< 48$h (idéalement $< 4$h) pendant 28 jours après exposition. \textbf{TasP / I=I} : Traitement du partenaire séropositif pour rendre charge virale indétectable. \textbf{PMTCT} : ARV mère + enfant + accouchement sécurisé + conseil alimentation. \textbf{Circoncision médicale volontaire} (réduit risque acquisition homme HSH/HSH de 60\%).
        \item \textbf{Structurelle :} Éducation sexuelle complète, réduction des stigmatisations, accès aux soins, programmes d'échange de seringues (RdR).
    \end{itemize}
    \item \textbf{Argumenter le choix :} Lier la méthode au mode de transmission ciblé (ex: PrEP ne protège pas des IST bactériennes $\rightarrow$ associer préservatif).
    \item \textbf{Préciser les conditions d'efficacité :} Observance (PrEP/PEP/ARV), délai (PEP), qualité du produit (préservatif non périmé, norme CE/ISO).
\end{enumerate}
\textbf{Formule clé :} \textbf{Prévention Combinée = Comportementale + Biomédicale + Structurelle}. Aucune méthode seule n'est à 100\% (sauf abstinence totale), la combinaison maximise la protection.
\end{methode}

\begin{exoresolu}[Élaboration d'un plan de prévention personnalisé]
\textbf{Énoncé :} \textbf{Paul}, 24 ans, étudiant à Douala, séronégatif au dernier test (il y a 3 mois). Il a une relation stable avec \textbf{Chantal}, 22 ans, qui vient d'apprendre sa séropositivité (charge virale = 45 000 copies/mL, CD4 = 450/mm$^3$). Chantal refuse de commencer les ARV « pour l'instant » car elle « ne se sent pas malade ». Paul ne veut pas utiliser de préservatif « pour faire un enfant rapidement ». Il consulte l'infirmière du CDV pour avoir la PrEP.
\begin{enumerate}
    \item Expliquer pourquoi la situation de Paul constitue un \textbf{risque élevé} d'acquisition du VIH.
    \item La PrEP est-elle indiquée ici ? Justifier en citant les conditions d'éligibilité et les limites.
    \item Proposer une \textbf{stratégie combinée} réaliste pour ce couple (sérodiscordant) en intégrant le désir d'enfant.
\end{enumerate}
\tcblower
\textbf{Solution détaillée :}
\begin{enumerate}
    \item \textbf{Risque élevé :} Couple \textbf{sérodiscordant} (Paul $-$, Chantal $+$). Chantal a une \textbf{charge virale élevée} (45 000 copies/mL) $\rightarrow$ forte infectiosité. Absence de préservatif (porte d'entrée muqueuses génitales exposées au sperme/sécrétions infectées). Désir d'enfant $\rightarrow$ rapports fréquents non protégés pendant la fenêtre de fertilité. Risque de transmission par acte vaginal $\approx$ 0,1\% à 1\% (plus élevé si charge virale haute, IST associées, phase aiguë).
    \item \textbf{PrEP pour Paul :} \textbf{OUI, fortement indiquée.}
    \begin{itemize}
        \item \textbf{Critères éligibilité (OMS/Cameroun) :} Personne séronégative, risque substantiel (partenaire séropositif non supprimé viralement, rapports non protégés).
        \item \textbf{Schémas :} Quotidien (TDF/3TC ou TDF/FTC) ou « à la demande » (2-1-1) \emph{uniquement pour HSH} $\rightarrow$ Paul hétérosexuel $\rightarrow$ \textbf{Quotidien obligatoire}.
        \item \textbf{Limites :} Ne protège \textbf{que contre le VIH}. Pas de protection contre IST (chlamydia, gonococcie, syphilis, HPV, herpès) ni grossesse non désirée (mais désiré ici). Nécessite observance stricte (7 jours pour protection rectale, 20 jours pour protection vaginale/cervicale). Suivi rénal (créatininémie) et sérologique trimestriel.
    \end{itemize}
    \item \textbf{Stratégie combinée réaliste (Approche couple) :}
    \begin{enumerate}
        \item \textbf{Urgence :} Paul démarre \textbf{PrEP quotidienne} immédiatement (protection effective après 20 jours pour voie vaginale). Conseils observance.
        \item \textbf{Objectif prioritaire :} \textbf{Convaincre Chantal de démarrer les ARV} (TasP). Arguments : « Tes ARV protègent ta santé (CD4 ne doivent pas baisser) ET protègent Paul \textbf{sans préservatif} dès que ta charge virale devient indétectable (généralement 3 à 6 mois). C'est la meilleure méthode pour faire un enfant sereinement ». \textbf{I=I (Indétectable = Intransmissible)}.
        \item \textbf{Projet grossesse (si Chantal sous ARV + CV indétectable $> 6$ mois) :} Rapports non protégés \textbf{uniquement pendant la fenêtre ovulatoire} (période fertile). Arrêt PrEP pour Paul si CV Chantal indétectable confirmée (avis médical).
        \item \textbf{Si Chantal refuse ARV :} PrEP Paul continue + \textbf{Préservatif obligatoire hors période fertile} + Dépistage Paul tous les 3 mois + Dépistage/traitement IST couple.
        \item \textbf{Suivi :} Consultation couple au CDV / Unité PMTCT pour planification conception sécurisée.
    \end{enumerate}
\end{enumerate}
\textbf{Conclusion :} La PrEP protège Paul \textbf{maintenant}. L'ARV pour Chantal protège le couple \textbf{durablement} et permet la conception naturelle sans risque (TasP). La prévention combinée est la seule réponse éthique et efficace.
\end{exoresolu}

\begin{methode}[Interpréter l'évolution de la charge virale (CV) et des lymphocytes T CD4 au cours de l'infection]
\textbf{Objectif :} Lire des courbes (CV vs temps, CD4 vs temps) ou des tableaux de valeurs pour identifier le stade clinique et l'efficacité thérapeutique.
\begin{enumerate}
    \item \textbf{Repérer les axes et unités :}
    \begin{itemize}
        \item Charge Virale (CV) : \textbf{copies/mL} (échelle souvent \textbf{logarithmique} : $10^2$, $10^3$, $10^4$, $10^5$, $10^6$). Seuil détection : 20-50 copies/mL. Seuil échec virologique OMS : \textbf{1000 copies/mL}.
        \item Lymphocytes T CD4 : \textbf{cellules/mm$^3$} (ou $\mu$L). Norme adulte : 500--1500. Seuils : $<$ 500 (immunodépression modérée), $<$ 350 (sévère), $<$ 200 (Sida déclaré, risque infections opportunistes majeures).
    \end{itemize}
    \item \textbf{Identifier les phases caractéristiques (infection naturelle non traitée) :}
    \begin{itemize}
        \item \textbf{Phase aiguë (Primo-infection / Séroconversion) :} Pic de CV \textbf{très élevé} ($> 10^6$ copies/mL), chute brutale des CD4 (souvent $< 500$), puis remontée partielle. Apparition anticorps (séroconversion).
        \item \textbf{Phase asymptomatique (Latence clinique) :} CV se stabilise au « \textbf{point de consigne viral} » (ex: $10^4$--$10^5$). CD4 déclinent \textbf{lentement et linéairement} ($\approx$ 50--80 cellules/mm$^3$/an).
        \item \textbf{Phase Sida (Symptomatique) :} CV remonte exponentiellement, CD4 s'effondrent ($< 200$), infections opportunistes (IO) / cancers.
    \end{itemize}
    \item \textbf{Analyser l'impact des ARV (Traitement) :}
    \begin{itemize}
        \item \textbf{Succès virologique :} CV $\rightarrow$ \textbf{indétectable} ($< 50$ copies/mL) en 3--6 mois. CD4 $\rightarrow$ \textbf{remontée} (gain 50--150 la 1$^{\text{re}}$ année, puis lente).
        \item \textbf{Échec virologique :} CV $> 1000$ copies/mL à 6 mois (ou rebond confirmé $> 200$ copies/mL après suppression). $\rightarrow$ Risque résistance $\rightarrow$ Génotypage / Changement de ligne.
        \item \textbf{Échec immunologique :} CV indétectable MAIS CD4 ne remontent pas ($< 350$ ou gain $< 50-100$/an) $\rightarrow$ Recherche cause (inflammation, co-infections, âge, toxicité).
    \end{itemize}
    \item \textbf{Rédiger l'interprétation :} « Entre $t_1$ et $t_2$, la CV passe de ... à ... (baisse de $\log_{10}$), les CD4 passent de ... à ... (hausse de ... cellules/mm$^3$). Cela traduit un \textbf{succès virologique et immunologique} sous ARV. »
\end{enumerate}
\textbf{Piège :} Ne pas confondre \textbf{log copies/mL} (ex: 4 log = 10 000 copies) et copies/mL brutes. Une baisse de 1 log = division par 10 de la CV.
\end{methode}

\begin{exoresolu}[Lecture de courbes d'évolution sous traitement]
\textbf{Énoncé :} Le graphique ci-dessous représente l'évolution de la charge virale (CV, en copies/mL, échelle log) et du taux de lymphocytes T CD4 (cellules/mm$^3$) chez \textbf{Mme Ngo}, 38 ans, découverte séropositive lors d'une grossesse (Mois 0). Elle démarre les ARV (TDF/3TC/DTG) immédiatement (Option B+). Les points de mesure sont M0, M3, M6, M12, M24.
\begin{popfigure}
\begin{tikzpicture}
\begin{axis}[
    width=\linewidth, height=8cm,
    xlabel={Temps (mois)}, ylabel={Charge Virale (copies/mL)},
    ymode=log, ymin=10, ymax=1e6,
    xtick={0,3,6,12,24}, xticklabels={M0, M3, M6, M12, M24},
    ytick={10,100,1000,10000,100000,1000000},
    yticklabels={$10^1$,$10^2$,$10^3$,$10^4$,$10^5$,$10^6$},
    legend pos=north east,
    grid=major,
]
\addplot[popcurve, mark=*, color=popblue, thick] coordinates {
    (0, 500000)
    (3, 2000)
    (6, 40)
    (12, 20)
    (24, 20)
};
\addlegendentry{Charge Virale}
\draw[dashed, popred, thick] (axis cs:0,50) -- (axis cs:24,50) node[above right, pos=0.95, popred, font=\small] {Seuil détection (50)};
\draw[dashed, poporange, thick] (axis cs:0,1000) -- (axis cs:24,1000) node[above right, pos=0.05, poporange, font=\small] {Seuil échec (1000)};
\end{axis}
\begin{axis}[
    width=\linewidth, height=8cm,
    xlabel={Temps (mois)}, ylabel={CD4 (cellules/mm$^3$)},
    ymin=0, ymax=800,
    xtick={0,3,6,12,24}, xticklabels={M0, M3, M6, M12, M24},
    ytick={0,200,350,500,650,800},
    axis y line*=right, axis x line=none,
    legend pos=south east,
]
\addplot[popcurve, mark=square*, color=poporange, thick] coordinates {
    (0, 280)
    (3, 320)
    (6, 380)
    (12, 480)
    (24, 580)
};
\addlegendentry{CD4}
\draw[dashed, poppurple, thick] (axis cs:0,200) -- (axis cs:24,200) node[above left, pos=0.95, poppurple, font=\small] {Seuil Sida (200)};
\draw[dashed, poppurple, thick] (axis cs:0,350) -- (axis cs:24,350) node[below left, pos=0.95, poppurple, font=\small] {Seuil sévère (350)};
\draw[dashed, poppurple, thick] (axis cs:0,500) -- (axis cs:24,500) node[below left, pos=0.05, poppurple, font=\small] {Norme (500)};
\end{axis}
\end{tikzpicture}
\legende{\textbf{Évolution biologique de Mme Ngo sous ARV (TDF/3TC/DTG) initiés à M0.} Charge virale (courbe bleue, échelle log) et CD4 (courbe orange). Seuil détection CV = 50 copies/mL (pointillés rouges). Seuil échec CV = 1000 copies/mL (pointillés orange). Seuils CD4 : 200 (Sida), 350 (sévère), 500 (norme) en pointillés violets.}
\end{popfigure}
\begin{enumerate}
    \item Décrire l'évolution de la charge virale entre M0 et M6. Calculer la baisse en $\log_{10}$ entre M0 et M3. Est-ce un succès virologique à M6 ?
    \item Décrire l'évolution des CD4. Quel était le stade immunologique à M0 ? À M24 ?
    \item À M3, la CV est à 2000 copies/mL. L'infirmière s'inquiète. Argumenter pour la rassurer en utilisant la notion de « décroissance biphasique ».
    \item Conclure sur l'efficacité du traitement à M24.
\end{enumerate}
\tcblower
\textbf{Solution détaillée :}
\begin{enumerate}
    \item \textbf{Évolution CV :} M0 : $\approx 5 \times 10^5$ copies/mL (500 000). Chute drastique à M3 : $\approx 2 \times 10^3$ (2 000). Devient \textbf{indétectable} ($< 50$) à M6 (valeur 40). Reste indétectable M12, M24.
    \textbf{Baisse en log M0-M3 :} $\log_{10}(500\,000) \approx 5,7$ ; $\log_{10}(2\,000) \approx 3,3$. Baisse = $5,7 - 3,3 = \textbf{2,4 log}_{10}$ (division par $\approx 250$). C'est une excellente réponse (attendu $> 1$ log à M1, $> 2$ log à M3).
    \textbf{Succès virologique à M6 :} \textbf{OUI}. CV $< 50$ copies/mL (seuil détection standard) atteint à M6 (délai standard $< 6$ mois).
    \item \textbf{Évolution CD4 :} M0 : 280 cellules/mm$^3$. Remontée lente mais régulière : 320 (M3), 380 (M6), 480 (M12), 580 (M24). Gain total = 300 cellules/mm$^3$ en 2 ans.
    \textbf{Stade M0 :} CD4 = 280 $<$ 350 $\rightarrow$ \textbf{Immunodépression sévère} (Stade 3 OMS / Sida immunologique). Risque IO modéré.
    \textbf{Stade M24 :} CD4 = 580 $>$ 500 $\rightarrow$ \textbf{Normalisation immunitaire} (Stade 1). Risque IO quasi nul.
    \item \textbf{Rassurer à M3 (CV = 2000) :} La décroissance de la CV sous ARV est \textbf{biphasique} :
    \begin{itemize}
        \item \textbf{Phase 1 (Jours 0--14) :} Chute très rapide (1/2 vie virions $\approx$ 6h). Blocage infection nouvelles cellules (Inhibiteurs Entrée/Intégrase/RT). CV divisée par 100--1000.
        \item \textbf{Phase 2 (Semaines 2--24) :} Pente plus lente. Élimination cellules infectées à vie longue (macrophages, réservoirs latents). 1/2 vie $\approx$ 2 semaines.
    \end{itemize}
    À M3 (12 sem), on est en \textbf{fin de phase 2}. Une CV à 2000 copies/mL est \textbf{attendue et normale} pour une CV initiale haute (500 000). La poursuite de la baisse vers l'indétectabilité à M6 confirme l'efficacité. L'inquiétude n'est justifiée que si CV $> 1000$ à M6 (échec).
    \item \textbf{Conclusion M24 :} \textbf{Succès thérapeutique complet.}
    \begin{itemize}
        \item \textbf{Virologique :} CV durablement indétectable ($< 50$) depuis M6 $\rightarrow$ \textbf{Suppression virale soutenue}. Risque transmission sexuelle nul (I=I). Prévention résistance.
        \item \textbf{Immunologique :} Reconstitution CD4 de 280 $\rightarrow$ 580 (gain +300). Sortie du stade Sida immunologique.
        \item \textbf{Clinique :} Patiente asymptomatique, grossesse menée à terme (enfant séronégatif si PMTCT respecté).
    \end{itemize}
    \textbf{Recommandation :} Maintenir observance, surveillance CV/CD4 tous 6 mois, dépistage IST, contraception/adaptation ARV si nouveau projet grossesse (DTG compatible).
\end{enumerate}
\end{exoresolu}

\begin{methode}[Interpréter un résultat de dépistage (Algorithme national Cameroun) et calculer le risque résiduel]
\textbf{Objectif :} Expliquer la démarche diagnostique (tests rapides, confirmation) et la notion de fenêtre sérologique.
\begin{enumerate}
    \item \textbf{Connaître l'algorithme national (stratégie 3 tests) :}
    \begin{itemize}
        \item \textbf{Test 1 (Dépistage) :} Test Rapide (TR) haute sensibilité (ex: Determine HIV-1/2, SD Bioline). Si \textbf{Négatif} $\rightarrow$ \textbf{Séronégatif} (sauf fenêtre).
        \item Si \textbf{Positif} $\rightarrow$ \textbf{Test 2 (Confirmation 1) :} TR haute spécificité (ex: Uni-Gold, Stat-Pak). Différent principe/antigènes que Test 1.
        \item Si Test 2 \textbf{Positif} $\rightarrow$ \textbf{Séropositif confirmé} (2 tests concordants positifs).
        \item Si Test 2 \textbf{Négatif} (Discordance) $\rightarrow$ \textbf{Test 3 (Arbitrage) :} TR ou ELISA (ex: Bioline, ELISA). Si Test 3 Positif $\rightarrow$ Positif. Si Test 3 Négatif $\rightarrow$ Négatif (ou indéterminé $\rightarrow$ nouveau prélèvement J+14).
    \end{itemize}
    \item \textbf{Gérer la « Fenêtre sérologique » (Période d'éclipse) :}
    \begin{itemize}
        \item Période post-contamination où le test est \textbf{négatif} malgré l'infection (pas encore d'anticorps détectables).
        \item Durée : \textbf{4 à 6 semaines} (tests 3$^{\text{e}}$/4$^{\text{e}}$ gén. Ag/Ab) à \textbf{3 mois} (tests rapides anticorps seuls).
        \item \textbf{Conduite à tenir :} Test négatif $\rightarrow$ \textbf{Recommander retrôtest à 6 semaines (test Ag/Ab) ou 3 mois (test rapide)}. Précaution prévention pendant ce délai.
    \end{itemize}
    \item \textbf{Calculer la Valeur Prédictive Positive (VPP) / Négative (VPN) :}
    \begin{align*}
        \text{VPP} &= \frac{\text{Vrais Positifs}}{\text{Vrais Positifs} + \text{Faux Positifs}} = \frac{\text{Se} \times \text{Prévalence}}{\text{Se} \times \text{Prévalence} + (1-\text{Sp}) \times (1-\text{Prévalence})} \\
        \text{VPN} &= \frac{\text{Vrais Négatifs}}{\text{Vrais Négatifs} + \text{Faux Négatifs}} = \frac{\text{Sp} \times (1-\text{Prévalence})}{\text{Sp} \times (1-\text{Prévalence}) + (1-\text{Se}) \times \text{Prévalence}}
    \end{align*}
    \textit{Se = Sensibilité, Sp = Spécificité.}
    \textbf{Règle :} Faible prévalence $\rightarrow$ VPP baisse (plus de faux positifs relatifs) $\rightarrow$ \textbf{Confirmation impérative}. Haute prévalence $\rightarrow$ VPN baisse (risque faux négatifs) $\rightarrow$ Vigilance fenêtre.
\end{enumerate}
\textbf{Terminologie :} « Test positif » $\neq$ « Séropositif confirmé ». On dit « Réactif au test de dépistage », « Confirmé par test de confirmation ».
\end{methode}

\begin{exoresolu}[Algorithme de dépistage et fenêtre sérologique]
\textbf{Énoncé :} Lors d'une campagne de dépistage dans une entreprise de Douala (prévalence estimée VIH = 3\%), on utilise l'algorithme national : Test 1 (Determine, Se=100\%, Sp=99,5\%), Test 2 (Uni-Gold, Se=100\%, Sp=99,8\%), Test 3 (Arbitrage ELISA, Se=99,9\%, Sp=99,9\%).
\begin{enumerate}
    \item Un employé a un Test 1 \textbf{Positif} et Test 2 \textbf{Négatif}. Que faire ? Quel est le statut probable ?
    \item Calculer la \textbf{VPP du Test 1 seul} dans cette population (Prévalence 3\%). Interpréter le résultat : pourquoi un test positif isolé ne suffit pas pour annoncer la séropositivité ?
    \item Un autre employé a eu un rapport à risque il y a \textbf{10 jours}. Il demande un test rapide (anticorps seul). Le test est négatif. Que lui répondre ? Quel test proposer et quand ?
\end{enumerate}
\tcblower
\textbf{Solution détaillée :}
\begin{enumerate}
    \item \textbf{Cas Discordance (T1 + / T2 -) :} \textbf{Protocole :} Réaliser le \textbf{Test 3 (Arbitrage)}.
    \begin{itemize}
        \item Si Test 3 \textbf{Positif} $\rightarrow$ Statut : \textbf{Séropositif confirmé} (2 tests positifs sur 3 : T1 et T3).
        \item Si Test 3 \textbf{Négatif} $\rightarrow$ Statut : \textbf{Séronégatif} (2 tests négatifs sur 3 : T2 et T3). L'employé est séronégatif, le T1 était un \textbf{faux positif}.
    \end{itemize}
    \textit{Note :} Ne jamais annoncer la séropositivité sur un seul test ou sur une discordance non résolue.
    \item \textbf{Calcul VPP Test 1 (Se=1, Sp=0,995, Prév=0,03) :}
    \[
    \text{VPP} = \frac{1 \times 0,03}{1 \times 0,03 + (1-0,995) \times (1-0,03)} = \frac{0,03}{0,03 + 0,005 \times 0,97} = \frac{0,03}{0,03 + 0,00485} = \frac{0,03}{0,03485} \approx \textbf{0,861} \ (86,1\%)
    \]
    \textbf{Interprétation :} Sur 100 tests positifs au Test 1, \textbf{seulement 86 sont de vrais positifs}. \textbf{14 sont des faux positifs} (spécificité 99,5\% $\rightarrow$ 0,5\% de faux positifs sur les 97\% séronégatifs = 0,485\% de la pop totale).
    \textbf{Conclusion :} Annoncer « Vous êtes séropositif » après un seul Test 1 exposerait \textbf{14\% des gens à une fausse annonce dramatique}. D'où l'exigence absolue de \textbf{confirmation par un 2$^{\text{e}}$ test de principe différent} (Sp 99,8\% $\rightarrow$ VPP finale $>$ 99,9\%).
    \item \textbf{Rapport à risque il y a 10 jours (Test Rapide Ab négatif) :}
    \begin{itemize}
        \item \textbf{Réponse :} « Ce test négatif \textbf{ne permet pas d'exclure l'infection} car vous êtes dans la \textbf{fenêtre sérologique} (période d'éclipse). Les anticorps ne sont pas encore détectables par ce test (apparition généralement 3-4 semaines, jusqu'à 3 mois). »
        \item \textbf{Test à proposer :} \textbf{Test combiné Ag/Ab (4$^{\text{e}}$ génération, ex: Alere Determine HIV-1/2 Ag/Ab Combo ou ELISA Ag/Ab)}. Il détecte l'antigène p24 (apparaît $\approx$ 15-20 jours) + anticorps.
        \item \textbf{Délai :} Refaire test Ag/Ab à \textbf{6 semaines (42 jours)} post-exposition. Si négatif $\rightarrow$ Fiabilité $> 99\%$. Si encore angoisse, test rapide anticorps à \textbf{3 mois} (12 semaines) pour lever tout doute.
        \item \textbf{Prophylaxie :} Évaluer éligibilité \textbf{PEP (Post-Exposition Prophylaxis)} : délai 10 jours $>$ 48h $\rightarrow$ \textbf{PEP plus possible}. Conseils prévention (préservatifs) en attendant le retrôtest.
    \end{itemize}
\end{enumerate}
\end{exoresolu}

\begin{methode}[Argumenter en faveur de l'observance thérapeutique et expliquer les conséquences de l'échec]
\textbf{Objectif :} Construire un argumentaire structuré (médical, individuel, collectif) pour soutenir l'observance.
\begin{enumerate}
    \item \textbf{Définir l'observance :} Prise des médicaments \textbf{au bon moment, à la bonne dose, sans oubli, sans interruption}, dans le respect des conditions (jeun/repas, interactions). Objectif : \textbf{$>$ 95\% des doses prises} (pour les anciennes molécules ; DTG plus tolérant mais viser 100\%).
    \item \textbf{Argument 1 : Efficacité virologique individuelle (Mécanisme) :}
    \begin{itemize}
        \item ARV bloquent la réplication $\rightarrow$ CV indétectable.
        \item \textbf{Oubli / Arrêt} $\rightarrow$ Réplication virale reprends $\rightarrow$ \textbf{Mutation aléatoire} lors de chaque cycle de réplication (RT sans preuve de lecture).
        \item Pression de sélection : Mutants résistants aux ARV pris $\rightarrow$ \textbf{Échec virologique} (CV $> 1000$ copies/mL).
    \end{itemize}
    \item \textbf{Argument 2 : Conséquences cliniques pour le patient :}
    \begin{itemize}
        \item Perte des acquis immunitaires (chute CD4), survenue Infections Opportunistes (IO), décès.
        \item \textbf{Résistance croisée :} Résistance à une molécule $\rightarrow$ souvent résistance à toute la classe (ex: NNRTI : EFV/NVP $\rightarrow$ résistance croisée totale). Perte d'options thérapeutiques (passage ligne 2, 3 plus chères, plus d'effets secondaires, pilulier plus lourd).
    \end{itemize}
    \item \textbf{Argument 3 : Santé publique / Prévention (TasP / I=I) :}
    \begin{itemize}
        \item Observance $\rightarrow$ CV indétectable $\rightarrow$ \textbf{Intransmissibilité sexuelle (I=I)}. Protection du/des partenaire(s).
        \item Non-observance $\rightarrow$ CV détectable $\rightarrow$ Risque transmission + \textbf{Transmission de souches résistantes} (résistance primaire chez nouveau infecté $\rightarrow$ limite options d'entrée).
    \end{itemize}
    \item \textbf{Argument 4 : Aspects pratiques / Soutien :} Outils (boîte pilulier, alarmes, applications), groupe de pairs (GAP, associations PLHIV), prise en charge effets secondaires (ne pas les cacher au soignant), dépistage dépression/addictions.
\end{enumerate}
\textbf{Phrase choc :} « \textbf{Prendre ses ARV tous les jours, c'est se protéger soi-même, protéger l'autre, et préserver ses traitements futurs.} »
\end{methode}

\begin{exoresolu}[Analyse d'un cas d'échec thérapeutique et conseil]
\textbf{Énoncé :} \textbf{Jean}, 32 ans, suit une trithérapie 1$^{\text{re}}$ ligne (TDF + 3TC + EFV) depuis 4 ans. Ses 3 derniers résultats de charge virale (CV) sont :
\begin{center}
\begin{tabular}{|c|c|c|}\hline
Date & CV (copies/mL) & CD4 (cell/mm$^3$) \\ \hline
M-12 & $< 20$ (Indétectable) & 620 \\ \hline
M-6 & 1 200 & 580 \\ \hline
M0 (aujourd'hui) & 15 000 & 490 \\ \hline
\end{tabular}
\end{center}
Jean avoue avoir « souvent oublié » ses comprimés ces 6 derniers mois à cause de son nouveau travail de nuit. Il ne présente pas de signe clinique d'IO.
\begin{enumerate}
    \item Qualifier l'évolution biologique (M-12 $\rightarrow$ M0). Parler d'\textbf{échec virologique} ? Justifier selon critères OMS.
    \item Expliquer le \textbf{mécanisme moléculaire} qui a conduit à cette situation (lien observance / réplication / mutation / sélection).
    \item La résistance au EFV (NNRTI) est suspectée. Pourquoi dit-on qu'il y a \textbf{résistance croisée} au sein des NNRTI ? Quelle conséquence sur la 2$^{\text{e}}$ ligne ?
    \item Proposer la \textbf{conduite à tenir} immédiate (conseil, examens, changement de traitement) selon recommandations OMS/Cameroun.
\end{enumerate}
\tcblower
\textbf{Solution détaillée :}
\begin{enumerate}
    \item \textbf{Évolution :} \textbf{Rebond virologique confirmé / Échec virologique.}
    \begin{itemize}
        \item M-12 : Succès (CV $< 20$).
        \item M-6 : \textbf{Blip / Bascule} : CV = 1 200 $>$ 1 000 copies/mL (seuil OMS échec). CD4 stables.
        \item M0 : \textbf{Confirmation échec} : CV = 15 000 (rebond $> 1$ log par rapport M-6), CD4 commencent à baisser (490).
    \end{itemize}
    \textbf{Critère OMS :} CV $> 1000$ copies/mL sur \textbf{deux mesures consécutives} (M-6 et M0) avec observance sous-optimale avérée $\rightarrow$ \textbf{Échec virologique confirmé}. Il ne s'agit pas d'un « blip » isolé (blip = CV faible transitoire $< 1000$ puis retour indétectable).
    \item \textbf{Mécanisme moléculaire :}
    \begin{enumerate}
        \item Oubli doses $\rightarrow$ Concentration plasmatique EFV/TDF/3TC chute sous CMI (Concentration Minimale Inhibitrice).
        \item Réplication virale reprend (cycle infection CD4 actif).
        \item \textbf{Erreur de la Transcriptase Inverse (RT) :} Pas d'activité exonucléase (proofreading) $\rightarrow$ Taux mutation $\approx 3 \times 10^{-5}$/base/cycle. Génome 10 kb $\rightarrow$ \textbf{Chaque virion =1 mutation en moyenne}.
                  \item \textbf{Sélection :} Mutants porteurs de résistance à l'EFV (ex: K103N, Y181C) ou 3TC (M184V) ont un \textbf{avantage sélectif} en présence du drug. Ils deviennent majoritaires (quasi-espèces).
                  \item Population virale $\rightarrow$ \textbf{Résistante aux ARV en cours}.
    \end{enumerate}
    \item \textbf{Résistance croisée NNRTI :}
    \begin{itemize}
        \item EFV (Efavirenz) et NVP (Nevirapine) sont des NNRTI de 1$^{\text{re}}$ génération. Ils se lient au même site allostérique de la RT.
        \item Mutations majeures (ex: \textbf{K103N}, Y181C, G190A) confèrent une \textbf{résistance de haut niveau à TOUS les NNRTI de 1$^{\text{e}}$ génération} (EFV, NVP).
        \item \textbf{Conséquence 2$^{\text{e}}$ ligne :} \textbf{Impossible de réutiliser un NNRTI (EFV/NVP)}. La 2$^{\text{e}}$ ligne OMS standard repose sur un \textbf{IP (Inhibiteur de Protéase) boosté au Ritonavir (ATV/r ou LPV/r) ou DTG (Inhibiteur Intégrase)} + 2 INRTI (souvent AZT + 3TC ou TDF + 3TC si M184V seul). Au Cameroun : \textbf{ATV/r (ou LPV/r) + AZT + 3TC} est la 2$^{\text{e}}$ ligne de référence si échec sur TDF/3TC/EFV.
    \end{itemize}
    \item \textbf{Conduite à tenir immédiate :}
    \begin{enumerate}
        \item \textbf{Conseil observance renforcé (Adherence Counseling) :} 3 séances rapprochées (J0, J14, J28) pour identifier barrières (travail nuit $\rightarrow$ adapter heure prise, boîte pilulier, alarme téléphone). \textbf{Ne PAS changer traitement avant d'avoir optimisé l'observance} (sauf CV très haute $> 100\,000$ ou signes cliniques).
        \item \textbf{Bilan pré-2$^{\text{e}}$ ligne :} Génotypage VIH (recherche mutations résistance) si disponible/accessible. Sinon, passage empirique 2$^{\text{e}}$ ligne après 3 mois d'observance optimisée si CV reste $> 1000$.
        \item \textbf{Changement traitement (si échec confirmé après ré-observance) :} Arrêt EFV (longue demi-vie $\rightarrow$ risque mono-thérapie « queue » EFV $\rightarrow$ donner TDF/3TC seuls 7-14 jours après arrêt EFV ou switch direct vers 2$^{\text{e}}$ ligne IP-based). \textbf{Nouveau régime : ATV/r (300/100 mg 1/j) + AZT (300 mg 2/j) + 3TC (150 mg 2/j)}. Surveillance hématologique (AZT : anémie/neutropénie), lipidique/rénale (ATV/r).
        \item \textbf{Suivi :} CV à 3 mois post-changement (objectif $< 50$).
    \end{enumerate}
\end{enumerate}
\textbf{Conclusion :} L'échec est évitable par l'observance. Une fois installé, il limite les options. L'éducation thérapeutique continue est le pilier de la pérennité de la 1$^{\text{re}}$ ligne.
\end{exoresolu}

\begin{methode}[Argumenter pour le dépistage volontaire (DV) et l'auto-dépistage]
\textbf{Objectif :} Structurer un plaidoyer (oral ou écrit) intégrant les bénéfices individuels et collectifs, et lever les freins.
\begin{enumerate}
    \item \textbf{Accroche :} Donnée épidémiologique locale (ex: « Au Cameroun, 2,7\% des adultes 15-49 ans vivent avec le VIH [EDS 2018], mais \textbf{1 sur 5 ignore son statut} »).
    \item \textbf{Bénéfice Individuel (Pour soi) :}
    \begin{itemize}
        \item \textbf{Négatif :} Rassurance, accès prévention (PrEP, préservatifs, conseil), tranquillité d'esprit.
        \item \textbf{Positif :} \textbf{Entrée précoce dans la prise en charge} $\rightarrow$ Démarrage ARV immédiat (Test \& Treat) $\rightarrow$ Espérance de vie \textbf{quasi normale}, qualité de vie préservée, évitement Sida/IO, projet parental sécurisé (PMTCT).
    \end{itemize}
    \item \textbf{Bénéfice Collectif (Pour les autres / Santé Publique) :}
    \begin{itemize}
        \item \textbf{TasP / I=I :} Personne traitée + CV indétectable = \textbf{Zéro transmission sexuelle}. « Undetectable = Untransmittable ».
        \item Casser la chaîne de transmission $\rightarrow$ Contrôle de l'épidémie (Objectifs 95-95-95 ONUSIDA : 95\% diagnostiqués, 95\% traités, 95\% supprimés).
        \item Protection du conjoint (sérodiscordance), de l'enfant (PMTCT).
    \end{itemize}
    \item \textbf{Levier : Auto-dépistage (VIH Auto-Test) :}
    \begin{itemize}
        \item Kit test rapide salivaire/sanguin (ex: OraQuick, Atomo) utilisable \textbf{seul, chez soi, confidentiel}.
        \item \textbf{Rôle :} Atteindre les « non-testés » (hommes, jeunes, populations clés, zones rurales).
        \item \textbf{Message clé :} « Test positif $\rightarrow$ \textbf{OBLIGATION} confirmation en structure de santé (CDV/Hôpital). Test négatif $\rightarrow$ Recommandation retrôtest si risque récent (fenêtre). »
    \end{itemize}
    \item \textbf{Lever les freins (Réponses aux objections) :}
    \begin{itemize}
        \item « J'ai peur du résultat » $\rightarrow$ « Savoir permet d'agir. Ne pas savoir, c'est risquer de découvrir trop tard (Sida). »
        \item « C'est stigmatisant » $\rightarrow$ « Confidentialité absolue (secret médical). Groupes de soutien (GAP) existent. »
        \item « Je n'ai pas de risque » $\rightarrow$ « Tout le monde a un statut. Le test est la seule façon de le connaître. Routinisation (dépistage proposé à toute consultation). »
    \end{itemize}
\end{enumerate}
\end{methode}

\begin{exoresolu}[Rédaction d'un plaidoyer pour le dépistage]
\textbf{Énoncé :} Vous êtes délégué(e) santé dans votre lycée à Bafoussam. Le proviseur vous demande de préparer une courte intervention (3 minutes) lors de la rentrée des classes pour inciter les élèves de Première à se faire dépister volontairement au CDV du district. Rédigez le texte de votre intervention en structurant vos arguments (Individuel, Collectif, Pratique) et en intégrant l'argument de l'auto-dépistage.
\tcblower
\textbf{Solution détaillée (Proposition de texte) :}
\textbf{« Chers camarades, chers aînés,}
\textbf{Aujourd'hui, au Cameroun, près de 500 000 personnes vivent avec le VIH. Le plus inquiétant ? \textbf{Près de 100 000 d'entre elles l'ignorent.} Elles sont en bonne apparence, vont en cours, travaillent, mais le virus détruit silencieusement leur système immunitaire.}
\textbf{Pourquoi se faire dépister MAINTENANT ?}
\textbf{1. Pour VOUS (Bénéfice individuel) :} Le test est \textbf{gratuit, confidentiel, rapide} (résultat en 20 min au CDV). Si vous êtes \textbf{négatif} : vous repartez rassuré(e), avec des préservatifs et des conseils pour le rester (PrEP si besoin). Si vous êtes \textbf{positif} : \textbf{ce n'est pas une condamnation}. Aujourd'hui, on commence les ARV \textbf{le jour même} (Test \& Treat). En les prenant bien, vous vivez \textbf{aussi longtemps et aussi bien} que tout le monde. Vous pourrez avoir des enfants séronégatifs, vous marier, réaliser vos rêves.
\textbf{2. Pour LES AUTRES (Bénéfice collectif) :} Une personne traitée avec une charge virale indétectable \textbf{NE TRANSMET PAS le virus sexuellement}. C'est la science : \textbf{I = I (Indétectable = Intransmissible)}. En vous testant et en vous soignant si besoin, vous protégez votre copain/copine, votre futur conjoint, votre futur bébé. Vous cassez la chaîne de l'épidémie.
\textbf{3. C'est FACILE et DISCRET (Auto-dépistage) :} Vous n'osez pas venir au CDV ? Il existe l'\textbf{auto-test VIH} (salive ou goutte de sang). Vous l'achetez en pharmacie ou on vous le donne gratuitement dans certaines associations. Vous le faites \textbf{chez vous, seul(e), en 15 min}. \textbf{Attention :} Si l'auto-test est positif, \textbf{COURREZ au CDV pour la confirmation} (c'est la règle : 2 tests positifs = séropositif). S'il est négatif mais que vous avez pris un risque récent (moins de 6 semaines), refaites-le plus tard (fenêtre sérologique).
\textbf{Camarades, le VIH n'a pas de visage. Il n'y a que VOTRE statut qui compte. Connaître son statut, c'est du POUVOIR. Le pouvoir de choisir sa vie, sa santé, son avenir.}
\textbf{Rendez-vous au CDV du District de Bafoussam, ou demandez-moi un auto-test. \textbf{Testez-vous. Protégez-vous. Aimez-vous en sécurité.} Merci. »}
\textbf{Analyse de la réponse :} Structure claire (Donnée locale $\rightarrow$ 3 piliers d'argumentation $\rightarrow$ Outil pratique $\rightarrow$ Conclusion mobilisatrice). Ton pair-à-pair, non moralisateur, scientifique (I=I, Fenêtre, Confirmation), contextualisé (Bafoussam, CDV, Gratuit). Respecte le temps (environ 3 min à l'oral).
\end{exoresolu}

\begin{attention}[Les 5 erreurs qui coûtent des points au Probatoire]
\begin{enumerate}
    \item \textbf{Confondre « Test positif » et « Séropositivité confirmée ».} \textbf{Erreur :} « Le patient est séropositif car son test Determine est positif. » \textbf{Correct :} « Le patient est \textbf{réactif au test de dépistage (T1)}. La séropositivité n'est confirmée qu'après \textbf{concordance positive d'un test de confirmation (T2)} (ou arbitrage T3 en cas de discordance). » $\rightarrow$ \textit{Notion clé : Algorithme national, VPP, Spécificité.}
    \item \textbf{Oublier la « Fenêtre sérologique » dans l'interprétation d'un test négatif récent.} \textbf{Erreur :} « Le test est négatif, il n'est pas infecté. » \textbf{Correct :} « Test négatif \textbf{mais exposition il y a 10 jours} $\rightarrow$ \textbf{Fenêtre sérologique}. Risque faux négatif. Recommander \textbf{retrôtest à 6 semaines (test Ag/Ab) ou 3 mois (test rapide Ab)}. » $\rightarrow$ \textit{Mots-clés : Période d'éclipse, p24, Ag/Ab, retrôtest.}
    \item \textbf{Dire que la PrEP ou le TasP protègent des IST ou de la grossesse.} \textbf{Erreur :} « Comme il prend la PrEP / sa charge virale est indétectable, il n'a pas besoin de préservatif. » \textbf{Correct :} « La PrEP et le TasP (I=I) protègent \textbf{UNIQUEMENT contre le VIH}. Ils ne protègent \textbf{PAS} contre les IST (syphilis, chlamydia, HPV, herpès, hépatite B) ni contre la grossesse non désirée. \textbf{Le préservatif reste la seule triple protection (VIH + IST + Grossesse).} » $\rightarrow$ \textit{Prévention combinée.}
    \item \textbf{Inverser l'évolution des CD4 et de la CV en phase asymptomatique (non traitée).} \textbf{Erreur :} « En phase asymptomatique, la charge virale baisse et les CD4 montent. » \textbf{Correct :} « En phase asymptomatique \textbf{sans traitement}, la charge virale se stabilise au \textbf{point de consigne} (plateau) et les CD4 \textbf{baissent lentement et régulièrement} ($\approx$ 50-80/an). Seule l'ARV fait baisser la CV et monter les CD4. » $\rightarrow$ \textit{Dynamique naturelle vs Dynamique sous ARV.}
    \item \textbf{Ne pas lier l'observance à la résistance (mécanisme RT) et à la transmission de souches résistantes.} \textbf{Erreur :} « S'il oublie ses médicaments, le virus revient et il tombe malade. » \textbf{Correct :} « L'inobservance $\rightarrow$ réplication virale non contrôlée $\rightarrow$ \textbf{mutations aléatoires par la Transcriptase Inverse (pas de proofreading)} $\rightarrow$ \textbf{sélection de mutants résistants} aux ARV pris $\rightarrow$ \textbf{échec virologique} (CV $>$ 1000) $\rightarrow$ \textbf{résistance croisée} (perte options 1$^{\text{re}}$ ligne) $\rightarrow$ risque \textbf{transmission de souche résistante} (résistance primaire chez nouveau infecté). » $\rightarrow$ \textit{Mécanisme moléculaire + Conséquence individuelle + Conséquence santé publique.}
\end{enumerate}
\end{attention}

\newpage

\coursec{Exercices}

\courssub{Je vérifie mes connaissances}

\begin{exercice}[QCM : Le VIH et son cycle]
Pour chaque question, une seule réponse est exacte. Recopie le numéro et la lettre correspondante.
\begin{enumerate}
    \item Le VIH est un rétrovirus. Son matériel génétique est :
    \begin{itemize}
        \item[A.] de l'ADN double brin
        \item[B.] de l'ARN simple brin
        \item[C.] de l'ARN double brin
        \item[D.] de l'ADN simple brin
    \end{itemize}
    \item L'enzyme qui permet la transcription inverse du génome viral en ADN est :
    \begin{itemize}
        \item[A.] l'ADN polymérase
        \item[B.] l'ARN polymérase
        \item[C.] la \cle{transcriptase inverse}
        \item[D.] l'intégrase
    \end{itemize}
    \item La fixation du VIH sur le lymphocyte CD4 se fait via l'interaction entre la glycoprotéine virale gp120 et le récepteur cellulaire :
    \begin{itemize}
        \item[A.] CCR5 uniquement
        \item[B.] CXCR4 uniquement
        \item[C.] CD4 et un co-récepteur (CCR5 ou CXCR4)
        \item[D.] le récepteur au CD8
    \end{itemize}
    \item Lors de la phase asymptomatique de l'infection à VIH :
    \begin{itemize}
        \item[A.] la charge virale est nulle et les CD4 sont stables
        \item[B.] le virus est éradiqué par le système immunitaire
        \item[C.] la réplication virale se poursuit activement dans les organes lymphoïdes et la charge virale est stable (point de consigne)
        \item[D.] les lymphocytes CD4 augmentent progressivement
    \end{itemize}
    \item Les antirétroviraux (ARV) ciblent différentes étapes du cycle viral. L'inhibiteur de l'intégrase bloque :
    \begin{itemize}
        \item[A.] la fixation du virus sur la cellule
        \item[B.] la transcription inverse
        \item[C.] l'intégration de l'ADN proviral dans le génome de la cellule hôte
        \item[D.] la maturation des nouvelles particules virales
    \end{itemize}
\end{enumerate}
\end{exercice}

\begin{exercice}[Vrai ou Faux ? Justifie chaque réponse]
Indique si l'affirmation est vraie (V) ou fausse (F) et justifie brièvement en mobilisant tes connaissances.
\begin{enumerate}
    \item On peut contracter le VIH en partageant un repas avec une personne séropositive.
    \item Le lait maternel est un liquide biologique pouvant transmettre le VIH.
    \item La séroconversion correspond au moment où la charge virale devient indétectable.
    \item Un patient sous traitement antirétroviral efficace (charge virale indétectable) ne transmet plus le VIH par voie sexuelle (concept I=I).
    \item Les lymphocytes CD8 cytotoxiques sont les principales cellules cibles du VIH.
\end{enumerate}
\end{exercice}

\begin{exercice}[Définitions et vocabulaire]
Donne une définition précise et scientifique pour chacun des termes suivants :
\begin{enumerate}
    \item \cle{VIH} (Virus de l'Immunodéficience Humaine)
    \item \cle{Lymphocyte T CD4} (ou lymphocyte T helper)
    \item \cle{Séroconversion}
    \item \cle{Charge virale}
    \item \cle{Trithérapie} (ou traitement antirétroviral combiné)
\end{enumerate}
\end{exercice}

\begin{exercice}[Calcul direct : Chute des CD4]
Au moment de l'infection primaire, un patient présente un taux de lymphocytes CD4 de \num{800} cellules/mm\textsuperscript{3}. Au stade Sida déclaré (sans traitement), ce taux est tombé à \num{150} cellules/mm\textsuperscript{3}.
\begin{enumerate}
    \item Calcule le pourcentage de chute du taux de CD4 entre ces deux stades.
    \item Rappelle le seuil de CD4 (en cellules/mm\textsuperscript{3}) en dessous duquel on parle de Sida déclaré selon la définition immunologique de l'OMS.
\end{enumerate}
\end{exercice}

\courssub{Je m'entraîne}

\begin{exercice}[Classification des risques de transmission]
Classe chacune des situations suivantes selon le risque de transmission du VIH : \textbf{Risque élevé}, \textbf{Risque faible/négligeable}, \textbf{Risque nul}. Justifie chaque classement en précisant le liquide biologique impliqué et sa concentration en virus.
\begin{enumerate}
    \item Rapport sexuel vaginal non protégé avec une personne séropositive non traitée.
    \item Baiser sur la joue (baiser « social »).
    \item Partage d'une seringue usagée entre usagers de drogues injectables.
    \item Allaitement maternel par une mère séropositive non traitée.
    \item Piqûre d'aiguille accidentelle chez un infirmier ayant prélevé du sang sur un patient séropositif.
    \item Partage d'une brosse à dents avec un colocataire séropositif (saignements gingivaux possibles).
    \item Morsure humaine sans effusion de sang.
    \item Transfusion sanguine non dépistée (contexte d'urgence sans test rapide).
    \item Rapport sexuel oral (fellation) avec éjaculation dans la bouche, partenaire séropositif sous trithérapie efficace (charge virale indétectable).
    \item Utilisation des mêmes toilettes qu'une personne séropositive.
\end{enumerate}
\end{exercice}

\begin{exercice}[Interprétation de courbes : Charge virale et CD4 sur 10 ans]
\begin{popfigure}
\begin{tikzpicture}
\begin{axis}[
    width=\linewidth,
    height=8cm,
    xlabel={Temps (années après infection)},
    ylabel={Valeurs},
    ymin=0, ymax=1200,
    xmin=0, xmax=12,
    xtick={0,2,4,6,8,10,12},
    ytick={0,200,400,600,800,1000},
    legend pos=north east,
    legend cell align=left,
    grid=major,
]
% CD4 curve (decreasing)
\addplot[popcurve, popblue, very thick, domain=0:12, samples=100] 
    {800 - 50*x - 10*sin(deg(x*3))}; % Approximation visuelle
\addlegendentry{Lymphocytes CD4 (cellules/mm\textsuperscript{3})}

% Viral Load curve (log scale representation on linear axis for simplicity, or use two axes)
% We'll plot Log10(VL) scaled to fit, or just a schematic curve.
% Let's plot a schematic Viral Load (arbitrary units 0-100 representing log copies)
\addplot[popcurve, poporange, very thick, dashed, domain=0:12, samples=100] 
    {50 + 400*exp(-(x-0.5)^2/0.5) + 50*x + 10*sin(deg(x*2))}; 
\addlegendentry{Charge virale (échelle arbitraire croissante $\sim$ log copies/mL)}
\end{axis}
\end{tikzpicture}
\legende{Évolution schématique de la charge virale (courbe orange pointillée) et du taux de lymphocytes CD4 (courbe bleue) au cours d'une infection à VIH non traitée sur 12 ans.}
\end{popfigure}

À partir du graphique ci-dessus :
\begin{enumerate}
    \item Identifie et nomme les trois grandes phases de l'infection à VIH (A, B, C) en te basant sur l'évolution des deux courbes.
    \item Décris l'évolution de la charge virale et des CD4 lors de la \textbf{phase A} (infection primaire). Quel événement immunologique majeur se produit à la fin de cette phase ?
    \item Explique la relation inverse observée entre la charge virale et le taux de CD4 tout au long de l'infection.
    \item Pourquoi la phase B est-elle appelée « phase asymptomatique » alors que la réplication virale est active ?
\end{enumerate}
\end{exercice}

\begin{exercice}[Cycle d'infection et cibles thérapeutiques]
Le schéma ci-dessous représente le cycle d'infection d'un lymphocyte CD4 par le VIH. Certaines étapes sont numérotées de 1 à 7.
\begin{popfigure}
\begin{tikzpicture}[node distance=1.2cm, >=Stealth, font=\small]
% Nodes
\node[draw, rounded corners, fill=popblueL, minimum width=3cm, minimum height=1.2cm, align=center] (virus){VIRION VIH\\ (ARN, RT, Intégrase, Protéase)};
\node[draw, rounded corners, fill=popgreenL, minimum width=3cm, minimum height=1.2cm, right=of virus, align=center] (cell){LYMPHOCYTE CD4\\ (Noyau, Ribosomes, Appareil de Golgi)};
\node[draw, dashed, rounded corners, fit=(virus) (cell), inner sep=5pt, label=above:\textbf{Cycle d'infection}] (frame) {};

% Steps - simplified text representation for LaTeX compilation without complex drawing
% We will use a textual description with numbered steps for the exercise.
\end{tikzpicture}
\legende{Schéma simplifié du cycle d'infection du VIH (représentation textuelle ci-dessous).}
\end{popfigure}

\textbf{Étapes du cycle (à ordonner et légender) :}
1. Fixation (gp120 sur CD4 + co-récepteur) $\rightarrow$ 2. Fusion et entrée $\rightarrow$ 3. \textbf{Transcription inverse} (ARN $\rightarrow$ ADN) $\rightarrow$ 4. \textbf{Intégration} (ADN proviral dans chromosome hôte) $\rightarrow$ 5. Transcription/Traduction (protéines virales) $\rightarrow$ 6. Assemblage $\rightarrow$ 7. \textbf{Budding/Maturation} (clivage par Protéase).

\begin{enumerate}
    \item Replace les 7 étapes dans l'ordre chronologique (ci-dessus elles le sont déjà, mais l'exercice consiste à les associer aux numéros sur un schéma vierge en examen). Ici, \textbf{identifie le nom de l'enzyme virale clé} pour les étapes 3, 4 et 7.
    \item Deux classes majeures d'antirétroviraux (ARV) agissent sur des étapes distinctes :
    \begin{itemize}
        \item Les \textbf{Inhibiteurs de la Transcriptase Inverse (ITI)} : sur quelle étape agissent-ils ? Donne un exemple de molécule (ex: Ténofovir, Lamivudine).
        \item Les \textbf{Inhibiteurs de la Protéase (IP)} : sur quelle étape agissent-ils ? Conséquence sur les virions produits ?
    \end{itemize}
    \item Cite une troisième classe d'ARV (autre que ITI et IP) et son étape d'action.
\end{enumerate}
\end{exercice}

\begin{exercice}[Cas pratique : Prévention de la Transmission Mère-Enfant (PTME)]
Au Cameroun, le programme national de PTME vise à réduire le taux de transmission mère-enfant à moins de 5\%.
\begin{enumerate}
    \item Cite les trois moments possibles de transmission du VIH de la mère à l'enfant.
    \item Une femme enceinte séropositive débute une trithérapie au 2\textsuperscript{e} trimestre. Son enfant reçoit une prophylaxie à la naissance (sirop de Névirine ou Zidovudine) pendant 6 semaines. L'allaitement est exclusif sous couverture ARV maternel pendant 6 mois.
    \begin{itemize}
        \item Explique le rôle de la trithérapie maternelle sur la charge virale plasmatique et le risque de transmission \emph{in utero} et \emph{péri-partum}.
        \item Pourquoi la prophylaxie du nouveau-né est-elle nécessaire même si la mère est traitée ?
        \item L'allaitement maternel exclusif sous ARV est-il recommandé au Cameroun ? Justifie par rapport aux risques de malnutrition/maladies diarrhéiques si lait artificiel (eau non potable).
    \end{itemize}
\end{enumerate}
\end{exercice}

\courssub{Je me prépare au Probatoire}

\begin{exercice}[Sujet type Probatoire : Dynamique de l'infection à VIH non traitée]
On étudie l'évolution de l'infection à VIH chez un patient non traité, suivi pendant 12 ans. Le document 1 présente l'évolution de la charge virale (copies/mL) et du taux de lymphocytes CD4 (cellules/mm\textsuperscript{3}).

\begin{popfigure}
\begin{tikzpicture}
\begin{axis}[
    name=axcd4,
    width=\linewidth, height=9cm,
    xlabel={Temps (années)}, ylabel={CD4 (cellules/mm\textsuperscript{3})},
    xmin=0, xmax=12, ymin=0, ymax=1000,
    xtick={0,1,2,3,4,5,6,7,8,9,10,11,12},
    ytick={0,200,400,600,800,1000},
    grid=major, legend pos=north east
]
% CD4 Data points (simulated)
\addplot[popblue, very thick, mark=*, mark size=2] coordinates {
    (0, 800) (0.5, 400) (1, 550) (2, 600) (3, 580) (4, 550) (5, 500) (6, 450) (7, 380) (8, 300) (9, 220) (10, 150) (11, 100) (12, 50)
};
\addlegendentry{CD4}
\end{axis}
% Second axis (non imbriqué, superposé) for Viral Load (Log scale)
\begin{axis}[
    at={(axcd4.south west)}, anchor=south west,
    width=\linewidth, height=9cm,
    xmin=0, xmax=12, ymin=1, ymax=7, % Log10 scale
    ytick={1,2,3,4,5,6,7},
    yticklabels={$10^1$, $10^2$, $10^3$, $10^4$, $10^5$, $10^6$, $10^7$},
    ylabel={Charge virale (copies/mL)},
    axis y line*=right, axis x line=none,
    legend pos=south east,
]
\addplot[poporange, very thick, mark=square*, mark size=2, dashed] coordinates {
    (0, 1) (0.2, 6.5) (0.5, 5.5) (1, 4.5) (2, 4.2) (3, 4.3) (4, 4.4) (5, 4.5) (6, 4.6) (7, 4.8) (8, 5.0) (9, 5.5) (10, 6.0) (11, 6.5) (12, 6.8)
};
\addlegendentry{Charge virale (log)}
\end{axis}
\end{tikzpicture}
\legende{Document 1 : Évolution de la charge virale (échelle logarithmique, courbe orange pointillée) et du taux de lymphocytes CD4 (courbe bleue) chez un patient non traité pendant 12 ans.}
\end{popfigure}

\textbf{Questions :}
\begin{enumerate}
    \item \textbf{Phase 1 (Année 0 à 1) :} Décris l'évolution simultanée de la charge virale et des CD4. Nomme cette phase. Quel mécanisme immunitaire explique la remontée transitoire des CD4 et la baisse de la charge virale vers le « point de consigne » ?
    \item \textbf{Phase 2 (Année 1 à 8 env.) :} Caractérise l'évolution des deux paramètres. Pourquoi cette phase est-elle dite « asymptomatique » cliniquement mais « dynamique » virologiquement ? Calcule la durée moyenne de cette phase.
    \item \textbf{Phase 3 (Année 8 à 12) :} Décris l'évolution. Nomme cette phase. Quel est le seuil de CD4 franchi vers l'année 9-10 qui signe le passage au Sida déclaré (définition immunologique) ?
    \item \textbf{Calcul :} Calcule le pourcentage de diminution du taux de CD4 entre le pic de l'infection primaire (an 0,5) et le stade Sida (an 10).
    \item \textbf{Synthèse :} Explique la relation inverse entre charge virale et taux de CD4 en détaillant les mécanismes directs (cytopathie virale) et indirects (activation immune chronique, apoptose des bystanders, destruction de l'architecture ganglionnaire).
\end{enumerate}
\end{exercice}

\begin{exercice}[Sujet type Probatoire : Le dépistage volontaire, arme de santé publique]
« Le dépistage volontaire du VIH est une arme majeure de santé publique pour contrôler l'épidémie. »
À partir de tes connaissances, discute cette affirmation en structurant ton argumentation autour des quatre piliers suivants :
\begin{enumerate}
    \item \textbf{La fenêtre sérologique et la séroconversion :} Définis ces deux notions. Pourquoi un test négatif ne garantit-il pas l'absence d'infection si l'exposition est récente ? Quelle conduite tenir ?
    \item \textbf{L'entrée précoce dans la prise en charge :} Explique en quoi le dépistage précoce (stade asymptomatique) améliore le pronostic individuel (reconstitution immunitaire, espérance de vie) et collectif (prévention).
    \item \textbf{La Prévention de la Transmission Mère-Enfant (PTME) :} En quoi le dépistage prénatal systématique (opt-out) est-il la condition \emph{sine qua non} de l'élimination de la transmission verticale ? Chiffres clés au Cameroun (taux de transmission sans/sans intervention).
    \item \textbf{L'observance thérapeutique et le concept I=I (Indétectable = Intransmissible) :} Explique le lien entre observance stricte, charge virale indétectable durable et absence de transmission sexuelle. Pourquoi le dépistage est-il la porte d'entrée de cette cascade de soins ?
\end{enumerate}
Ta réponse sera rédigée, structurée en paragraphes argumentés et illustrée par des exemples concrets (ex: campagne « Holiday Without HIV » au Cameroun, index testing).
\end{exercice}

\begin{exercice}[Sujet type Probatoire : Rupture d'observance et risque de résistance]
Mme M., 34 ans, suit une trithérapie (Ténofovir + Lamivudine + Dolutégravir) depuis 4 ans. Sa charge virale est indétectable ($< 20$ copies/mL) et ses CD4 à 650/mm\textsuperscript{3}. Pour des raisons socio-économiques (rupture de stock à la pharmacie de district, frais de transport), elle interrompt son traitement pendant 3 mois. Elle reprend ensuite régulièrement. Un contrôle à 1 mois de la reprise montre : Charge virale = 45 000 copies/mL ; CD4 = 420/mm\textsuperscript{3}.

\begin{enumerate}
    \item \textbf{Rebond viral :} Explique le mécanisme physiopathologique du rebond viral massif observé après l'arrêt du traitement (réservoirs viraux, réactivation de la transcription du provirus intégré).
    \item \textbf{Chute des CD4 :} Interprète la baisse des CD4 de 650 à 420 en 3 mois d'arrêt puis 1 mois de reprise. Relie-la à la réplication virale non contrôlée.
    \item \textbf{Résistance :} Le Dolutégravir (Inhibiteur d'Intégrase de 2\textsuperscript{e} génération) a un \textbf{barrière génétique à la résistance élevée}. Cependant, la Lamivudine (3TC) a une barrière faible (mutation M184V rapide).
    \begin{itemize}
        \item Pourquoi l'arrêt \emph{simultané} des trois molécules (arrêt brutal) limite-t-il théoriquement le risque de sélection de mutants résistants par rapport à un arrêt progressif ou une monothérapie déguisée ?
        \item Malgré tout, pourquoi le rebond viral à 45 000 copies/mL sous traitement repris expose-t-il au risque d'émergence de résistances si l'observance n'est pas parfaite par la suite ?
    \end{itemize}
    \item \textbf{Argumentaire d'éducation thérapeutique :} En tant qu'agent de santé, rédige un message clair et bienveillant pour Mme M. afin de l'aider à comprendre l'importance cruciale de l'observance \textbf{à vie, sans omission}, en mobilisant les notions de \emph{charge virale indétectable}, \emph{réservoirs}, \emph{résistance} et \emph{transmissibilité}. Propose une solution concrète pour éviter les ruptures de stock (ex: dispensation communautaire, multi-mois).
\end{enumerate}
\end{exercice}

\newpage

\coursec{Corrigés des exercices}

\coursec{Exercices corrigés et sujets type Probatoire — Lutte contre le VIH/Sida}

\begin{corrige}[Exercice 1 — QCM : Le VIH et son cycle]
\textbf{1 — B.} Le VIH est un \cle{rétrovirus} dont le génome est constitué de \textbf{deux molécules d'ARN simple brin} de sens positif, encapsidées dans la capside. C'est précisément parce que son génome est de l'ARN qu'il a besoin d'une transcriptase inverse.

\textbf{2 — C.} La \cle{transcriptase inverse} (ou \emph{reverse transcriptase}, RT) est l'enzyme virale qui synthétise un ADN complémentaire à partir de l'ARN viral. L'\cle{intégrase} (D) intervient plus tard, lors de l'intégration de l'ADN proviral.

\textbf{3 — C.} La glycoprotéine d'enveloppe \textbf{gp120} se fixe d'abord sur le récepteur \textbf{CD4}, ce qui provoque un changement de conformation permettant la liaison à un \textbf{co-récepteur} (CCR5 ou CXCR4). La fixation sur CD4 seul ou sur le co-récepteur seul est insuffisante : les deux interactions sont nécessaires à la fusion.

\textbf{4 — C.} Pendant la phase asymptomatique, le virus se réplique activement dans les organes lymphoïdes (ganglions) ; la charge virale plasmatique se stabilise à un niveau appelé \emph{point de consigne} (set point), tandis que les CD4 déclinent lentement. Le virus n'est jamais éradiqué (B faux) et les CD4 n'augmentent pas (D faux).

\textbf{5 — C.} Les \cle{inhibiteurs d'intégrase} (ex. Dolutégravir, Raltégravir) bloquent l'\textbf{intégration} de l'ADN proviral dans le chromosome de la cellule hôte. La maturation (D) est la cible des \cle{inhibiteurs de la protéase}.
\end{corrige}

\begin{corrige}[Exercice 2 — Vrai ou Faux ?]
\begin{enumerate}
    \item \textbf{FAUX.} La salive, la sueur et les larmes ne contiennent pas de quantités infectieuses de VIH. Seuls le \textbf{sang}, le \textbf{sperme}, les \textbf{sécrétions vaginales et rectales} et le \textbf{lait maternel} transmettent le virus. Partager un repas, un verre ou des couverts ne présente \emph{aucun risque}.
    \item \textbf{VRAI.} Le lait maternel contient des lymphocytes infectés et des particules virales libres : la transmission peut survenir pendant l'allaitement (transmission post-natale). D'où la prophylaxie ARV maternelle et infantile pendant toute la durée de l'allaitement au Cameroun.
    \item \textbf{FAUX.} La \cle{séroconversion} est l'\emph{apparition} d'anticorps anti-VIH détectables dans le sang (2 à 12 semaines après l'infection). Elle coïncide avec une \emph{baisse} de la charge virale (contrôle immunitaire partiel), mais celle-ci ne devient indétectable que sous traitement antirétroviral efficace.
    \item \textbf{VRAI.} Le concept \textbf{I = I} (Indétectable = Intransmissible), validé par les études PARTNER et HPTN 052, établit qu'une personne sous ARV avec une charge virale durablement indétectable ($< 200$ copies/mL, voire $< 50$) ne transmet pas le VIH par voie sexuelle. Condition : observance parfaite et contrôle régulier de la charge virale.
    \item \textbf{FAUX.} Les cibles principales du VIH sont les \textbf{lymphocytes T CD4} (helper), ainsi que les macrophages et les cellules dendritiques qui expriment CD4. Les CD8 cytotoxiques participent au contraire à la \emph{défense} anti-VIH en détruisant les cellules infectées.
\end{enumerate}
\end{corrige}

\begin{corrige}[Exercice 3 — Définitions et vocabulaire]
\begin{enumerate}
    \item \textbf{VIH} : rétrovirus (famille des \emph{Retroviridae}, genre \emph{Lentivirus}) à enveloppe, dont le génome est constitué de deux molécules d'ARN simple brin associées à la transcriptase inverse. Il infecte préférentiellement les cellules portant le récepteur CD4 (lymphocytes T4, macrophages) et provoque une immunodépression progressive.
    \item \textbf{Lymphocyte T CD4} : globule blanc du système immunitaire adaptatif portant le marqueur membranaire CD4. Véritable « chef d'orchestre » de la réponse immunitaire, il active les lymphocytes B (production d'anticorps) et les lymphocytes T CD8 cytotoxiques. Taux normal : 500 à \num{1200} cellules/mm\textsuperscript{3}.
    \item \textbf{Séroconversion} : passage d'un état séronégatif à un état séropositif, c'est-à-dire apparition dans le sang d'anticorps spécifiques anti-VIH détectables par les tests sérologiques (ELISA, tests rapides), en général 3 à 12 semaines après la contamination.
    \item \textbf{Charge virale} : quantité de VIH circulant dans le plasma, exprimée en nombre de copies d'ARN viral par millilitre de sang (copies/mL). Elle reflète l'intensité de la réplication virale et l'efficacité du traitement. Elle est dite « indétectable » en dessous du seuil de détection de la technique ($< 20$ à $50$ copies/mL).
    \item \textbf{Trithérapie} : traitement antirétroviral \emph{combiné} associant au moins trois molécules, selon le schéma standard \textbf{2 INRTI (inhibiteurs nucléos(t)idiques de la transcriptase inverse) + 1 molécule d'une autre classe} (INNTI, inhibiteur d'intégrase ou inhibiteur de la protéase boosté). L'association vise à bloquer simultanément plusieurs étapes du cycle viral, à obtenir une charge virale indétectable et à empêcher l'émergence de souches résistantes.
\end{enumerate}
\end{corrige}

\begin{corrige}[Exercice 4 — Calcul direct : chute des CD4]
\begin{enumerate}
    \item La chute absolue est :
    \[ 800 - 150 = 650 \text{ cellules/mm}^3 \]
    Le pourcentage de chute par rapport à la valeur initiale est :
    \[ \frac{650}{800} \times 100 = \frac{65\,000}{800} = \num{81,25}\,\% \]
    Le taux de CD4 a donc chuté de \textbf{\num{81,25} \%} entre l'infection primaire et le stade Sida déclaré.
    \item Selon la définition immunologique de l'OMS, on parle de \textbf{Sida déclaré} (stade d'immunodéficience sévère) lorsque le taux de CD4 passe \textbf{en dessous de 200 cellules/mm\textsuperscript{3}}. Ici, 150 cellules/mm\textsuperscript{3} : le patient est bien en dessous du seuil.
\end{enumerate}
\begin{attention}[Point de vigilance]
Le pourcentage de chute se calcule toujours par rapport à la valeur \emph{initiale} (800), et non à la valeur finale. Une erreur classique consiste à écrire $650/150 \times 100$.
\end{attention}
\end{corrige}

\begin{corrige}[Exercice 5 — Classification des risques de transmission]
\begin{enumerate}
    \item \textbf{Risque élevé.} Rapport vaginal non protégé : le sperme et les sécrétions vaginales/cervicales d'une personne non traitée ont une charge virale élevée ; les micro-lésions muqueuses facilitent l'entrée du virus. C'est le premier mode de transmission au Cameroun.
    \item \textbf{Risque nul.} Le baiser social met en jeu la peau intacte et éventuellement la salive, liquide non contaminant.
    \item \textbf{Risque élevé.} Partage de seringue : le sang résiduel, très concentré en virus, est injecté directement dans la circulation sanguine — voie la plus efficace de transmission.
    \item \textbf{Risque élevé.} Allaitement par une mère non traitée : le lait maternel contient des cellules infectées ; le risque cumulé de transmission post-natale atteint 10 à 15 \% en l'absence d'intervention.
    \item \textbf{Risque faible mais réel ($\approx \num{0,3}\,\%$).} Piqûre accidentelle avec aiguille souillée de sang : exposition sanguine directe mais volume faible. Conduite à tenir : traitement post-exposition (TPE) en urgence, dans les 48 à 72 h.
    \item \textbf{Risque faible/négligeable mais déconseillé.} La brosse à dents peut être souillée de sang gingival ; le virus survit mal à l'air libre, mais le risque théorique existe : on ne partage jamais les objets de toilette pouvant être tachés de sang.
    \item \textbf{Risque nul à négligeable.} Sans effusion de sang, seule la salive est en contact : pas de transmission documentée.
    \item \textbf{Risque très élevé ($\approx 90$--$95\,\%$).} Transfusion de sang non dépisté : inoculation massive de sang contaminé. D'où le dépistage systématique des poches de sang au Cameroun (CNTS).
    \item \textbf{Risque nul à négligeable (I = I).} Le partenaire ayant une charge virale durablement indétectable ne transmet pas le VIH par voie sexuelle. (En l'absence de traitement, la fellation avec éjaculation présenterait un risque faible mais réel.)
    \item \textbf{Risque nul.} Aucun liquide contaminant n'est échangé ; le VIH ne se transmet pas par les surfaces, l'eau ou l'air.
\end{enumerate}
\begin{aretenir}[À retenir]
Quatre liquides seulement transmettent le VIH : \textbf{sang, sperme, sécrétions vaginales/rectales, lait maternel}. Trois portes d'entrée : voie \textbf{sexuelle}, voie \textbf{sanguine}, transmission \textbf{mère-enfant}. Toute situation sans ces liquides ni ces voies est à risque nul.
\end{aretenir}
\end{corrige}

\begin{corrige}[Exercice 6 — Interprétation de courbes : charge virale et CD4 sur 10 ans]
\begin{enumerate}
    \item \textbf{Phase A (0 à $\approx$ 1 an) : infection primaire} (primo-infection) — pic de charge virale puis chute ; chute transitoire des CD4. \textbf{Phase B (1 à $\approx$ 8--10 ans) : phase asymptomatique} (latence clinique) — charge virale stabilisée au point de consigne, déclin lent des CD4. \textbf{Phase C (au-delà) : phase Sida} — remontée brutale de la charge virale et effondrement des CD4.
    \item En phase A, la charge virale explose (pic à $10^6$--$10^7$ copies/mL) car le virus se réplique sans opposition, puis elle diminue fortement ; les CD4 chutent puis remontent partiellement. À la fin de cette phase survient la \textbf{séroconversion} : apparition des anticorps anti-VIH et mise en place de la réponse immunitaire (lymphocytes CD8 cytotoxiques) qui contrôle partiellement la réplication et fixe le « point de consigne ».
    \item Le VIH infecte et détruit les CD4 : plus la réplication est intense (charge virale élevée), plus la destruction des CD4 est rapide. Réciproquement, quand la réponse immunitaire contient le virus (charge virale basse), les CD4 se reconstituent partiellement. D'où l'évolution en miroir des deux courbes : la charge virale est le \emph{moteur}, le taux de CD4 la \emph{conséquence} de l'infection.
    \item Elle est « asymptomatique » car le patient ne présente \emph{aucun signe clinique} : l'immunité résiduelle suffit à prévenir les infections. Mais virologiquement, des milliards de virions sont produits chaque jour dans les ganglions, et les CD4 déclinent régulièrement : c'est une \emph{latence clinique}, pas une latence virologique. Le patient asymptomatique est \textbf{contaminant}.
\end{enumerate}
\end{corrige}

\begin{corrige}[Exercice 7 — Cycle d'infection et cibles thérapeutiques]
\begin{enumerate}
    \item Enzymes virales clés :
    \begin{itemize}
        \item Étape 3 (transcription inverse) : la \textbf{transcriptase inverse} (RT), qui synthétise l'ADN proviral à partir de l'ARN viral.
        \item Étape 4 (intégration) : l'\textbf{intégrase}, qui insère l'ADN proviral dans le chromosome de la cellule hôte.
        \item Étape 7 (maturation) : la \textbf{protéase}, qui clive les polyprotéines virales pour rendre les nouveaux virions infectieux.
    \end{itemize}
    \item \textbf{INRTI/INNTI} (ex. Ténofovir, Lamivudine, Efavirenz) : ils agissent sur l'\textbf{étape 3}, en bloquant la transcriptase inverse ; l'ADN proviral ne peut pas être synthétisé, le cycle s'arrête avant l'intégration. \textbf{IP} (ex. Lopinavir, Atazanavir, Darunavir) : ils agissent sur l'\textbf{étape 7} ; les polyprotéines ne sont pas clivées, les virions produits sont \textbf{immatures et non infectieux}.
    \item Troisième classe : les \textbf{inhibiteurs d'intégrase} (ex. Dolutégravir, Raltégravir), qui bloquent l'étape 4 (intégration). On peut aussi citer les inhibiteurs d'entrée/fusion (ex. Maraviroc, antagoniste de CCR5 ; Enfuvirtide, inhibiteur de fusion), qui bloquent les étapes 1--2.
\end{enumerate}
\begin{methode}[Astuce de mémorisation]
Retenir l'ordre « \textbf{F}ixation -- \textbf{E}ntrée -- \textbf{T}ranscription inverse -- \textbf{I}ntégration -- \textbf{T}raduction -- \textbf{A}ssemblage -- \textbf{M}aturation » (F-E-T-I-T-A-M) et associer chaque classe d'ARV à son enzyme cible : INRTI/INNTI $\rightarrow$ RT, II $\rightarrow$ intégrase, IP $\rightarrow$ protéase.
\end{methode}
\end{corrige}

\begin{corrige}[Exercice 8 — Cas pratique : PTME]
\begin{enumerate}
    \item Les trois moments de transmission mère-enfant :
    \begin{itemize}
        \item \emph{in utero}, pendant la grossesse (passage transplacentaire) ;
        \item \emph{péri-partum}, au moment de l'accouchement (contact avec le sang et les sécrétions cervico-vaginales) — moment le plus à risque ;
        \item \emph{post-natal}, pendant l'allaitement (lait maternel).
    \end{itemize}
    \item \textbf{Rôle de la trithérapie maternelle :} en réduisant la charge virale plasmatique jusqu'à l'indétectabilité, elle diminue la quantité de virus susceptible de franchir le placenta (\emph{in utero}) et de contaminer l'enfant lors de l'accouchement (\emph{péri-partum}). Sans intervention, le taux de transmission est de 25 à 35 \% ; avec trithérapie + prophylaxie néonatale, il tombe en dessous de 2 à 5 \%.
    
    \textbf{Nécessité de la prophylaxie du nouveau-né :} même traitée, la mère peut présenter une charge virale résiduelle au moment de l'accouchement ; le nouveau-né a pu être exposé \emph{in utero} ou pendant le passage dans les voies génitales. Le sirop ARV (Névirapine ou Zidovudine, 6 semaines) agit comme un \emph{traitement post-exposition} qui empêche l'installation de l'infection chez l'enfant.
    
    \textbf{Allaitement sous ARV :} oui, l'OMS et le programme camerounais recommandent l'\textbf{allaitement maternel exclusif pendant 6 mois sous couverture ARV maternelle} (poursuivi jusqu'à 12 mois). Justification : dans un contexte où l'eau potable n'est pas garantie, le lait artificiel expose le nourrisson aux \textbf{diarrhées infectieuses} et à la \textbf{malnutrition}, causes majeures de mortalité infantile. Le bénéfice de l'allaitement (protection immunitaire, nutrition) dépasse le risque résiduel de transmission, devenu très faible sous trithérapie efficace. L'allaitement \emph{mixte} (lait + autres aliments avant 6 mois) est en revanche déconseillé car il irrite la muqueuse digestive et augmente le risque de transmission.
\end{enumerate}
\end{corrige}

\begin{corrige}[Exercice 9 — Sujet type Probatoire : dynamique de l'infection non traitée]
\begin{enumerate}
    \item \textbf{Phase 1 (0 à 1 an) — infection primaire (primo-infection) :} la charge virale culmine très tôt (pic à $\approx 10^{6,5}$ copies/mL vers 0,2 an) puis chute rapidement pour se stabiliser vers $10^{4,2}$--$10^{4,5}$ copies/mL ; simultanément, les CD4 chutent de 800 à 400 cellules/mm\textsuperscript{3} puis remontent partiellement (550--600). \textbf{Mécanisme :} la mise en place de la réponse immunitaire adaptative — lymphocytes CD8 cytotoxiques détruisant les cellules infectées et anticorps neutralisants (séroconversion) — contrôle partiellement la réplication virale, d'où la remontée transitoire des CD4 et la fixation du \emph{point de consigne} de la charge virale.
    \item \textbf{Phase 2 (1 à $\approx$ 8 ans) — phase asymptomatique :} la charge virale reste quasi stable (légère remontée de $10^{4,2}$ à $10^{5}$ copies/mL) et les CD4 déclinent lentement et régulièrement (de 600 à 300 cellules/mm\textsuperscript{3}). « Asymptomatique » cliniquement car le patient ne présente aucun signe de maladie ; « dynamique » virologiquement car la réplication se poursuit activement dans les organes lymphoïdes (turnover de milliards de virions/jour). \textbf{Durée moyenne :} $8 - 1 = 7$ \textbf{ans} (en général 8 à 10 ans en l'absence de traitement).
    \item \textbf{Phase 3 (8 à 12 ans) — stade Sida :} la charge virale remonte fortement (de $10^{5}$ à près de $10^{7}$ copies/mL) et les CD4 s'effondrent (de 300 à 50 cellules/mm\textsuperscript{3}). Le seuil immunologique du \textbf{Sida déclaré est franchi vers l'année 9--10, lorsque les CD4 passent sous 200 cellules/mm\textsuperscript{3}} (ici 150 à l'an 10). Apparaissent alors les infections opportunistes (tuberculose, pneumocystose, candidoses...).
    \item \textbf{Calcul :} CD4 au pic de l'infection primaire (an 0,5) : 400 cellules/mm\textsuperscript{3} ; au stade Sida (an 10) : 150 cellules/mm\textsuperscript{3}.
    \[ \text{Diminution} = \frac{400 - 150}{400} \times 100 = \frac{250}{400} \times 100 = \num{62,5}\,\% \]
    Le taux de CD4 a diminué de \textbf{\num{62,5} \%} entre ces deux points. (Si l'on prend la valeur initiale de 800 à l'an 0, la chute totale atteint $\frac{800-150}{800}\times100 = \num{81,25}\,\%$.)
    \item \textbf{Synthèse — relation inverse charge virale / CD4 :}
    \begin{itemize}
        \item \emph{Mécanisme direct (cytopathie virale) :} la réplication du VIH dans le lymphocyte CD4 aboutit au bourgeonnement de nouveaux virions et à la mort de la cellule hôte (lyse, formation de syncytiums). Plus la charge virale est élevée, plus le nombre de CD4 détruits chaque jour est grand.
        \item \emph{Mécanismes indirects :} l'activation immunitaire chronique provoque l'\textbf{apoptose} des lymphocytes CD4 \emph{non infectés} (« bystanders ») ; l'inflammation persistante détruit l'\textbf{architecture des ganglions lymphatiques} (fibrose), empêchant la régénération des CD4 ; la production médullaire et thymique ne compense plus les pertes. Le déséquilibre destruction/production explique le déclin progressif, qui s'accélère quand la charge virale remonte en phase Sida.
    \end{itemize}
\end{enumerate}
\textbf{Barème indicatif (sur 20) :} Q1 : 4 pts (description 2 + mécanisme 2) ; Q2 : 4 pts (caractérisation 2 + paradoxe 1 + calcul durée 1) ; Q3 : 3 pts ; Q4 : 3 pts (calcul détaillé exigé) ; Q5 : 6 pts (direct 2 + indirects 4).
\begin{attention}[Points de vigilance]
Exiger : la lecture quantitative des courbes (valeurs citées), le vocabulaire exact (primo-infection, point de consigne, séroconversion, latence \emph{clinique}), le seuil de 200 CD4/mm\textsuperscript{3}, et un calcul de pourcentage rapporté à la valeur de départ. Ne pas confondre « asymptomatique » avec « absence de réplication ».
\end{attention}
\end{corrige}

\begin{corrige}[Exercice 10 — Sujet type Probatoire : le dépistage volontaire, arme de santé publique]
\textbf{Introduction.} L'affirmation est pleinement justifiée : le dépistage est la \emph{porte d'entrée} de toute la cascade de prévention et de soins — sans connaissance de son statut, aucune prise en charge ni aucune prévention efficace n'est possible.

\textbf{1. Fenêtre sérologique et séroconversion.} La \textbf{séroconversion} est l'apparition d'anticorps anti-VIH détectables dans le sang, 3 à 12 semaines après la contamination. La \textbf{fenêtre sérologique} (ou période silencieuse) est l'intervalle entre l'infection et la positivation des tests : pendant cette période, la personne est \emph{infectée et très contaminante} (charge virale très élevée de la primo-infection) alors que le test d'anticorps est encore négatif. Un test négatif après une exposition récente ne garantit donc pas l'absence d'infection : la conduite à tenir est de \textbf{refaire un test de contrôle 3 mois après la dernière prise de risque} et d'adopter entre-temps des rapports protégés.

\textbf{2. Entrée précoce dans la prise en charge.} Dépisté au stade asymptomatique (CD4 encore élevés), le patient débute la trithérapie tôt : la charge virale devient indétectable, les CD4 se reconstituent (\textbf{reconstitution immunitaire}) et l'espérance de vie rejoint presque celle de la population générale. Collectivement, chaque personne traitée efficacement devient \emph{non contaminante} : le dépistage précoce tarit les sources de transmission et fait reculer l'épidémie (stratégie « Tester et Traiter », objectifs ONUSIDA 95-95-95).

\textbf{3. PTME.} Le dépistage prénatal systématique (proposé à toute femme enceinte en consultation prénatale, approche \emph{opt-out}) est la condition \emph{sine qua non} de l'élimination de la transmission verticale : seule une mère connue séropositive peut recevoir la trithérapie, la prophylaxie du nouveau-né et les conseils d'allaitement. Sans intervention, le taux de transmission mère-enfant est de \textbf{25 à 35 \%} ; avec la cascade PTME complète, il tombe \textbf{sous 2 à 5 \%} — objectif du programme national camerounais.

\textbf{4. Observance et I = I.} Une observance stricte et quotidienne maintient une pression constante sur le virus : la charge virale devient et reste \textbf{indétectable} ($< 50$ copies/mL). Les grandes études (PARTNER, HPTN 052) ont démontré qu'une personne durablement indétectable \textbf{ne transmet pas le VIH par voie sexuelle} : \textbf{I = I}. Le dépistage est la première marche de cette cascade : dépistage $\rightarrow$ mise sous traitement $\rightarrow$ observance $\rightarrow$ indétectabilité $\rightarrow$ intransmissibilité. Les campagnes communautaires au Cameroun (dépistage mobile, \emph{index testing} autour des cas identifiés, campagnes festives de dépistage) illustrent cette logique de santé publique.

\textbf{Conclusion.} Le dépistage volontaire transforme une infection autrefois mortelle en maladie chronique contrôlée et non transmissible : c'est bien une arme majeure, à la fois individuelle et collective, contre l'épidémie.

\textbf{Barème indicatif (sur 20) :} définitions et fenêtre sérologique : 4 pts ; bénéfices individuel/collectif : 5 pts ; PTME avec chiffres : 4 pts ; I = I et cascade de soins : 5 pts ; cohérence et exemples : 2 pts.
\begin{attention}[Points de vigilance]
Attendus : la distinction nette fenêtre sérologique/séroconversion, la conduite à tenir (contrôle à 3 mois), les chiffres 25--35 \% vs $< 5$ \% pour la PTME, et la condition « charge virale \emph{durablement} indétectable » pour I = I. Éviter l'affirmation « un test négatif = non infecté » sans nuance.
\end{attention}
\end{corrige}

\begin{corrige}[Exercice 11 — Sujet type Probatoire : rupture d'observance et résistance]
\begin{enumerate}
    \item \textbf{Rebond viral.} Les ARV bloquent la réplication du virus mais n'éliminent pas l'\textbf{ADN proviral intégré} dans le génome des cellules hôtes : des lymphocytes CD4 \emph{quiescents} (mémoire) hébergent des provirus silencieux et constituent les \textbf{réservoirs viraux}. À l'arrêt du traitement, la pression pharmacologique disparaît : les provirus des réservoirs sont \emph{réactivés} (transcription de l'ADN proviral en ARN viral), de nouveaux virions sont produits et infectent de nouvelles cellules CD4. En quelques semaines, la charge virale remonte massivement (ici 45\,000 copies/mL), car la réplication repart sans aucun frein.
    \item \textbf{Chute des CD4.} La réplication virale non contrôlée détruit directement les CD4 infectés (cytopathie) et indirectement les cellules voisines (activation immune chronique, apoptose). La baisse de 650 à 420 cellules/mm\textsuperscript{3} (soit une perte de 230 cellules, environ 35 \%) reflète cette destruction massive pendant les 3 mois d'arrêt ; 1 mois de reprise a fait chuter la charge virale mais la \textbf{reconstitution des CD4 est lente} (plusieurs mois), d'où un taux encore bas au contrôle.
    \item \textbf{Résistance.}
    \begin{itemize}
        \item L'arrêt \emph{simultané} des trois molécules supprime d'un coup toute pression de sélection : le virus rebondit sous forme de souche \emph{sauvage} (sensible), car aucun mutant résistant n'a d'avantage sélectif en l'absence de médicament. En revanche, un arrêt progressif ou une « monothérapie déguisée » (une seule molécule encore active, ex. la Lamivudine à demi-vie longue persistant seule) expose le virus à une \emph{pression de sélection partielle} : les mutants résistants (ex. M184V pour la 3TC) sont alors sélectionnés rapidement, car la barrière génétique de cette molécule est faible (une seule mutation suffit).
        \item Malgré tout, au rebond viral, des milliards de virions sont produits ; la transcriptase inverse étant très \emph{mutagène} (pas de relecture), chaque réplication génère des mutants. Si l'observance après reprise est imparfaite (doses oubliées, concentrations sous-thérapeutiques), les mutants partiellement résistants — notamment à la Lamivudine — disposent d'un avantage sélectif et peuvent s'installer, compromettant durablement l'efficacité du traitement.
    \end{itemize}
    \item \textbf{Message d'éducation thérapeutique (exemple attendu) :} « Madame M., votre traitement est votre bouclier : pris chaque jour, il maintient le virus \emph{indétectable}, ce qui protège vos défenses immunitaires et vous rend \emph{non contaminante} pour votre entourage (I = I). Mais le virus se cache dans des \emph{réservoirs} que le médicament ne peut pas atteindre : dès qu'on arrête, il se réveille et se multiplie à nouveau — c'est ce qui s'est passé. Pire, un traitement pris irrégulièrement peut rendre le virus \emph{résistant}, et nous perdrions alors des options de traitement précieuses. Le traitement se prend \textbf{à vie, chaque jour, sans omission}. Pour éviter les ruptures : nous pouvons vous prescrire une \textbf{dispensation multi-mois} (3 mois de médicaments) et vous orienter vers la \textbf{distribution communautaire} de proximité ou un groupe de soutien de pairs, afin de réduire vos déplacements et frais de transport. »
\end{enumerate}
\textbf{Barème indicatif (sur 20) :} Q1 : 5 pts (réservoirs 3 + réactivation 2) ; Q2 : 4 pts ; Q3 : 6 pts (pression de sélection 3 + mutagenèse/observance 3) ; Q4 : 5 pts (4 notions mobilisées + solution concrète).
\begin{attention}[Points de vigilance]
Exiger les mots-clés : \emph{ADN proviral intégré}, \emph{réservoirs}, \emph{pression de sélection}, \emph{barrière génétique}, \emph{indétectable = intransmissible}. Le message de la Q4 doit être bienveillant, non culpabilisant, et proposer une solution concrète et réaliste dans le contexte camerounais (dispensation multi-mois, approche communautaire).
\end{attention}
\end{corrige}

\newpage

\coursec{Fiche bilan}

\begin{popfigure}
\begin{tikzpicture}[
  mindmap,
  grow cyclic,
  every node/.style={concept, execute at begin node=\hskip0pt},
  concept color=popblue,
  level 1/.append style={level distance=4.5cm, sibling angle=360/6},
  level 2/.append style={level distance=3cm, sibling angle=360/6},
  text=white,
  font=\small\bfseries
]
\node[concept, scale=1.2, align=center]{VIH/Sida\\Lutte \& Prévention}
  child[concept color=poporange] { node[concept, align=center]{Structure \& Cycle\\Infection} }
  child[concept color=popgreen] { node[concept, align=center]{Transmission\\Voies réelles} }
  child[concept color=poppurple] { node[concept, align=center]{Prévention\\ABC + PrEP/PEP} }
  child[concept color=poppink] { node[concept, align=center]{Évolution\\3 phases cliniques} }
  child[concept color=popgold] { node[concept, align=center]{Dépistage\\Séroconversion} }
  child[concept color=popdark] { node[concept, align=center]{Prise en charge\\ARV + Observance} };
\end{tikzpicture}
\legende{Carte mentale récapitulative de la leçon ND08 : les 6 piliers de la lutte contre le VIH/Sida en 1ère D.}
\end{popfigure}

\begin{bilan}

\courssub{Définitions clés à connaître par cœur}

\begin{itemize}
  \item \cle{VIH (Virus de l'Immunodéficience Humaine)} : rétrovirus à ARN, enveloppé, de la famille des \emph{Lentiviridae}. Deux types : VIH-1 (mondial, plus virulent) et VIH-2 (Afrique de l'Ouest, moins virulent).
  \item \cle{CD4 (Lymphocyte T4 / T helper)} : cellule cible principale du VIH, porteur du récepteur CD4 et des co-récepteurs CCR5/CXCR4. Chef d'orchestre de l'immunité adaptative.
  \item \cle{Séroconversion} : apparition des anticorps anti-VIH détectables dans le sang (fenêtre sérologique : 3 à 12 semaines post-infection).
  \item \cle{Charge virale (CV)} : nombre de copies d'ARN VIH par mL de plasma. Marqueur de la réplication virale et de l'efficacité du traitement.
  \item \cle{Sida (Syndrome d'Immunodéficience Acquise)} : stade terminal (phase 3) défini par un taux de CD4 $< 200$/mm$^3$ et/ou la survenue d'infections opportunistes (IO) ou de cancers définissant le Sida.
  \item \cle{ARV (Antirétroviraux)} : molécules bloquant la réplication du VIH (IP, INNT, INTI, II, IE). Prescrits en \cle{trithérapie} (association de 3 molécules de 2 classes minimum).
  \item \cle{Observance} : prise correcte et continue du traitement (objectif $\ge 95\%$). Condition \textbf{sine qua non} de la suppression virale durable (CV $< 50$ copies/mL = \cle{indétectable}).
  \item \cle{TasP (Treatment as Prevention)} : personne sous ARV efficace, CV indétectable $\Rightarrow$ \textbf{risque nul de transmission sexuelle} (Indétectable = Intransmissible, I=I).
\end{itemize}

\courssub{Structure du VIH et cycle d'infection (lymphocyte CD4)}

\begin{center}
\begin{tabularx}{\linewidth}{|l|X|}\hline
\textbf{Composant viral} & \textbf{Rôle / Cible thérapeutique} \\\hline
\ce{GP120} (glycoprotéine de surface) & Fixation au récepteur \cle{CD4} puis co-récepteur \cle{CCR5} (souches R5, précoces) ou \cle{CXCR4} (souches X4, tardives). $\rightarrow$ Cible : \textbf{Inhibiteurs d'entrée (IE)} ex. Enfuvirtide, Maraviroc. \\\hline
\ce{GP41} (transmembranaire) & Fusion des membranes virale et cellulaire. $\rightarrow$ Cible : \textbf{IE} (Enfuvirtide). \\\hline
\cle{Capside (p24)} & Protéine structurelle majeure, antigène détecté par test p24 (diagnostic précoce nourrisson, fenêtre sérologique). \\\hline
\cle{ARN génomique (2 brins +)} & Matériel génétique + enzymes virales embarquées. \\\hline
\cle{Transcriptase Inverse (TI)} & Synthèse ADN $\leftarrow$ ARN (rétrotranscription). $\rightarrow$ Cible : \textbf{INTI} (Lamivudine, Emtricitabine) \& \textbf{INNTI} (Névirapine, Éfavirenz, Dolutégravir\footnote{Dolutégravir est un INTI de 2e génération, barrière génétique élevée.}). \\\hline
\cle{Intégrase} & Intégration de l'ADN viral (provirus) dans l'ADN hôte. $\rightarrow$ Cible : \textbf{II} (Raltégravir, Dolutégravir, Bictégravir). \\\hline
\cle{Protéase} & Clivage des polyprotéines virales (maturation). $\rightarrow$ Cible : \textbf{IP} (Lopinavir/r, Atazanavir/r, Darunavir/r). \\\hline
\end{tabularx}
\end{center}

\textbf{Cycle d'infection (6 étapes) :}
\begin{enumerate}
  \item \textbf{Fixation} : gp120 $\rightarrow$ CD4 + co-récepteur (CCR5/CXCR4).
  \item \textbf{Fusion/Pénétration} : gp41 médie la fusion $\rightarrow$ libération capside dans le cytoplasme.
  \item \textbf{Rétrotranscription} : TI $\rightarrow$ ADN double brin (provirus). \emph{Étape error-prone $\rightarrow$ mutations/résistance}.
  \item \textbf{Intégration} : Intégrase $\rightarrow$ insertion dans chromosome hôte (latence possible).
  \item \textbf{Transcription/Traduction} : Machinerie hôte $\rightarrow$ ARNm viraux $\rightarrow$ polyprotéines (Gag, Pol, Env).
  \item \textbf{Assemblage/Budding/Maturation} : Protéase clive polyprotéines $\rightarrow$ virions matures infectieux.
\end{enumerate}

\courssub{Transmission et Prévention}

\begin{center}
\begin{tabularx}{\linewidth}{|l|X|X|}\hline
\textbf{Voie} & \textbf{Mécanisme / Fluides infectants} & \textbf{Prévention efficace} \\\hline
\textbf{Sexuelle (80\% cas Cmr)} & Sperme, sécrétions vaginales, sang menstruel, lésions muqueuses (IST non traitées $\uparrow$ risque $\times 10$). & \textbf{Préservatif} (M/F) \textbf{systématique \& correct}. \textbf{PrEP} (Truvada) si risque élevé. \textbf{TasP} (partenaire séropositif sous ARV efficace). \\\hline
\textbf{Sanguine} & Sang total, produits sanguins, matériel d'injection contaminé (usagers de drogues, accidents d'exposition AES, scarifications, tatouages non stériles). & Sécurité transfusionnelle (dépistage donneurs). \textbf{Matériel stérile / à usage unique}. \textbf{PEP} (Prophylaxie Post-Exposition) $< 48$h (idéal $< 4$h) sur 28 jours. \\\hline
\textbf{Mère-Enfant (VTME)} & Grossesse (transplacentaire), Accouchement (contact sang/sécrétions), Allaitement (lait maternel). & \textbf{Dépistage ANC1} (1ère consultation prénatale). \textbf{ARV mère} (Option B+ : trithérapie à vie dès diagnostic). \textbf{ARV nourrisson} (Névirapine/Zidovudine 6 sem). \textbf{Allaitement exclusif 6 mois} + ARV mère $\rightarrow$ risque $< 2\%$. \\\hline
\end{tabularx}
\end{center}

\courssub{Évolution de l'infection (Histoire naturelle sans traitement)}

\begin{center}
\begin{tabularx}{\linewidth}{|c|X|X|X|}\hline
\textbf{Phase} & \textbf{Clinique} & \textbf{Biologique (CD4 / Charge Virale)} & \textbf{Durée moyenne} \\\hline
\textbf{1. Primo-infection / Séroconversion} (J+15 à J+90) & Syndrome pseudo-grippal / mononucléosique (fièvre, adénopathies, éruption, méningite aseptique). Souvent méconnu. & \textbf{CV : Pic extrême} ($> 10^6$ copies/mL). \textbf{CD4 : Chute brutale}. \textbf{Séroconversion} : Ac anti-VIH apparaissent (fin phase). & 2--6 semaines \\\hline
\textbf{2. Phase asymptomatique / Latence clinique} & \textbf{Aucun signe} (ou adénopathies généralisées persistantes PAG). Patient se sent bien. & \textbf{CV : Plateau} (point de set-point viral). \textbf{CD4 : Déclin lent} ($\sim 50-80$/mm$^3$/an). & \textbf{5--10 ans} (variable) \\\hline
\textbf{3. Phase Sida déclaré} & \textbf{Infections Opportunistes (IO)} : Candidose œsophagienne, Pneumocystose, Toxoplasmose cérébrale, Tuberculose (forme disséminée), Cryptococcose, CMV, Herpès sévère... \textbf{Cancers} : Sarcome de Kaposi, Lymphome, Cancer col utérin. \textbf{AMA} (Amaigrissement Majeur). & \textbf{CD4 $< 200$/mm$^3$} (seuil définition Sida OMS). \textbf{CV : Très élevée}. & Issue fatale sans ARV \\\hline
\end{tabularx}
\end{center}

\courssub{Dépistage, Prise en charge et Traitements}

\begin{methode}[Algorithme de dépistage (Stratégie OMS / Cameroun)]
\begin{enumerate}
  \item \textbf{Test de dépistage (TDR / ELISA)} : Détecte Ac anti-VIH-1/2 + Ag p24 (4e génération = combiné). \textbf{Négatif} $\rightarrow$ Non infecté (si hors fenêtre). \textbf{Positif} $\rightarrow$ Étape 2.
  \item \textbf{Test de confirmation (Western Blot ou 2e TDR différent)} : \textbf{Positif} $\rightarrow$ \textbf{Séropositivité confirmée}. \textbf{Indéterminé} $\rightarrow$ Contrôle à J+14 / J+30 (séroconversion en cours ?).
  \item \textbf{Conseil pré/post-test} obligatoire (confidentialité, consentement, orientation).
\end{enumerate}
\end{methode}

\begin{methode}[Bilan initial avant mise sous ARV (Protocole Cameroun)]
\begin{enumerate}
  \item \textbf{Confirmation sérologique} + Typage VIH-1/2.
  \item \textbf{Numération CD4} (stade immunologique, urgence traitement si $< 200$ ou stade 3/4 OMS).
  \item \textbf{Charge Virale (CV) de référence} (si dispo).
  \item \textbf{Recherche co-infections / comorbidités} : VIH/TB (recherche BK systématique), Hépatites B/C (HBsAg, Anti-HCV), IST, Grossesse, Fonction rénale (Créatinine/DFG pour TDF), Hb/NFS (anémie pour AZT).
  \item \textbf{Conseil observance} (éducation thérapeutique, choix schéma, effets indésirables).
  \item \textbf{Initiation ARV} : \textbf{Treat All} (tout séropositif, quelque soit CD4). Schéma 1ère ligne recommandé : \textbf{TDF + 3TC (ou FTC) + DTG} (Dolutégravir). Alternative : TDF+3TC+EFV 400mg.
\end{enumerate}
\end{methode}

\begin{propriete}[Critères de succès virologique (Suivi sous ARV)]
\begin{itemize}
  \item \textbf{Suppression virale} : CV $< 50$ copies/mL (indétectable) à J+6 mois (but ultime).
  \item \textbf{Échec virologique} : CV $> 1000$ copies/mL à deux reprises (espacées de 3 mois) avec bonne observance $\rightarrow$ \textbf{Test de résistance} + passage 2ème ligne (ex: AZT+3TC+ATV/r ou LPV/r).
  \item \textbf{Reconstitution immunitaire} : Hausse CD4 $> 100-150$/mm$^3$ la 1ère année. Risque \cle{SRII} (Syndrome de Réconstitution Immunitaire Inflammatoire) au démarrage si IO préexistante (TB, Crypto).
\end{itemize}
\end{propriete}

\courssub{Formules / Relations quantitatives utiles}

\[
\text{Ratio CD4/CD8} < 1 \text{ (inversion, signe d'activation immunitaire chronique)}
\]
\[
\text{Log drop CV} = \log_{10}(\text{CV}_{\text{initiale}}) - \log_{10}(\text{CV}_{\text{suivi}}) \quad (\text{viser } \ge 1 \log \text{ à 1 mois, } \ge 2 \log \text{ à 3 mois})
\]
\[
\text{Incidence} = \frac{\text{Nouvelles infections}}{\text{Population saine} \times \text{Durée}} \quad ; \quad \text{Prévalence} = \frac{\text{Cas existants}}{\text{Population totale}}
\]

\end{bilan}

\begin{aretenir}[Checklist avant le Probatoire]
\begin{itemize}
  \item $\square$ Dessiner et légender le VIH (gp120, gp41, capside p24, ARN, TI, Intégrase, Protéase) et nommer la classe d'ARV ciblant chaque enzyme.
  \item $\square$ Décrire dans l'ordre les 6 étapes du cycle d'infection du lymphocyte CD4 (Fixation $\rightarrow$ Maturation).
  \item $\square$ Citer les 3 voies de transmission réelles + 1 fluide infectant par voie + 1 mesure préventive spécifique par voie (Préservatif, PrEP, PEP, Sécurité transfusionnelle, PTME/Option B+).
  \item $\square$ Différencier \cle{fenêtre sérologique}, \cle{séroconversion} et \cle{phase asymptomatique} (durée, signes, CV, CD4, sérologie).
  \item $\square$ Construire et interpréter un tableau de l'évolution CV / CD4 / Clinique sur les 3 phases (pic CV phase 1, plateau phase 2, effondrement CD4 phase 3).
  \item $\square$ Expliquer le principe du test combiné 4e génération (Ag p24 + Ac) et l'algorithme de confirmation (Test 1 $\rightarrow$ Test 2).
  \item $\square$ Lister le bilan pré-thérapeutique minimal (CD4, CV, HBsAg, Anti-HVC, BK, Créatinine, NFS, Grossesse).
  \item $\square$ Écrire le schéma 1ère ligne ARV recommandé au Cameroun (TDF+3TC+DTG) et citer 2 effets indésirables majeurs à surveiller par classe (TDF: rénal/os ; DTG: prise de poids/neuro ; EFV: neuro/teratogène).
  \item $\square$ Définir \cle{Observance} ($\ge 95\%$), \cle{Échec virologique} (CV $> 1000$ copies/mL confirmée), \cle{TasP / I=I} (CV indétectable = risque nul transmission).
  \item $\square$ Argumenter : « Pourquoi dépister précocement ? » (Bénéfice individuel : espérance vie normale ; Bénéfice collectif : prévention TasP, élimination PTME).
  \item $\square$ Citer 3 infections opportunistes définissant le Sida (stade 4 OMS) et l'agent responsable (ex: \emph{Pneumocystis jirovecii}, \emph{Toxoplasma gondii}, \emph{Mycobacterium tuberculosis}).
  \item $\square$ Calculer une prévalence ou une incidence à partir de données brutes et interpréter le sens pour la santé publique.
\end{itemize}
\end{aretenir}

\findecours

\end{document}
